\documentclass[11pt,oneside,openany]{book}
\usepackage[paperwidth=7.5in,paperheight=10in,top=0.78in,bottom=0.78in,left=0.82in,right=0.82in,headheight=15pt,headsep=0.22in]{geometry}
\ifdefined\pdfinfoomitdate\pdfinfoomitdate=1\fi
\ifdefined\pdftrailerid\pdftrailerid{}\fi
\ifdefined\pdfsuppressptexinfo\pdfsuppressptexinfo=15\fi
\usepackage[T1]{fontenc}
\usepackage[utf8]{inputenc}
\usepackage{amsmath,amsthm,mathtools}
\usepackage{newtxtext,newtxmath}
\usepackage{microtype}
\usepackage[dvipsnames]{xcolor}
\usepackage{enumitem,booktabs,tabularx,longtable,array}
\usepackage{tikz}
\usetikzlibrary{arrows.meta,positioning,calc,decorations.pathreplacing}
\usepackage[most]{tcolorbox}
\usepackage{listings,fancyhdr,titlesec,needspace}
\usepackage[unicode,hidelinks]{hyperref}
\usepackage{bookmark}
\definecolor{ink}{HTML}{17364A}
\definecolor{accent}{HTML}{336F74}
\definecolor{papergray}{HTML}{F2F5F5}
\hypersetup{pdftitle={Mathematics in Collaboration: A One-Semester Course},pdfauthor={Mathematics in Collaboration},pdfsubject={Proof literacy, mathematical structures, AI collaboration, and research apprenticeship}}
\setlength{\parindent}{1em}
\setlength{\parskip}{3pt}
\setlength{\emergencystretch}{2em}
\linespread{1.035}
\raggedbottom
\setlist{topsep=5pt,itemsep=3pt,parsep=1pt,leftmargin=1.65em}
\setcounter{tocdepth}{1}
\setcounter{secnumdepth}{1}
\newcommand{\chapterkind}{CHAPTER}
\titleformat{\chapter}[display]{\color{ink}\normalfont\Large\bfseries\raggedright\hyphenpenalty=10000\exhyphenpenalty=10000}{\large \chapterkind\ \thechapter}{8pt}{\Huge}
\titlespacing*{\chapter}{0pt}{-14pt}{22pt}
\titleformat{\section}{\large\bfseries\color{ink}}{\thesection}{0.7em}{}
\titleformat{\subsection}{\normalsize\bfseries}{\thesubsection}{0.7em}{}
\pagestyle{fancy}\fancyhf{}
\fancyhead[L]{\small\nouppercase{\leftmark}}
\fancyfoot[L]{\small Mathematics in Collaboration}
\fancyfoot[R]{\thepage}
\renewcommand{\headrulewidth}{0.3pt}
\renewcommand{\chaptermark}[1]{\markboth{\thechapter\enspace #1}{}}
\fancypagestyle{plain}{\fancyhf{}\fancyfoot[L]{\small Mathematics in Collaboration}\fancyfoot[R]{\thepage}\renewcommand{\headrulewidth}{0pt}}
\theoremstyle{plain}
\newtheorem{theorem}{Theorem}[chapter]
\newtheorem{lemma}[theorem]{Lemma}
\newtheorem{proposition}[theorem]{Proposition}
\theoremstyle{definition}
\newtheorem{definition}[theorem]{Definition}
\newtheorem{example}[theorem]{Example}
\newtheorem{exercise}{Exercise}[chapter]
\newtcolorbox{principle}[1]{enhanced,colback=papergray,colframe=accent,boxrule=0pt,borderline west={2pt}{0pt}{accent},arc=0pt,left=9pt,right=8pt,top=7pt,bottom=7pt,title={#1},coltitle=ink,fonttitle=\bfseries,attach title to upper={\par\smallskip}}
\newenvironment{lab}[1]{\section*{AI laboratory: #1}\addcontentsline{toc}{subsection}{AI laboratory: #1}}{}
\newcommand{\pauseq}[2]{\par\medskip\noindent\textbf{Pause and decide.} #1\par\smallskip\noindent\textit{Check.} #2\par\medskip}
\newcommand{\goal}[1]{\begin{principle}{This week}#1\end{principle}}
\newcommand{\exref}[1]{Exercise~\ref{#1}}
\newcommand{\sol}[1]{\subsection*{Solution to Exercise #1}\phantomsection\label{sol:#1}}
\newcommand{\NN}{\mathbb N}\newcommand{\ZZ}{\mathbb Z}\newcommand{\QQ}{\mathbb Q}\newcommand{\RR}{\mathbb R}
\newcommand{\abs}[1]{\left|#1\right|}\newcommand{\norm}[1]{\left\|#1\right\|}
\newcommand{\set}[1]{\left\{#1\right\}}
\newcommand{\eps}{\varepsilon}
\DeclareMathOperator{\im}{im}\DeclareMathOperator{\id}{id}
\DeclareMathOperator{\dist}{dist}\DeclareMathOperator{\diam}{diam}
\DeclareMathOperator{\E}{\mathbb E}\DeclareMathOperator{\Prob}{\mathbb P}
\lstset{basicstyle=\ttfamily\small,columns=fullflexible,keepspaces=true,breaklines=true,showstringspaces=false,frame=lines,rulecolor=\color{accent},aboveskip=9pt,belowskip=9pt,xleftmargin=3pt,upquote=true,keywordstyle=\bfseries,commentstyle=\itshape}
\newcommand{\code}[1]{\texttt{#1}}
\begin{document}
\frontmatter
\begin{titlepage}
\color{ink}
\noindent{\small A ONE-SEMESTER COURSE}\par
\vspace{1.05in}
\noindent{\fontsize{33}{38}\selectfont\bfseries Mathematics\\in Collaboration\par}
\vspace{0.35in}
\noindent{\Large Learning to ask, understand, prove,\\and investigate with AI\par}
\vspace{0.5in}\noindent\rule{1.25in}{1.5pt}
\vspace{0.35in}
\noindent{\large From first mathematical sentences to\\complete theorem proofs and research reading\par}
\vfill
\noindent Eighteen weeks\quad\textbullet\quad Worked proofs\quad\textbullet\quad Guided laboratories\\
Exercises and complete solutions\quad\textbullet\quad Lean and computational companions
\par\vspace{0.25in}
\noindent Reader edition 3.0\enspace /\enspace 7 October 2026
\end{titlepage}
\chapter*{To the reader}
\addcontentsline{toc}{chapter}{To the reader}
You are not expected to have taken a university mathematics course. You are expected to be able to add, subtract, multiply, and divide ordinary numbers. Chapter~1 rebuilds the algebraic notation we need. No calculus, linear algebra, programming, or previous experience with proofs is assumed. We explain the language of definitions, theorems, and proofs before asking you to use it. Ordinary arithmetic is our starting point; Chapter~11 introduces and explains the additional property of real numbers needed for limits. We do not construct the number system from formal logic.

This book asks a different question from ``How much mathematics can I memorize before I begin?'' The question is: \emph{How can I enter a mathematical conversation, discover what is missing from my understanding, and learn enough to participate responsibly?} AI is a participant in that conversation, not a replacement for your understanding. It can propose examples, explanations, arguments, and code. You must learn how to ask for the missing bridge, identify a changed assumption, and explain what the resulting argument actually establishes.

The course is organized around three recurring problems. Can a collection of people be assigned distinct acceptable tasks? Can repeated approximate updates converge to an exact solution? Can a sequence of plus and minus signs have prescribed correlations? These problems begin with finite examples. They eventually lead to Hall's matching theorem, the contraction theorem, and a concrete entry into research on circulant matrices. A fourth strand, finite probability, shows how existence can follow from averaging rather than from a direct construction.

You will see definitions first as answers to problems. You will see proofs first as plans, then as complete arguments, then as objects to modify and challenge. The guiding habits---choose a revealing representation, track assumptions, find an obstruction, exploit symmetry, and control error---are practiced repeatedly rather than collected in a detached list.

\section*{How to use a chapter}
Read the opening problem before asking an assistant for a solution. Try a small example. Read the explanation, stopping at every ``Pause and decide.'' Cover the check before answering. In a proof, distinguish a line whose truth you understand from a line whose purpose you understand; you need both. Each laboratory contains a scripted exchange written for teaching, not a record of an empirical AI evaluation. The deliberately flawed replies are labeled. Actual models may respond differently.

Finish the four core exercises. The two transfer exercises revisit the same ideas with a change of setting or hypothesis. Complete solutions to all six appear at the end of the book. Consult a solution after a genuine attempt, then reconstruct it with the book closed. Research activities have worked baselines and explicit assessment criteria; they do not have a promised new theorem hidden in the answer key.

\section*{The endpoint}
By the end, you should be able to reconstruct several complete proofs, formulate a minimal prerequisite plan for an unfamiliar statement, run and audit small mathematical experiments, read elementary Lean proofs, and investigate a carefully bounded question with AI. You will not automatically understand every research paper. A specialized proof may require substantial additional study. The skill developed here is to discover precisely which study is needed, carry it out, and reconnect it to the target proof.

Novelty is a separate question from correctness. The worked theorems and capstone baseline are established mathematics or elementary consequences, not claims of new research. A student investigation becomes a research contribution only after its relation to prior work and its mathematical evidence have been examined.

\chapter*{The semester as one argument}
\addcontentsline{toc}{chapter}{The semester as one argument}
The schedule assumes eighteen teaching weeks, three contact hours per week, and approximately four to six hours of individual reading, exercises, and laboratory work. These are planning assumptions, not measured completion times. Chapters~1--6 deliberately proceed slowly. Chapters~8--9 and 11--12 form consecutive two-week developments; do not separate their proofs from their preparation. Weeks~17--18 integrate the earlier work rather than introduce a new list of subjects.

The chapter is a week's reading, not a script to finish in one sitting. Divide it at its section headings. For a long proof, one sitting may consist of understanding the statement and one smaller argument. On returning, first say what has already been established and what remains. The worked examples are part of the route, not optional background for a supposedly slower reader. In a taught course, use the two transfer exercises for additional practice when the four core exercises are secure; do not increase the required workload merely because the explanations are longer.

\begingroup\small
\begin{longtable}{@{}p{0.34in}p{2.12in}p{2.87in}@{}}
\toprule Week & Mathematical development & Concrete output\\\midrule\endhead
1 & Variables, representation, and explanation & A precise claim and a fully explained parity proof.\\
2 & Logic and dependence & A theorem translated, negated, and tested.\\
3 & Sets and maps & A composition proof and a counterexample to a converse.\\
4 & Building and reading arguments & A proof annotated by goal, fact, and justification.\\
5 & Induction and invariants & A recurrence proof and a terminating algorithm.\\
6 & Equivalence and symmetry & A well-defined construction on residue classes.\\
7 & Graphs, counting, and extremal choice & A counting proof and a shortest-path argument.\\
8 & Assignments and augmenting paths & An algorithm explained with its invariant.\\
9 & Hall's theorem and proof reconstruction & A complete long proof and the midterm portfolio.\\
10 & Vectors, linearity, and representation & A kernel criterion and an inner-product inequality.\\
11 & Distance, limits, and completeness & A convergence argument with explicit tolerances.\\
12 & The contraction theorem & A full proof with a numerical error certificate.\\
13 & Averaging and finite probability & An existence proof and a deterministic construction.\\
14 & Sign sequences and research experiments & Exact correlations and a necessary-order theorem.\\
15 & Lean: statements and proof states & Small proofs reconstructed in a formal language.\\
16 & Lean: definitions, trust, and repair & A formalization with an assumption audit.\\
17 & Learning an unfamiliar proof with AI & A prerequisite map and a repository reading dossier.\\
18 & A bounded research investigation & A complete capstone baseline and a defended extension.\\
\bottomrule
\end{longtable}
\endgroup

\section*{A repeatable weekly rhythm}
In the first meeting, unpack the opening problem and jointly derive one example. In the second, reconstruct the main proof and compare two possible next moves. In the third, carry out the AI laboratory and critique a proposed argument. For independent study, use the same sequence: first attempt, guided reading, reconstruction, transfer. A student who already knows an example should explain it under a changed assumption rather than merely skip it.

\section*{Assessment without outsourcing understanding}
A suggested allocation is 30\% for weekly core exercises, 20\% for laboratory records, 20\% for the Week~9 portfolio, and 30\% for the final investigation. AI assistance is allowed when its role is disclosed. Each submission must contain a short unaided explanation of its central inference, one tested boundary case, and a statement of what remains borrowed or unchecked. Grade the mathematical content and the student's demonstrated understanding, not the length of the conversation log or the polish of AI-generated prose.

The Week~9 portfolio contains a correct Hall proof, an instance that satisfies the condition, an instance that fails it, and an explanation of why a greedy assignment can fail. The final dossier contains an exact claim, a minimal prerequisite map, a complete baseline proof, an audited experiment, and one justified variant. Chapter~18 supplies the worked model; Appendix~\ref{app:guide} supplies the rubric.

\section*{Keeping track of evidence}
Use two separate columns in your notebook. The first records understanding: \emph{can explain}, \emph{can follow}, or \emph{not yet understood}. The second records support: \emph{proved here}, \emph{externally cited}, \emph{computed in a finite range}, or \emph{not established}. A formally checked proof can still be poorly understood; a clearly understood experimental pattern can still be unproved. Do not collapse these two dimensions into one confidence score.

\tableofcontents
\mainmatter


% ===== Week 1 =====

\chapter{A first mathematical conversation}\label{ch:start}
\goal{Turn a familiar numerical pattern into a statement and a proof. Learn the notation needed to describe a pattern without assuming it is true. Begin using AI to expose a missing step rather than to hide it.}

\section{What does a pattern promise?}
Suppose two people collect an odd number of tokens. If they combine their collections, must the total be even? Try $3+5=8$, $7+9=16$, and $1+1=2$. The examples agree. They do not yet explain what happens for two odd numbers we have not tried.

A mathematical statement makes its scope explicit: \emph{For every pair of odd integers, their sum is even.} An integer is a whole number, possibly negative or zero: $\ldots,-2,-1,0,1,2,\ldots$. We write their collection as $\ZZ$. The statement includes negative odd integers as well as positive ones. For example, $-3+5=2$. A calculation checks a particular pair. A proof explains why every allowed pair works.

The first collaborative question is therefore not ``Can you prove this?'' but ``What does the claim range over, and what exactly must hold?'' An assistant that quietly replaces all integers by positive integers has changed the question. The new statement may still be true, but it is not the same statement.

\section{Letters are names, not mysterious numbers}
A letter such as $n$ can stand for an integer. The expression $2n$ means $2\times n$; $n+1$ is one more than $n$; and $n^2$ means $n\times n$. Parentheses group an expression: $2(n+1)$ is twice the entire quantity $n+1$. It equals $2n+2$, not $2n+1$.

An equality says that its two sides name the same value. It does not mean ``now calculate the next thing.'' For example,
\[
2(3+1)=2\cdot4=8
\]
is legitimate because all three expressions have the same value. The chain $2+3=5\cdot2=10$ is not legitimate: its first expression is $5$, but its middle expression is $10$. Write two sentences or two separate equations when the values differ.

The arithmetic rules we use are familiar. Addition can be reordered and regrouped, and multiplication distributes over addition:
\[
a+b=b+a,\qquad (a+b)+c=a+(b+c),\qquad a(b+c)=ab+ac.
\]
The same integer may have many expressions. For example, $12=3\cdot4=2\cdot6$. Choosing an expression is not merely cosmetic. The second product immediately displays evenness, whereas the first does not.

A fraction $a/b$ means $a$ divided by $b$ and requires $b\ne0$. The symbol $\ne$ means ``is not equal to.'' Multiplying both sides of an equality by the same number preserves equality. Dividing both sides does too, provided the divisor is not zero. This condition will be responsible for several important counterexamples.

\pauseq{An assistant starts from $x^2=x$, divides by $x$, and concludes $x=1$. Has it found every solution?}{No. Division assumes $x\ne0$, but $x=0$ satisfies the original equation. A complete argument separates $x=0$ from $x\ne0$. In the second case division gives $x=1$. Both $0$ and $1$ are solutions.}

\section{The small algebra moves we will use}
An expression such as $2k+1$ is built from operations you already know. To evaluate it at $k=3$, replace every occurrence of $k$ by $3$: $2k+1$ becomes $2\cdot3+1=7$. This replacement is called \emph{substitution}. Substitution does not say that $k$ always equals three; it checks one chosen input.

Sometimes we give the rule itself a name. Writing $f(n)=n+1$ means that the rule named $f$ adds one to its input $n$; for example, $f(3)=4$. Read $f(n)$ as ``the output of $f$ at input $n$,'' not as $f$ multiplied by $n$. Chapter~3 develops functions fully. For now this notation lets us refer to a formula and its values without rewriting the formula each time.

To rewrite an expression without changing its value, use an arithmetic rule in either direction. Distributing gives $2(k+1)=2k+2$. Reading that equality in reverse gives $2k+2=2(k+1)$; this is called \emph{factoring}. In a proof, factoring is often useful because it reveals a common multiplier.

For a slightly longer example, apply distribution one bracket at a time:
\begin{align*}
 (2k+1)(2\ell+1)
 &=2k(2\ell+1)+1(2\ell+1)\\
 &=4k\ell+2k+2\ell+1.
\end{align*}
The first line gives each term of the first bracket a copy of the second bracket. The second line multiplies out those two copies. The letter $\ell$, read ``ell,'' is a different name from $k$, not the digit one. We need two names because the inputs may differ.

There is also a difference between an identity and an equation to solve. The equality $2(k+1)=2k+2$ holds for every integer $k$; it is an \emph{identity}. The equation $2k+1=7$ asks which values of $k$ make it true. Subtracting one from both sides gives $2k=6$; dividing by two gives $k=3$. Substitution into the original equation checks that $2\cdot3+1=7$.

A power records repeated multiplication: $2^3=2\cdot2\cdot2=8$. The exponent three counts the factors, not additions. We use $a^0=1$ for $a\ne0$, and $a^{-n}=1/a^n$ for positive integers $n$ and nonzero $a$. In particular, $2^{-3}=1/8$. Joining a row of $n$ factors and a row of $m$ factors gives $a^{n+m}=a^na^m$. For nonzero $a$, canceling $m$ factors gives $a^n/a^m=a^{n-m}$. We will revisit these rules with actual numbers whenever they carry a new argument.

\pauseq{Why is $(a+b)^2$ not generally $a^2+b^2$?}{The square means $(a+b)(a+b)$. Distribution gives $a^2+ab+ba+b^2=a^2+2ab+b^2$. For $a=1,b=2$, the two proposed values are nine and five. The missing cross terms account for the difference.}

\section{A definition that makes a proof possible}
A \emph{definition} gives an exact meaning to a word or symbol. It is an agreement about how that word will be used, not a claim that a pattern has been discovered. We need a meaning precise enough to decide which numbers qualify.

Here is our agreement. An integer is \emph{even} when it is twice an integer. It is \emph{odd} when it is one more than twice an integer. Thus $8$ is even because $8=2\cdot4$, and $7$ is odd because $7=2\cdot3+1$. Negative integers are allowed: $-3=2(-2)+1$ is odd. Zero is even because $0=2\cdot0$.

When a definition is displayed with a label, it has the same role as this explanation. The number on the label is only a location: Definition~1.1 is the first numbered item in Chapter~1. It does not indicate difficulty or a new mathematical operation:
\par\noindent\begin{minipage}{\linewidth}
\begin{definition}[Even and odd integers]\label{def:parity}
An integer $a$ is even if it can be written as $a=2k$ for an integer $k$. It is odd if it can be written as $a=2k+1$ for an integer $k$.
\end{definition}
\end{minipage}\par

The phrase ``for an integer $k$'' is essential. Without requiring $k$ to be an integer, every integer $a$ would have the form $2k$, simply by choosing $k=a/2$. The definition is not merely asking whether a number can be divided by two. It asks whether the result is an integer.

A particular value showing that a definition's requirement is met is called a \emph{witness}. For the odd integer $7$, the witness is $k=3$. It is not the number $7$ itself: it is the auxiliary integer that makes $7=2k+1$ true. For $-3$, the witness is $-2$. Different inputs may need different witnesses.

A definition can be used in two directions. If someone tells us $a$ is odd, we may introduce an integer $k$ with $a=2k+1$. Conversely, if we produce such an integer and equality, we have shown $a$ is odd. The first direction unpacks information; the second recognizes when the required information has been supplied. These are the two uses that our first proof will combine.

\section{Definition, example, claim, theorem, proof}
Before reading a numbered result, distinguish five jobs that mathematical sentences can do.

\emph{A definition sets the meaning.} ``Odd means one more than twice an integer'' supplies a rule for using a word. There is no pattern to verify before adopting the rule. However, showing that a particular object meets the rule may require work.

\emph{An example shows the meaning in one instance.} The equation $7=2\cdot3+1$ shows how the definition applies. A \emph{nonexample} does not meet the definition. For instance, $4$ is not odd: if $4=2k+1$, solving gives $k=3/2$, not an integer. Both kinds of example help us understand the boundary of a definition.

\emph{A claim says something that may need justification.} ``Adding two odd integers always gives an even integer'' is not part of the definitions. We must establish that the output has the required form whenever the inputs have their required forms. Until this is done, agreeing with a few examples is not enough.

\emph{A theorem is a claim supplied with a proof.} The labels \emph{proposition} and \emph{lemma} use the same standard of justification. As a heading, \emph{Proposition} labels an assertion presented with a proof; a lemma is an assertion used to support another result. A theorem is usually a result the author wishes to emphasize. These are organizational labels, not increasing levels of truth or permissions to skip steps. A \emph{corollary} is a result deduced from an earlier result. A \emph{conjecture} is a proposed statement not yet established in the discussion.

There is a second use of the word \emph{proposition} that we will meet in logic and Lean: it can mean any statement that can be true or false, whether or not we have proved it. Thus calling an expression a proposition does not assert its truth. The numbered heading in a textbook and the general word for a logical statement have different roles.

\emph{A proof supplies the justification.} It starts from the allowed inputs and facts, makes justified deductions, and arrives at exactly the promised conclusion. Its purpose is not to reassure us that the author is confident. Its purpose is to let us see why no allowed case has been omitted. The small square at the end of a proof means ``the argument ends here''; it is not an operation to perform.

A \emph{hypothesis}, also called an assumption or premise, is a condition under which a result is asserted. In our first proposition, the hypotheses are that the two inputs are integers and are odd. The \emph{conclusion} is that their sum is even. We may use the hypotheses during the proof. We may not simply assume the conclusion to obtain it.

Some facts are accepted as the starting rules of a theory; these are called \emph{axioms}. In this course we also explicitly accept familiar arithmetic facts without rebuilding their foundations. That does not permit us to declare a difficult claim true whenever a proof is missing. We distinguish a stated starting fact from something the argument is meant to establish.

\pauseq{Classify these sentences: ``Even means twice an integer''; ``Six is even because $6=2\cdot3$''; ``The sum of two odd integers is even.''}{The first is a definition, the second verifies an example, and the third is a general claim. The third becomes an established result in this book when the proof below has been completed.}

\section{How to enter the first proof}
A mathematical writer often begins with ``let'' or ``take.'' Here ``let $a$ and $b$ be any odd integers'' means: give names to arbitrary allowed inputs. It does not mean choose your favorite two examples. We must use only their stated properties so that the argument works whichever odd integers were supplied.

The word ``some'' has a different job. Once $a$ is given and is known to be odd, \emph{some} integer $k$ makes $a=2k+1$ true. We need not know its numerical value to use its existence. ``By the definition'' tells us where that information came from. ``Hence'' or ``therefore'' announces a deduction; the reason for the deduction must still be available.

Read the next result in three passes. First identify the input conditions and the promised output. Then read the plan to see why a particular representation is being chosen. Finally follow the proof, checking both the arithmetic and the last return to the definition. The plan explains how we might find the proof; the proof is the argument that actually establishes the claim.

\par\noindent\begin{minipage}{\linewidth}
\begin{proposition}\label{thm:odd-sum}
The sum of any two odd integers is even.
\end{proposition}
\end{minipage}\par

\noindent\textit{Plan.} Represent each input so that oddness is visible. Rearrange their sum until the definition of evenness is visible.
\begin{proof}
Let the two odd integers be $a$ and $b$. By the definition of oddness, there are integers $k$ and $\ell$ with $a=2k+1$ and $b=2\ell+1$. The letters $k$ and $\ell$ need not denote the same integer. Substituting gives
\begin{align*}
a+b&=(2k+1)+(2\ell+1) &&\text{use the two representations}\\
&=2k+2\ell+2 &&\text{regroup the four terms}\\
&=2(k+\ell+1) &&\text{factor out }2.
\end{align*}
The quantity $k+\ell+1$ is an integer because sums of integers are integers. Hence $a+b$ is twice an integer. This is exactly the definition of evenness.
\end{proof}

The final sentence matters. It connects the expression to the property we promised. A proof should not end simply because the algebra looks tidy. It should end because the required conclusion has been reached.

To see how one general line covers many examples, return to $a=7$ and $b=-3$. Their witnesses are $k=3$ and $\ell=-2$. The proof constructs the evenness witness $k+\ell+1=2$, and indeed $7+(-3)=4=2\cdot2$. The numerical check is not a second proof. It lets us inspect what each symbol in the general proof was doing.

The proof does not say that every sum is even. Its conclusion is conditional on \emph{both} inputs being odd. If one input is not odd, this particular result makes no promise. Learning to keep the assumptions attached to a result is as important as remembering the result itself.

\begin{principle}{Habit 1: choose a representation that exposes the target}
If the goal is evenness, try to display a factor of two. If the goal is an assignment, display who receives what. If the goal is an error bound, display the difference between the approximation and the target. Ask an assistant why a proposed representation helps before letting it continue a calculation.
\end{principle}

\section{Reading a formula without losing the picture}
Sometimes we need to add many quantities. The notation
\[
\sum_{j=1}^{4}a_j=a_1+a_2+a_3+a_4
\]
means add the entries named $a_1$ through $a_4$. The symbol $j$ is a temporary index: it takes the values $1,2,3,4$ in turn. A subscript is a label, not multiplication. For example, $a_2$ is the second entry, whereas $2a$ is twice $a$.

Suppose the entries are $1,3,5,7$. Their total is $16$. Changing the temporary letter from $j$ to $r$ does not change the sum. Changing the upper limit from $4$ to $3$ does: we would obtain $1+3+5=9$. Small pieces of notation can express substantial changes in a claim.

We will use $\NN=\{0,1,2,\ldots\}$ for natural numbers. Braces enclose the members of a collection, and the dots mean the displayed pattern continues. We use $\QQ$ for numbers expressible as fractions of integers with nonzero denominator, and $\RR$ for real numbers on the number line. Integers and fractions are real numbers; real numbers also include lengths not expressible as such fractions. When a letter is said to be real, these are the allowed inputs, not merely the positive whole numbers. When counting a positive number of objects, we will say $n\ge1$ explicitly. The notation $x\ge y$ means $x$ is at least $y$; $x>y$ means strictly larger. The difference between strict and non-strict inequalities will matter at boundaries.

The word \emph{positive} means strictly greater than zero. \emph{Nonnegative} means greater than or equal to zero, so it includes zero. Likewise, $x\le y$ means $x$ is at most $y$, and $x<y$ means strictly smaller. A chain such as $0\le x\le1$ states two conditions together: $0\le x$ and $x\le1$.

Later we will use square roots. For a nonnegative real $t$, the symbol $\sqrt t$ denotes the \emph{nonnegative} real number whose square is $t$. For example, $\sqrt9=3$, although both $3$ and $-3$ have square nine. We use the usual real-number fact that this nonnegative square root exists and is unique. The expression $\sqrt{x^2}$ therefore equals the nonnegative magnitude of $x$, not always $x$ itself. At $x=-3$ it equals three. We give this magnitude its own notation in the next chapter.

\section{A proof is also an explanation of dependence}
In the parity proof, the two inputs may vary. The witnesses $k$ and $\ell$ vary with them. The final witness for evenness is built from those earlier witnesses. This dependence is part of the argument, not an incidental choice of letters.

Suppose an assistant writes both odd inputs as $2k+1$ using the same $k$. That only describes the special case in which the two inputs agree. The algebra could be correct and the argument still fail to cover the intended claim. A useful question is: ``Which quantities are independent, and which are chosen in response to others?'' Chapter~2 turns this question into a general method.

\begin{lab}{Asking for the missing bridge}
\noindent\textbf{Your first attempt.} Explain to a partner why three odd integers have an odd sum. Use different letters for the three witnesses.

\noindent\textbf{Learner.} ``Here is my argument: three odd numbers have sum $2k+2\ell+2m+3$. I believe the answer is odd, but I cannot see the final representation. Explain only that step.''

\noindent\textbf{Teaching reply.} ``Split $3$ as $2+1$, then factor the terms divisible by two:
\[
2k+2\ell+2m+3=2(k+\ell+m+1)+1.
\]
The quantity in parentheses is an integer, so the result has the defining form of an odd integer.''

\noindent\textbf{Your follow-up.} ``Why is it not enough to say that the expression looks odd? Which exact definition finished the proof?''

\noindent\textbf{Record.} Write the general claim, your original stopping point, and the repaired final step. Then repeat the argument without the reply. The purpose of the exchange is to isolate a small learning need, not to replace your attempt with a complete generated answer.
\end{lab}

\section*{Exercises}
Exercises~1--4 are core; 5--6 are transfer. Full solutions appear in the solution part.
\begin{exercise}\label{ex:1.1}Evaluate $2(x+3)$ and $2x+3$ at $x=4$. Explain the difference in words.\end{exercise}
\begin{exercise}\label{ex:1.2}Prove that the sum of an even integer and an odd integer is odd. State the integer that witnesses the conclusion.\end{exercise}
\begin{exercise}\label{ex:1.3}An assistant claims that the product of two odd integers is odd because $(2k+1)^2$ is odd. Identify the coverage problem and give a complete proof.\end{exercise}
\begin{exercise}\label{ex:1.4}Let $a_1=2$, $a_2=-1$, and $a_3=4$. Expand and evaluate $\sum_{j=1}^3(2a_j+1)$.\end{exercise}
\begin{exercise}\label{ex:1.5}Find every real solution of $x(x-2)=0$. Explain why division by $x$ alone is not a complete method. You may use the arithmetic fact that a product of real numbers is zero only when at least one factor is zero.\end{exercise}
\begin{exercise}\label{ex:1.6}You check a numerical claim for inputs $1$ through $10$. Construct a formula that agrees with the claim $f(n)=0$ on all these inputs but fails at $11$. Hint: use a product with one factor for each checked input.\end{exercise}

\section*{Before moving on}
You should be able to explain the parity proof without remembering its precise wording. You should also be able to name the difference between a definition, a checked example, and a statement about every allowed input. Those distinctions will remain necessary when the examples become research-level objects.


% ===== Week 2 =====

\chapter{Statements, choices, and assumptions}\label{ch:logic}
\goal{Read the exact strength of a claim. Track which choices may depend on which inputs. Negate statements correctly and use counterexamples to ask better questions.}

\section{An implication is a conditional promise}
The statement ``If an integer is divisible by four, then it is even'' has a condition and a conclusion. We abbreviate their logical form as $P\Rightarrow Q$, read ``$P$ implies $Q$.'' Here $P$ is divisibility by four and $Q$ is evenness. Divisibility by four means $n=4k$ for an integer $k$. Then $n=2(2k)$, which proves evenness.

The implication makes no promise about a number that is not divisible by four. Such a number may be even, like $6$, or odd, like $7$. A common reading error is to treat the condition as necessary for the conclusion. It is sufficient, but may not be necessary.

The \emph{converse} is $Q\Rightarrow P$. In this example it says every even integer is divisible by four, which is false because $6$ is even but not divisible by four. A single counterexample is enough: it satisfies the proposed assumption and violates the proposed conclusion.

The statement $P\Leftrightarrow Q$, read ``$P$ if and only if $Q$,'' promises both implications. To prove it, check each direction separately. The statement ``$n$ is even if and only if $n^2$ is even'' is one we will prove in Chapter~4. Its two directions require different choices of representation.

\subsection*{A theorem is a rule you must qualify to use}
To use ``divisible by four implies even'' at $n=20$, first verify the input condition: $20=4\cdot5$, with five an integer. The theorem then gives evenness. There are two distinct pieces here: a general implication already justified and evidence that this particular input satisfies its hypothesis. Merely naming a theorem does not provide the second piece.

``Sufficient'' means enough to obtain the conclusion. Divisibility by four is sufficient for evenness. ``Necessary'' means required whenever the property in question holds. Evenness is necessary for divisibility by four. Neither sentence says the two properties coincide. At six, evenness holds while divisibility by four fails.

To test an implication, look for an input where its hypothesis is true and its conclusion false. An input outside the hypothesis does not refute it. Seven is not a counterexample to ``divisible by four implies even'': it does not satisfy the required input condition. This simple distinction is the difference between asking an AI for random tests and asking it to search for a genuine counterexample.

\section{And, or, and not}
The statement $P\land Q$ means both $P$ and $Q$ hold. The statement $P\lor Q$ means at least one holds; both are allowed. Mathematical ``or'' is inclusive unless the author explicitly says otherwise. The symbol $\neg P$ means $P$ is false.

Consider a number that is even or divisible by three. The number $6$ qualifies for both reasons. To disprove this property, one must show that neither condition holds. Thus
\[
\neg(P\lor Q)\quad\Longleftrightarrow\quad (\neg P)\land(\neg Q).
\]
By contrast, to disprove ``even and divisible by three,'' it is enough for either condition to fail:
\[
\neg(P\land Q)\quad\Longleftrightarrow\quad (\neg P)\lor(\neg Q).
\]
These rules are not arbitrary symbol manipulations. They describe exactly what would make the original promises fail.

An implication fails precisely when its assumption holds and its conclusion fails. Therefore
\[
\neg(P\Rightarrow Q)\quad\Longleftrightarrow\quad P\land\neg Q.
\]
This is the specification of a counterexample to a conditional theorem.

\subsection*{Read the grouping before changing the symbols}
The letters $P,Q$ stand for whole statements, not numbers. If $P$ is ``$n$ is even'' and $Q$ is ``$n$ is divisible by three,'' then $P\land Q$ is the single statement that both requirements hold. Parentheses say which statements belong together, just as they group numerical operations.

For $n=4$, $P$ is true and $Q$ is false. Thus $P\lor Q$ is true but $P\land Q$ is false. For $n=5$, both are false, so both combinations are false. For $n=6$, both are true. These three checks give concrete meaning to the earlier negation rules. They are not proofs of every logical identity, but they help prevent a symbol from replacing its meaning.

A universal statement such as ``every allowed input has the property'' is not a command to list all inputs. It specifies the coverage a proof must have. An existence statement asks for at least one input with the property; it does not promise only one. We now introduce symbols for these two kinds of coverage.

\section{The order of choices changes the theorem}
The symbol $\forall$ means ``for every,'' and $\exists$ means ``there exists.'' Write $x\in\RR$ to say that $x$ belongs to the real numbers. Compare
\begin{align*}
&\forall x\in\RR\ \exists y\in\RR:\quad y>x,\tag{A}\\
&\exists y\in\RR\ \forall x\in\RR:\quad y>x.\tag{B}
\end{align*}
In (A), someone supplies $x$ and we may respond with a suitable $y$. Taking $y=x+1$ works. The response depends on the input. In (B), we must first choose one $y$ that exceeds every real number. Once we choose it, the input $x=y+1$ defeats it. The two sentences contain the same inequality but make different promises.

Read a quantified statement as a sequence of choices. A later existential choice may use earlier universally chosen inputs. It cannot secretly depend on a later input when it was required to be fixed beforehand. This is one of the most useful questions to ask when a proof involving constants becomes confusing.

\subsection*{Tolerance and the inequality moves it needs}
A \emph{tolerance} is a positive amount of error we are prepared to allow. The letter $\eps$, read ``epsilon,'' is often used for it, but it is just a name. It might stand for $1/10$, $1/1000$, or any other positive real number. ``For every positive $\eps$'' includes all of those choices; it does not refer to one unspecified fixed decimal.

Before using fractions in an inequality, check signs. Adding the same number to both sides preserves order. Multiplying or dividing by a positive number preserves order too. Multiplying by a negative number reverses order: from $2<5$ we get $-2>-5$. For $0<a<b$, divide $a<b$ by the positive product $ab$ to obtain $1/b<1/a$. This explains, rather than merely quotes, the rule that reciprocals reverse order for positive inputs.

For example, to make $1/N<1/100$, choosing $N=101$ works, while $N=100$ gives equality and fails the strict inequality. For a general tolerance, we will choose an integer beyond $1/\eps$. Integers are not bounded above on the real number line: there is an integer beyond any prescribed real number. This accepted property is called \emph{Archimedean}; the name is not an extra method to learn.

\begin{example}
The statement ``For every positive real $\eps$, there is a positive integer $N$ with $1/N<\eps$'' allows $N$ to depend on $\eps$. A smaller requested tolerance generally needs a larger integer. The existence of an integer exceeding any specified real number is the Archimedean property of the usual number systems, which we take as a starting fact. Choose $N>1/\eps$. Because both numbers are positive, reciprocation reverses their order and gives $1/N<\eps$.
\end{example}

\pauseq{Can one positive integer $N$ satisfy $1/N<\eps$ for every positive $\eps$?}{No. After $N$ is fixed, take $\eps=1/(2N)$. Then $1/N=2\eps>\eps$. The failed claim demands one universal choice instead of a choice adapted to each tolerance.}

\section{Negation as the design of a failure}
To negate a universal statement, find one failing input. To negate an existence statement, rule out every candidate:
\[
\neg(\forall x\ P(x))\Longleftrightarrow\exists x\ \neg P(x),
\qquad
\neg(\exists x\ P(x))\Longleftrightarrow\forall x\ \neg P(x).
\]
For instance, the negation of (A) is
\[
\exists x\in\RR\ \forall y\in\RR:\quad y\le x.
\]
It claims that some real number is an upper bound for all real numbers. Notice that $y>x$ became $y\le x$, not $y<x$. Equality belongs to the failure of a strict inequality.

Now take a more elaborate statement: ``For every student, there is an acceptable task.'' Its negation is ``There is a student with no acceptable task.'' This says nothing yet about assigning different tasks to different students. Two students might have the same single acceptable task. Thus individually feasible choices need not combine into a globally feasible assignment. This will become our first long theorem.

Here is a slower negation of ``for every integer $n$, there is an even integer $m$ greater than $n$.'' Fixing $n$, the desired response $m$ must satisfy \emph{both} ``$m$ is even'' and ``$m>n$.'' A response fails when at least one requirement fails: it is not even, or it is at most $n$. For the whole promise to fail, there must be one input $n$ for which \emph{every} candidate response fails. Thus the negation is: there exists an integer $n$ such that every integer $m$ is either not even or at most $n$. We reversed the two choices and then described failure of their final test.

This is a useful order of work: negate the plain-language promise first, then write its symbols. Asking an assistant to ``flip the quantifiers'' without explaining what counts as failure can leave a correct-looking but misunderstood sentence.

\section{Contraposition and vacuous truth}
The \emph{contrapositive} of $P\Rightarrow Q$ is $\neg Q\Rightarrow\neg P$. These are equivalent in ordinary classical logic. If $Q$ fails, then $P$ could not have held, because the original implication would have forced $Q$. Conversely, if the contrapositive holds and $P$ is true, $Q$ cannot be false.

This is different from the converse. ``Divisible by four implies even'' is equivalent to ``not even implies not divisible by four.'' It is not equivalent to ``even implies divisible by four.''

An implication with an impossible assumption has no counterexample. For example, ``Every integer that is both even and odd is larger than $100$'' is true in this limited logical sense because no integer is both even and odd. We call this \emph{vacuous truth}. It does not tell us that such an integer exists. If a model produces a theorem whose assumptions cannot simultaneously hold, a proof may be formally correct but uninformative for the intended application.

\section{Magnitude and ranges of real inputs}\label{sec:magnitude}
The \emph{absolute value} $|x|$ is the distance of the real number $x$ from zero. Thus $|3|=3$, $|-3|=3$, and $|0|=0$. Explicitly, $|x|=x$ for $x\ge0$, and $|x|=-x$ for $x<0$. The second formula gives a nonnegative result because the negative of a negative number is positive.

Changing the sign of a number does not change its distance from zero, so $|-x|=|x|$. Multiplication obeys $|cx|=|c|\,|x|$. To check this rule, if both factors are nonnegative, both sides equal $cx$. If both are negative, their product is positive and $|c|\,|x|=(-c)(-x)=cx$. If their signs differ, the product is negative and both sides equal $-cx$. If either factor is zero, both sides are zero. Thus all sign cases are covered. We will use this scaling rule when bounding errors.

For $b\ge0$, the statement $|x|\le b$ says that $x$ lies between $-b$ and $b$, including the endpoints. Thus $|x|\le2$ means $-2\le x\le2$. Its negation, $|x|>2$, means $x<-2$ or $x>2$. This is the two-sided conclusion needed when an equation or inequality involves a square.

Squaring preserves order for nonnegative numbers. If $0\le a<b$, then
\[
 b^2-a^2=(b-a)(b+a)>0,
\]
so $a^2<b^2$. But we cannot square an arbitrary signed inequality without checking this condition: $-3<2$ while $9>4$. Since $x^2=|x|^2$, the condition $x^2>4$ is equivalent to $|x|>2$, not merely to $x>2$.

An \emph{interval} names a range on the real line. We write $[a,b]$ for all $x$ with $a\le x\le b$, and $(a,b)$ for all $x$ with $a<x<b$. A square bracket includes its endpoint; a round bracket excludes it. The notation $[0,\infty)$ means all nonnegative reals; the infinity symbol says there is no upper endpoint, not that infinity is another real number. Parentheses may also denote an ordered pair; the surrounding sentence tells you whether we mean a range or a pair of entries.

\begin{principle}{Habit 2: keep a statement contract}
Before investigating, write the domain, hypotheses, conclusion, and order of choices. After every proposed repair or generalization, compare with that contract. Ask which assumption was weakened or strengthened and whether the new conclusion is stronger or weaker. Truth alone does not establish fidelity to the original question.
\end{principle}

\begin{lab}{A constant that knows too much}
\noindent\textbf{Target.} Decide whether $\forall x>0\ \exists C>0:\ x\le C$ and $\exists C>0\ \forall x>0:\ x\le C$ mean the same thing.

\noindent\textbf{Deliberately flawed assistant reply.} ``Both are true: just set $C=x$.''

\noindent\textbf{Learner.} ``At what stage is $C$ chosen in each sentence? Your choice uses $x$. Is that dependency permitted in the second sentence?''

\noindent\textbf{Repair.} The first sentence is true using $C=x$. The second is false: for any proposed $C$, choose $x=C+1$. This positive input violates $x\le C$.

\noindent\textbf{Transfer.} Replace ``a constant'' by ``an algorithm setting,'' ``a task assigned to everyone,'' or ``an error tolerance.'' In each version distinguish a setting adapted to each input from one setting that must work for all inputs. Write the quantifiers before discussing plausibility.
\end{lab}

\section*{Exercises}
\begin{exercise}\label{ex:2.1}State the converse and contrapositive of ``If $x=3$, then $x^2=9$,'' for real $x$. Decide which are true.\end{exercise}
\begin{exercise}\label{ex:2.2}Negate: ``For every integer $n$, there is an integer $m$ with $m>n$ and $m$ even.'' Preserve the domain and negate the conjunction correctly.\end{exercise}
\begin{exercise}\label{ex:2.3}Decide whether $\forall x\in\RR\ \exists y\in\RR:\ x+y=0$ and $\exists y\in\RR\ \forall x\in\RR:\ x+y=0$ are true. Supply a witness or a refutation.\end{exercise}
\begin{exercise}\label{ex:2.4}Refute: ``If a real number satisfies $x^2>4$, then $x>2$.'' Repair the conclusion without restricting the domain.\end{exercise}
\begin{exercise}\label{ex:2.5}Two students each accept at least one task, but cannot receive distinct acceptable tasks. Give an example. What stronger condition would detect its failure?\end{exercise}
\begin{exercise}\label{ex:2.6}For each positive $\eps$, choose a positive integer $N$ such that every integer $n\ge N$ satisfies $1/n<\eps$. Explain both why $N$ may depend on $\eps$ and why it must not depend on the later input $n$.\end{exercise}

\section*{Before moving on}
Translate a quantified sentence into a sequence of choices, and translate its negation into an attempted failure. In later chapters, this will clarify convergence, uniform bounds, existence, and formal theorem statements without adding a new logical language each time.


% ===== Week 3 =====

\chapter{Objects, maps, and information}\label{ch:maps}
\goal{Use sets to specify allowed objects and functions to describe transformations. Prove a statement by unpacking its definition. Distinguish information lost by a map from information lost only outside its actual inputs.}

\section{Sets express the boundary of a question}
A set is a collection whose membership matters. The set $A=\{1,3,5\}$ has three elements. Repeating a member does not create a new element: $\{1,1,3\}=\{1,3\}$. Order does not matter either. By contrast, an ordered pair $(1,3)$ is different from $(3,1)$. The distinction between a collection and an ordered arrangement is part of the mathematical model.

Write $x\in A$ when $x$ is a member and $x\notin A$ otherwise. The empty set $\varnothing$ has no elements. The notation $A\subseteq B$ says that every member of $A$ is a member of $B$. To prove this containment, take an arbitrary $x\in A$ and show $x\in B$. To prove equality of sets, prove containment in both directions.

The union $A\cup B$ contains elements in at least one set. The intersection $A\cap B$ contains elements in both. The difference $A\setminus B$ contains elements in $A$ but not in $B$. For $A=\{1,2,3\}$ and $B=\{3,4\}$, these are $\{1,2,3,4\}$, $\{3\}$, and $\{1,2\}$ respectively.

The notation $\{x\in A:P(x)\}$ selects the members of $A$ satisfying a property $P$. For example, $\{n\in\ZZ:n\text{ is even}\}$ is the set of even integers. Read the colon as ``such that.''

\subsection*{Why a membership proof proves an entire set equality}
The statements $1\in\{1,2\}$ and $\{1\}\subseteq\{1,2\}$ have different subjects. The first talks about one number. The second talks about a set whose members must all occur in another set. In particular, ``$x$ is a member'' and ``the one-element set $\{x\}$ is contained'' are related but differently written claims.

Two sets are equal precisely when they have the same members. To prove $D=E$, it is therefore enough to show that anything in $D$ is in $E$, and anything in $E$ is in $D$. Start the first direction with an arbitrary $x\in D$; do not assume at the outset that it also belongs to $E$. Producing that second membership is the work of the argument.

If the starting set is empty, there is no member to defeat the proposed containment. This is why $\varnothing\subseteq E$ for every set $E$. The sentence ``take an arbitrary member'' is conditional on a member being supplied; it does not assert that the set is nonempty. This is the set version of vacuous truth.

For a concrete preview of the next identity, use $A=\{1,2,3\}$, $B=\{2,4\}$, and $C=\{3,4\}$. First $B\cup C=\{2,3,4\}$; retaining members also in $A$ leaves $\{2,3\}$. On the other side, $A\cap B=\{2\}$ and $A\cap C=\{3\}$; their union is again $\{2,3\}$. The proof below explains why this agreement persists for arbitrary sets.

\par\noindent\begin{minipage}{\linewidth}
\begin{proposition}
For any sets $A,B,C$, we have $A\cap(B\cup C)=(A\cap B)\cup(A\cap C)$.
\end{proposition}
\end{minipage}\par

\begin{proof}
Suppose $x\in A\cap(B\cup C)$. Then $x\in A$ and $x$ belongs to $B$ or to $C$. In the first case $x\in A\cap B$; in the second $x\in A\cap C$. In either case $x$ belongs to their union. This proves one containment.

Conversely, suppose $x\in(A\cap B)\cup(A\cap C)$. In the first case $x\in A$ and $x\in B$; in the second $x\in A$ and $x\in C$. Both cases imply $x\in A$ and $x\in B\cup C$. Thus $x\in A\cap(B\cup C)$, proving the reverse containment.
\end{proof}
The proof did not depend on what the elements were. Its content was the pattern of membership. This is a first example of a structure that survives a change of subject matter.

\subsection*{Counting members and forming pairs}
For a finite set $S$, the notation $|S|$ means the number of its members. For example, if $S=\{2,4,6\}$, then $|S|=3$, not twelve. We count members, not their sum. The empty set has size zero. The word \emph{cardinality} is another name for this size. Bars around a real number mean its absolute value; bars around a finite set mean its cardinality. The kind of object inside the bars tells us which meaning is being used.

Two sets are \emph{disjoint} if they have no member in common. Counting a disjoint union adds the two sizes. If the sets overlap, simply adding their sizes counts the common members twice. For example, $|\{1,2\}|+|\{2,3\}|=4$ but $|\{1,2,3\}|=3$.

The \emph{Cartesian product} $A\times B$ is the set of ordered pairs $(a,b)$ with $a\in A$ and $b\in B$. If $A=\{1,2\}$ and $B=\{p,q\}$, it contains $(1,p),(1,q),(2,p),(2,q)$. A rule whose input consists of two points of $X$ can therefore be described as a function on $X\times X$. The multiplication-looking sign forms pairs of objects; it does not multiply their numerical values.

\section{A function includes its input and output sets}
A function $f:A\to B$ assigns exactly one output $f(x)\in B$ to each input $x\in A$. The set $A$ is the domain; $B$ is the codomain. Not every element of $B$ needs to occur as an output. The outputs that actually occur form the image
\[
f(A)=\{f(x):x\in A\}.
\]
For $f:\RR\to\RR$ given by $f(x)=x^2$, the image is the nonnegative real numbers, written $[0,\infty)$. Its codomain is all real numbers. We could define a different function with the same rule but codomain $[0,\infty)$; that distinction changes whether every target is attained.

The identity map $\id_A:A\to A$ sends $x$ to itself. If $f:A\to B$ and $g:B\to C$, their composition $g\circ f:A\to C$ applies $f$ first and $g$ second:
\[
(g\circ f)(x)=g(f(x)).
\]
The order is easy to reverse when reading notation. Say ``$g$ after $f$'' aloud until the direction is familiar.

\begin{center}
\begin{tikzpicture}[>=Stealth,node distance=1.65in]
\node (a) {$A$};\node[right=of a] (b) {$B$};\node[right=of b] (c) {$C$};
\draw[->] (a)--node[above]{$f$}(b);\draw[->](b)--node[above]{$g$}(c);
\draw[->,bend right=22](a)to node[below]{$g\circ f$}(c);
\end{tikzpicture}
\end{center}

\subsection*{Follow two concrete maps before composing general ones}
Let $f(x)=x+1$ and $g(y)=2y$, with all inputs and outputs real. Starting at three, $f$ first gives four and $g$ then gives eight. Thus $(g\circ f)(3)=8$. Starting instead with $g$ gives six, and then $f$ gives seven, so $(f\circ g)(3)=7$. The two orders need not agree.

The general formulas are $(g\circ f)(x)=2(x+1)$ and $(f\circ g)(x)=2x+1$. This connects composition to the bracket calculation from Chapter~1. In $g(f(x))$, evaluate the inner parentheses first. Before any evaluation, check that the output of the inner function is an allowed input of the outer one.

To call a map \emph{reversible}, we mean that a second map can undo it in both directions. More precisely, an inverse of $f:A\to B$ is a map $h:B\to A$ with $h(f(a))=a$ for every $a\in A$ and $f(h(b))=b$ for every $b\in B$. For $f(x)=x+1$, the map $h(y)=y-1$ works: $(x+1)-1=x$ and $(y-1)+1=y$. A recipe that undoes one selected input is not yet an inverse for a whole function.

\section{Injective, surjective, and reversible}
A function is \emph{injective} when distinct inputs never merge. Its proof-friendly definition is
\[
\forall x,y\in A,\qquad f(x)=f(y)\Rightarrow x=y.
\]
A function is \emph{surjective} when every element of its codomain is an output:
\[
\forall b\in B\ \exists a\in A:\qquad f(a)=b.
\]
A function with both properties is \emph{bijective}. An injective function preserves the distinction between inputs. A surjective function reaches every specified target. These are different obligations.

Consider the finite map $f:\{a,b\}\to\{1,2,3\}$ with $f(a)=1$ and $f(b)=2$. Its inputs do not merge, so it is injective. It misses the target three, so it is not surjective. Conversely, the map $g:\{a,b,c\}\to\{1,2\}$ with outputs $1,1,2$ is surjective but not injective. Its two preimages of one prevent a uniquely determined reverse assignment.

Why does bijectivity supply an inverse? Given a target $b$, surjectivity supplies at least one preimage $a$. Injectivity says that two such preimages would have to be equal. We can therefore define $h(b)$ to be that unique preimage. By construction $f(h(b))=b$, and for any input $a$, the unique preimage of $f(a)$ is $a$ itself, giving $h(f(a))=a$. Both inverse identities have now been checked.

The square map on all real inputs is not injective because $f(1)=f(-1)$. Restricted to nonnegative inputs and nonnegative outputs, it is bijective: each nonnegative real has exactly one nonnegative square root. Here we use the usual square-root property of real numbers. Changing the domain removes the collision, and changing the codomain removes unreachable negative targets.

\par\noindent\begin{minipage}{\linewidth}
\begin{theorem}\label{thm:injcomp}
If $f:A\to B$ and $g:B\to C$ are injective, then $g\circ f$ is injective.
\end{theorem}
\end{minipage}\par

\noindent\textit{Plan.} Equality at the final output should let us recover equality at the intermediate stage, and then at the input stage.
\begin{proof}
Take any $x,y\in A$ and assume $g(f(x))=g(f(y))$. The elements $f(x)$ and $f(y)$ belong to $B$, so injectivity of $g$ applies and gives $f(x)=f(y)$. Injectivity of $f$ then gives $x=y$. This is exactly the implication needed for injectivity of the composition.
\end{proof}

An assistant may describe the two steps as ``cancel $g$ and then cancel $f$.'' This is acceptable shorthand only after the hypotheses are clear. There is no general rule that arbitrary functions can be canceled.

\par\noindent\begin{minipage}{\linewidth}
\begin{proposition}\label{thm:compconverse}
If $g\circ f$ is injective, then $f$ is injective. The function $g$ need not be injective on all of $B$.
\end{proposition}
\end{minipage}\par

\begin{proof}
Suppose $f(x)=f(y)$. Applying the same function $g$ to equal objects gives $g(f(x))=g(f(y))$. Injectivity of $g\circ f$ then gives $x=y$, so $f$ is injective.

For the second assertion, take $A=\{0,1\}$, $B=\{0,1,2\}$, $C=\{0,1\}$, and define
\[
f(0)=0,\quad f(1)=1,\qquad g(0)=0,\quad g(1)=g(2)=1.
\]
The composition sends $0$ to $0$ and $1$ to $1$, so it is injective. Yet $g$ is not injective because the distinct inputs $1,2$ have equal output.
\end{proof}
The unused input $2$ explains the failure. The composition only tests $g$ on the image $f(A)$, not on the whole set $B$. A counterexample can reveal the exact correction: $g$ is injective when restricted to $f(A)$.

\section{Images and preimages answer different questions}
For $S\subseteq A$, the image $f(S)$ collects the outputs generated by $S$. For $T\subseteq B$, the preimage
\[
f^{-1}(T)=\{x\in A:f(x)\in T\}
\]
collects inputs landing in $T$. The notation does not require an inverse function. For the square map, $f^{-1}(\{1\})=\{-1,1\}$, even though the square map is not invertible on $\RR$.

For any $S,T\subseteq B$,
\[
f^{-1}(S\cap T)=f^{-1}(S)\cap f^{-1}(T).
\]
Read the left side as ``inputs whose outputs land in the intersection.'' We can justify the identity one membership condition at a time:
\begin{align*}
 x\in f^{-1}(S\cap T)
 &\Longleftrightarrow f(x)\in S\cap T &&\text{definition of preimage}\\
 &\Longleftrightarrow f(x)\in S\text{ and }f(x)\in T
 &&\text{definition of intersection}\\
 &\Longleftrightarrow x\in f^{-1}(S)\text{ and }x\in f^{-1}(T)
 &&\text{definition of preimage}\\
 &\Longleftrightarrow x\in f^{-1}(S)\cap f^{-1}(T)
 &&\text{definition of intersection}.
\end{align*}
Each double arrow is valid in both directions, so the chain establishes both containments. The input $x$ was arbitrary; consequently the two sets have exactly the same members. No inverse function was used. Chapter~16 will express these same membership conditions in Lean.

\begin{principle}{Habit 3: identify the information actually used}
A statement about the whole codomain may be stronger than a statement about reachable outputs. A statement about labeled objects may be stronger than a statement about their structure. Before asking for a generalization, locate the part of the object that the proof actually visits.
\end{principle}

\begin{lab}{Repairing a converse precisely}
\noindent\textbf{Learner.} ``I understand why injective maps compose. Is the converse true? Please separate the claims about $f$ and $g$ rather than answering yes or no.''

\noindent\textbf{Task.} Produce the finite counterexample above without consulting it, and explain why $g$ still behaves injectively on the image of $f$.

\noindent\textbf{Follow-up.} ``What additional assumption on $f$ would make injectivity of $g$ on its entire domain follow?''

\noindent\textbf{Model conclusion.} Surjectivity of $f$ is enough. For any $b_1,b_2\in B$ with $g(b_1)=g(b_2)$, choose $a_1,a_2$ with $f(a_1)=b_1$ and $f(a_2)=b_2$. The composition is injective, so $a_1=a_2$; applying $f$ gives $b_1=b_2$. Each choice and each hypothesis has a visible role.
\end{lab}

\section*{Exercises}
\begin{exercise}\label{ex:3.1}For $A=\{1,2,4\}$ and $B=\{2,3,4\}$, find $A\cup B$, $A\cap B$, and $A\setminus B$.\end{exercise}
\begin{exercise}\label{ex:3.2}Prove that a composition of two surjective functions is surjective. Choose witnesses in the correct order.\end{exercise}
\begin{exercise}\label{ex:3.3}Classify $f:\ZZ\to\ZZ$, $f(n)=2n$, as injective or surjective. Then change only its codomain to make it bijective.\end{exercise}
\begin{exercise}\label{ex:3.4}Prove $f^{-1}(S\cup T)=f^{-1}(S)\cup f^{-1}(T)$ by explaining membership on both sides.\end{exercise}
\begin{exercise}\label{ex:3.5}Suppose $h:B\to A$ satisfies $h\circ f=\id_A$. Prove $f$ is injective. Suppose instead $f\circ h=\id_B$. Prove $f$ is surjective.\end{exercise}
\begin{exercise}\label{ex:3.6}Is $f(S\cap T)=f(S)\cap f(T)$ always true for subsets $S,T$ of the domain? Prove the always-valid containment and give a counterexample to equality.\end{exercise}


% ===== Week 4 =====

\chapter{How a proof finds its next step}\label{ch:proof}
\goal{Work backward from a target and forward from hypotheses. Understand direct proof, contrapositive, contradiction, and uniqueness as different ways to organize obligations, not as names to memorize.}

\section{Two directions of attention}
A proof has a target and a supply of facts. Working backward asks what would be sufficient to establish the target. Working forward asks what follows from the available facts. A useful intermediate claim connects the two directions.

To prove an integer is even, the target tells us to find a representation $2k$. To prove a function is injective, the target tells us to start with two inputs having equal outputs. To prove existence, we need either a candidate or an argument excluding the possibility that no candidate exists. Much of proof construction is recognizing what the conclusion asks us to produce.

For example, to prove $n^2$ is even when $n$ is even, the hypothesis gives $n=2k$. Squaring gives $n^2=4k^2=2(2k^2)$. The witness required by the target is $2k^2$. The proof works forward from the definition and backward from the desired form at the same time.

\subsection*{A proof state on paper}
We can record a short argument using two columns: facts already available and what still has to be shown. For ``if $n$ is even then $n^2$ is even,'' the starting fact is ``$n$ is an even integer.'' The target is not another fact; it is ``find an integer $w$ such that $n^2=2w$.'' The name $w$ is new because we have not yet found the witness.

Unpacking the hypothesis supplies $n=2k$ with $k$ an integer. Squaring gives $n^2=4k^2$. Factoring gives $n^2=2(2k^2)$. We can now choose $w=2k^2$. It is an integer, and the required equality holds. Only at this point does the target become an established fact.

This distinction prevents a circular argument such as ``since $n^2$ is even, write $n^2=2w$.'' That sentence assumes the very conclusion being sought. Working backward is allowed during exploration---we can ask what would be enough---but the final proof must justify how the needed condition was actually obtained.

\section{When the contrapositive reveals a usable representation}
Now reverse the statement: if $n^2$ is even, must $n$ be even? The equation $n^2=2k$ is less convenient because it does not directly express $n$ as twice an integer. A different approach is to suppose $n$ is not even.

Every integer is either even or odd, but not both. This follows from division by two with remainder $0$ or $1$, a basic integer-arithmetic fact. Thus the contrapositive says: if $n$ is odd, then $n^2$ is odd. Write $n=2k+1$ and calculate
\begin{align*}
n^2&=(2k+1)(2k+1)\\
&=4k^2+4k+1\\
&=2(2k^2+2k)+1.
\end{align*}
The final expression is odd. Therefore $n^2$ cannot be even when $n$ is odd. Combining this with the first direction proves
\[
n\text{ is even}\quad\Longleftrightarrow\quad n^2\text{ is even}.
\]

The strategic lesson is specific: the negation of the desired conclusion supplied a more useful representation than the original hypothesis did. Contraposition is not a ritual to apply whenever a proof feels difficult.

\subsection*{What a contradiction permits, and what it does not}
To prove a claim by contradiction, temporarily suppose the claim is false. Retain the original hypotheses. If their consequences force a statement and its denial, the temporary supposition cannot hold. We then conclude the claim is true. This is ordinary classical reasoning, the setting of this book.

For example, suppose an integer were both even and odd. Then $2k=2\ell+1$ for integers $k,\ell$, and subtraction would give $2(k-\ell)=1$. Division gives $k-\ell=1/2$, impossible because the left side is an integer. The contradiction is that the same number would be an integer and equal to a noninteger. Saying merely ``that looks unlikely'' would not finish the proof.

In the next argument the contradiction is not a false numerical equality. It is a conflict with a deliberately chosen \emph{smallest} denominator. We will explain why that choice exists before using it. A contradiction can violate a selection property just as decisively as it can violate arithmetic.

\section{A complete contradiction proof}
A rational number is a fraction $a/b$ of integers with $b\ne0$. A real number that is not rational is irrational. The positive square root of two, written $\sqrt2$, is the positive real number whose square is two. Its existence belongs to the usual real-number system. We will prove that it is not a rational number.

\par\noindent\begin{minipage}{\linewidth}
\begin{theorem}\label{thm:sqrt2}
The number $\sqrt2$ is irrational.
\end{theorem}
\end{minipage}\par

\noindent\textit{Plan.} If a fraction represented $\sqrt2$, squaring would force both its numerator and denominator to be even. That would allow a smaller positive denominator, contradicting a minimal choice.
\begin{proof}
Suppose $\sqrt2$ were rational. Then it could be written as $a/b$ with integers $a,b$ and $b>0$. Among all such representations, choose one with the smallest positive denominator $b$. This choice is possible because every nonempty set of positive integers has a least element; this is the well-ordering property of the integers.

The choice does not assume there are only finitely many fractions representing the number. Start from any one candidate denominator $b_0$. Among the finite list $1,2,\ldots,b_0$, at least $b_0$ occurs as a denominator. Select the first member of that list that occurs. A still smaller denominator would also be on the list, so this selected one is smallest among all candidates. If the initial representation had a negative denominator, multiplying numerator and denominator by $-1$ would first make it positive without changing the fraction.

Squaring $\sqrt2=a/b$ and multiplying by $b^2$ gives $a^2=2b^2$. Hence $a^2$ is even. By the parity result just proved, $a$ is even, so $a=2c$ for some integer $c$. Substitution gives
\[
(2c)^2=2b^2,\qquad 4c^2=2b^2,\qquad b^2=2c^2.
\]
Thus $b^2$ is even, and the same parity result shows that $b$ is even. Write $b=2d$. Since $b>0$, the integer $d=b/2$ is positive and smaller than $b$. But
\[
\sqrt2=\frac{a}{b}=\frac{2c}{2d}=\frac{c}{d}.
\]
This is a representation with a smaller positive denominator, contradicting the choice of $b$. The assumed rational representation therefore does not exist.
\end{proof}

There are two separate ideas: a local arithmetic fact about parity, and a global choice of the smallest denominator. The contradiction is not simply that numerator and denominator are even; fractions with both even certainly exist. The contradiction is that a representation chosen with minimal denominator could be reduced further.

\pauseq{Where was positivity of the denominator used?}{It made the set of candidate denominators a nonempty set of positive integers with a least member, and it ensured $b/2$ was still positive and strictly smaller. Minimality would not have the same meaning over unrestricted negative denominators.}

\section{Existence is not uniqueness}
The equation $x^2=1$ has solutions, but not a unique real solution. It has exactly two: $1$ and $-1$. ``There is a unique solution'' makes two claims. First, at least one candidate satisfies the specification. Second, any two candidates satisfying it must agree.

Consider the equation $x=(x+1)/2$ over the real numbers. Existence follows by checking $x=1$. For uniqueness, suppose both $x$ and $y$ satisfy the equation. Subtracting their equations gives
\[
x-y=\frac{x-y}{2}.
\]
Multiplying by two and subtracting $x-y$ gives $x-y=0$, so $x=y$. We did not need to discover the solution again. We compared two hypothetical solutions and proved they could not differ. Chapter~12 will generalize this comparison into a theorem about iterative methods.

The two-candidate argument proves \emph{at most one} solution. It would also be valid in a problem with no solutions at all: there would be no two different solutions to compare. That is why the earlier substitution of $x=1$ is indispensable for ``exactly one.'' In general, ``exactly one'' combines ``at least one'' from existence and ``at most one'' from uniqueness.

At the subtraction step, write both equations before combining them:
\[
 x=\frac{x+1}{2},\qquad y=\frac{y+1}{2}.
\]
Subtract the second from the first. On the right the numerators become $(x+1)-(y+1)=x-y$, because the two constant terms cancel. This explains why comparing solutions can be simpler than solving again: shared terms disappear.

When a proof says ``let $x$ be the solution,'' check that existence has already been established or is being used only conditionally. Otherwise the wording may smuggle the main difficulty into an opening sentence.

\section{Reading a long proof as a sequence of smaller promises}
A \emph{lemma} is a supporting result. A \emph{theorem} is a result emphasized as a main conclusion. A numbered \emph{proposition} is another presentation of an assertion with a proof. In logic, the same word can also mean a statement whose truth has not yet been established. The labels describe the author's organization, not different standards of truth.

When a paper invokes a lemma, pause to identify its inputs and output. In the irrationality proof, the parity lemma accepts the fact that $a^2$ is even and returns that $a$ is even. The hypothesis is visibly satisfied. A reference such as ``by a standard result'' is not enough for a beginner unless the result and the matching assumptions are supplied.

A proof can also use temporary assumptions. In a contradiction proof, the false statement being assumed must remain clearly marked. We may deduce consequences from it inside that argument, but we must not export it as an established fact afterward.

\begin{principle}{Habit 4: ask for a local contract, not just more words}
At a difficult step, ask: What is known here? What is the next target? Which exact result is used? Why do its assumptions hold? Which quantity is being constructed? These questions reveal missing mathematics more reliably than asking for an explanation to be ``simpler.''
\end{principle}

\begin{lab}{Finding the actual contradiction}
\noindent\textbf{Deliberately incomplete reply.} ``Assume $\sqrt2=a/b$. Then $a$ and $b$ are even, a contradiction.''

\noindent\textbf{Learner.} ``Why can a numerator and denominator not both be even? For example, $2/4$ is a valid fraction. Identify the missing condition.''

\noindent\textbf{Repair.} Choose a representation with the smallest positive denominator, or explicitly choose one in lowest terms and justify that such a representation exists. Then show that both numerator and denominator being even contradicts that chosen property.

\noindent\textbf{Your output.} A three-part annotation: the assumption introduced for contradiction; the supporting parity result; and the exact statement contradicted. Then explain why checking many decimal digits of $\sqrt2$ would not establish irrationality.
\end{lab}

\section*{Exercises}
\begin{exercise}\label{ex:4.1}Prove directly that if $a,b$ are even integers, then $ab$ is divisible by four.\end{exercise}
\begin{exercise}\label{ex:4.2}Prove by contrapositive: if $3n+1$ is even, then the integer $n$ is odd.\end{exercise}
\begin{exercise}\label{ex:4.3}Prove that if $r$ is rational and $u$ is irrational, then $r+u$ is irrational. State why the difference of two rational numbers is rational.\end{exercise}
\begin{exercise}\label{ex:4.4}Prove that $x=3x-4$ has exactly one real solution. Separate existence and uniqueness.\end{exercise}
\begin{exercise}\label{ex:4.5}An assistant assumes $x^2=2$, then writes $x=\sqrt2$. What is missing over the real domain? Give a corrected classification.\end{exercise}
\begin{exercise}\label{ex:4.6}Explain how the parity proof for $\sqrt2$ changes when the target is $\sqrt8$. Give a full contradiction argument or reduce it to the theorem already proved.\end{exercise}


% ===== Week 5 =====

\chapter{Induction, invariants, and progress}\label{ch:induction}
\goal{Explain why a finite procedure works for every input size. Separate what a process preserves from what forces it to stop. Learn how the choice of an induction statement can expose the right structure.}

\section{From a few sums to every sum}
The first odd number is $1$, the first two sum to $4$, and the first three sum to $9$. The pattern suggests
\[
1+3+\cdots+(2n-1)=n^2\qquad(n\ge1).
\]
The left side grows by the next odd number. The right side changes from $n^2$ to $(n+1)^2$, and
\[
(n+1)^2-n^2=n^2+2n+1-n^2=2n+1.
\]
Both sides therefore change in the same way. This observation is the mechanism behind an induction proof.

Induction uses a starting case and a transition. To prove a statement $P(n)$ for every natural number $n$, prove $P(0)$ and then show that for an arbitrary $k\ge0$, the assumption $P(k)$ implies $P(k+1)$. The second step is conditional. It does not assume the conclusion for all $n$.

\subsection*{Seeing the transition at one actual index}
For $n=3$, the formula says $1+3+5=3^2=9$. The next sum adds seven, so it becomes $9+7=16=4^2$. A transition at a general index $k$ does the same thing: it retains the sum of the first $k$ odd numbers and adds the next odd number $2k+1$.

The \emph{induction hypothesis} is the temporary assumption that the formula works at that particular arbitrary index $k$. The \emph{induction step} shows it then works at $k+1$. We separately prove a \emph{base case} from which these transitions begin. If the base is zero, the transition gives one, then two, then three, and so on. For any specified finite index, that chain reaches it after finitely many transitions.

The sum with no terms is assigned the value zero because it changes nothing when added to another sum. Thus starting our formula at $n=0$ is meaningful: there are no odd numbers to add, and the right side $0^2$ is also zero. We have not introduced a mysterious zeroth odd number.

\par\noindent\begin{minipage}{\linewidth}
\begin{theorem}\label{thm:odds}
For every $n\in\NN$, $\sum_{j=1}^{n}(2j-1)=n^2$, where a sum with no terms is defined to be zero.
\end{theorem}
\end{minipage}\par

\begin{proof}
For $n=0$, both sides are zero. Suppose the formula holds for $n=k$. Split the sum for $k+1$ into the first $k$ terms and its final term:
\begin{align*}
\sum_{j=1}^{k+1}(2j-1)
&=\sum_{j=1}^{k}(2j-1)+(2(k+1)-1)\\
&=k^2+2k+1 &&\text{by the induction hypothesis}\\
&=(k+1)^2.
\end{align*}
This is the formula for $k+1$, completing the induction.
\end{proof}

The proof preserves the original meaning of $k$ throughout. If you do not understand an induction argument, write out its transition explicitly at $k=2$. Then return to the arbitrary $k$ and identify which part did not depend on that particular value.

\section{Why induction is not circular}
Imagine a counterexample existed. Choose the smallest natural number $m$ for which $P(m)$ fails. The starting case shows $m\ne0$, so $m-1$ is a smaller natural number. Minimality gives $P(m-1)$, and the transition gives $P(m)$, a contradiction. This is the \emph{minimal-counterexample} view of induction.

Sometimes a step needs all smaller cases, not just the immediately preceding one. Strong induction allows us to assume $P(0),\ldots,P(k)$ when proving $P(k+1)$. The same smallest-counterexample argument justifies it: every smaller case is true by minimality. In Chapter~9, a graph may split into two parts whose sizes are not predetermined, so we will need this form of induction.

\section{A recursively defined process}
A recurrence defines the next value from earlier values. Let $x_0=0$ and
\[
x_{n+1}=\frac{x_n+1}{2}.
\]
Then $x_1=1/2$, $x_2=3/4$, and $x_3=7/8$. The formula
\[
x_n=1-2^{-n}
\]
fits the examples. Here $2^{-n}=1/2^n$. Prove it by induction: it holds at $n=0$ because $1-1=0$. If it holds at $n$, substitute the proposed formula into the update:
\begin{align*}
 x_{n+1}
 &=\frac{x_n+1}{2} &&\text{definition of the next term}\\
 &=\frac{(1-2^{-n})+1}{2} &&\text{induction hypothesis}\\
 &=\frac{2-2^{-n}}2\\
 &=1-\frac{2^{-n}}2\\
 &=1-2^{-(n+1)}.
\end{align*}
The last equality uses $2^{-n}/2=1/(2^n\cdot2)=1/2^{n+1}$. Notice that $x_{n+1}$ is the name of the next term, whereas $x_n+1$ means one added to the current term.
This proves the exact formula for every finite step. It also identifies the error relative to $1$ as $2^{-n}$. A pattern in the outputs has become an explanation of their behavior.

\section{An invariant answers what cannot change}
An invariant is a property or quantity preserved by every legal step. Suppose a row contains plus and minus signs, and a legal move flips exactly two signs. Associate the numbers $1$ and $-1$ with the signs. Flipping two signs multiplies the product of the row by $(-1)^2=1$, so the product is unchanged.

Starting with all plus signs gives product $1$. A row with exactly one minus sign has product $-1$ and cannot be reached by these moves. This is an impossibility proof: it does not enumerate all possible move sequences. It identifies a barrier every sequence must respect.

An invariant does not guarantee termination. A process may keep moving among states with the same invariant. To show that a process stops, we need a separate progress argument.

\section{Euclid's algorithm: meaning and termination}
A positive integer $d$ divides an integer $a$ if $a=dk$ for some integer $k$. Write $d\mid a$. For nonnegative integers $a,b$ not both zero, their greatest common divisor, $\gcd(a,b)$, is the largest positive integer dividing both. The set of common positive divisors is nonempty because it contains $1$, and finite because at least one of $a,b$ is positive and bounds its divisors.

For $b>0$, integer division gives unique integers $q,r$ with
\[
a=qb+r,\qquad 0\le r<b.
\]
The integer $r$ is the remainder, also written $a\bmod b$. Here $q$ counts the whole copies of $b$ removed from $a$; $r$ is what remains. For example, $17=3\cdot5+2$, so the quotient is three and the remainder is two. The condition $0\le r<b$ identifies the permitted remainder; $17=2\cdot5+7$ is a true equality but not division with its remainder in the required range.

To show that the replacement preserves common divisors, let $d$ divide $a$ and $b$. Then $a=ds$ and $b=dt$ for some integers $s,t$. Substitution gives
\[
 r=a-qb=ds-qdt=d(s-qt).
\]
The bracket is an integer, so $d$ divides $r$. Thus every common divisor of $a,b$ is a common divisor of $b,r$.

For the reverse direction, suppose $b=dt$ and $r=du$. Then
\[
 a=qb+r=qdt+du=d(qt+u),
\]
so $d$ divides $a$ as well. Both directions are needed: only then do we know that the two pairs have \emph{exactly} the same common divisors. Their greatest common divisors consequently agree, giving
\[
\gcd(a,b)=\gcd(b,r).
\]
This is the invariant behind the algorithm. Repeatedly replace $(a,b)$ by $(b,a\bmod b)$ while $b>0$.

For $(a,b)=(84,30)$, the divisions are
\[
84=2\cdot30+24,\qquad30=1\cdot24+6,\qquad24=4\cdot6+0.
\]
The successive pairs are $(84,30)$, $(30,24)$, $(24,6)$, $(6,0)$. At the last pair, the greatest common divisor is $6$, because the positive divisors common to $6$ and $0$ are exactly the divisors of $6$.

Why must the process stop? Its second coordinate is a nonnegative integer and strictly decreases at every step that occurs. There cannot be an infinite strictly decreasing sequence of nonnegative integers. Once the second coordinate is zero, the algorithm stops and returns the first coordinate. The invariant proves the answer is right; the decreasing integer proves an answer is eventually returned.

\begin{principle}{Habit 5: separate preservation from progress}
For an iterative construction, ask two different questions. What property remains true after a step? What quantity forces progress or controls the accumulated error? A preserved quantity may prove correctness without proving termination. A decreasing real quantity may converge without ever reaching zero.
\end{principle}

\pauseq{Does a strictly decreasing nonnegative real quantity prove that a process stops?}{No. The values $1,1/2,1/4,\ldots$ decrease forever and remain positive. The Euclidean termination proof depends on the decreasing quantity being an integer, not merely a real number.}

\begin{lab}{Finding the right induction statement}
\noindent\textbf{Learner.} ``The iterates $x_{n+1}=(x_n+1)/2$ look as though they approach $1$. Instead of asserting convergence, help me find an exact error formula.''

\noindent\textbf{Useful reply.} Subtract the proposed limit from both sides:
\[
x_{n+1}-1=\frac{x_n-1}{2}.
\]
Thus the error is halved each time. Induction gives $x_n-1=2^{-n}(x_0-1)$ for any starting real value $x_0$.

\noindent\textbf{Your check.} Verify the base case and the multiplication step. Then recover the earlier formula when $x_0=0$. The more general error statement explains more and is no harder to prove. This is an example of strengthening a statement to make its structure simpler.
\end{lab}

\section*{Exercises}
\begin{exercise}\label{ex:5.1}Prove $\sum_{j=1}^{n}j=n(n+1)/2$ for $n\ge0$ by induction. Show the algebra in the transition.\end{exercise}
\begin{exercise}\label{ex:5.2}If $y_0=3$ and $y_{n+1}=2y_n$, find and prove a formula for $y_n$.\end{exercise}
\begin{exercise}\label{ex:5.3}Run Euclid's algorithm on $(1071,462)$. Record each quotient, remainder, and the final greatest common divisor.\end{exercise}
\begin{exercise}\label{ex:5.4}Starting with five plus signs and flipping two signs at a time, can you reach a configuration with exactly three minus signs? Justify without listing all moves.\end{exercise}
\begin{exercise}\label{ex:5.5}For the recurrence $x_{n+1}=(x_n+1)/2$ with arbitrary real $x_0$, prove the error formula in the laboratory. Explain why the error never becomes exactly zero when $x_0\ne1$.\end{exercise}
\begin{exercise}\label{ex:5.6}An induction proof uses $P(k)$ to prove $P(k+2)$ and checks only $P(0)$. Which indices has it actually covered? Give the additional starting case needed to cover all natural numbers.\end{exercise}


% ===== Week 6 =====

\chapter{Sameness, symmetry, and structure}\label{ch:structure}
\goal{Recognize when different descriptions represent the same relevant object. Check that a construction does not depend on arbitrary choices. Use symmetry to simplify a problem without silently discarding valid cases.}

\section{Which differences matter?}
The fractions $1/2$ and $2/4$ look different but represent the same rational number. A graph with vertices named Alice, Ben, and Chen may describe the same connection pattern as a graph named $1,2,3$. A necklace does not change when we rotate it, although its written sequence changes if we choose a different starting bead.

In each example, we decide which distinctions to ignore. This decision should be expressed precisely. Otherwise ``these are the same'' can conceal a change to the problem.

A binary relation on a set $A$ specifies which pairs of elements are related. Write $x\sim y$ for one such relation. An \emph{equivalence relation} is a relation with three requirements. \emph{Reflexivity} says each element is related to itself. \emph{Symmetry} says that if $x\sim y$, then $y\sim x$. \emph{Transitivity} says that if $x\sim y$ and $y\sim z$, then $x\sim z$. These are parts of the definition. Calling a new relation an equivalence relation creates three obligations: check every requirement for arbitrary allowed elements, not just a favorite example.

Consider students grouped by their birth month. Each student has the same birth month as themself. If Alice has the same birth month as Ben, Ben has the same birth month as Alice. If Alice and Ben share a birth month and Ben and Chen do too, Alice and Chen share it. The three checks justify regarding students with a common month as one group for this particular comparison. They do not say the students are identical in every other respect.

The relation ``is a friend of'' is generally not transitive, so it is not automatically an equivalence relation. An explanation that says merely ``a relation groups similar things'' is insufficient: similarity may fail the exact conditions needed for a class construction.

\section{Residues: a small model of a large idea}
Fix a positive integer $m$. Say $a\equiv b\pmod m$ when $m$ divides $a-b$. This means that $a$ and $b$ differ by a whole number of periods of length $m$. For $m=4$, we have $1\equiv5\equiv-3\pmod4$.

This is an equivalence relation. Reflexivity follows because $a-a=0=m\cdot0$. If $a-b=mk$, then $b-a=m(-k)$, proving symmetry. If $a-b=mk$ and $b-c=m\ell$, then $a-c=m(k+\ell)$, proving transitivity.

The equivalence class $[a]_m$ consists of all integers equivalent to $a$. Division with remainder shows that every class has one representative among $0,1,\ldots,m-1$. These $m$ classes form the set $\ZZ/m\ZZ$. The notation can initially be read simply as ``integers with differences by multiples of $m$ ignored.'' We do not need the general theory of quotient objects to use this example correctly.

For modulus four, the class $[1]_4$ is the entire set
\[
 \{\ldots,-7,-3,1,5,9,\ldots\}.
\]
The integer one is a \emph{representative}: one member used to name that set. Five is another representative, so $[1]_4=[5]_4$, even though $1\ne5$. Three levels must stay separate: the integer, the set of integers equivalent to it, and the collection of all such sets.

Why do these classes form nonoverlapping groups? If a number $z$ belonged both to the class of $a$ and to the class of $b$, symmetry and transitivity would give $a\sim z\sim b$, hence $a\sim b$. Any number related to $a$ would then be related to $b$, and conversely. The two classes would be identical. Thus classes can be equal or disjoint, but cannot overlap only partially. Reflexivity ensures every number belongs to its own class, so none is omitted.

\section{Well-definedness: does the answer depend on the representative?}
We would like to add classes by adding representatives:
\[
[a]_m+[b]_m=[a+b]_m.
\]
But a class has many representatives. We must show that replacing them does not change the resulting class.

Suppose $a'=a+mk$ and $b'=b+m\ell$. Then
\[
a'+b'=a+b+m(k+\ell).
\]
Thus $a'+b'\equiv a+b\pmod m$, so the two choices give the same class. The operation is \emph{well-defined}: it is genuinely an operation on classes rather than an operation on a hidden choice of their names.

Test this with $[1]_4+[2]_4$. Using representatives one and two gives $[3]_4$. Using five instead of one and $-2$ instead of two gives $[5+(-2)]_4=[3]_4$ again. Choosing five and six gives $[11]_4$, which is the same class as $[3]_4$. The ordinary integer sum may change; the \emph{output class}, which is what the operation promises to produce, does not.

This explains the point of the general calculation. It is not trying to prove $a'+b'=a+b$, which would usually be false. It proves that the two sums differ by a multiple of $m$, exactly the criterion for equal output classes. Always ask what kind of object the output is before asking whether two outputs agree.

Now consider the proposed ordinary-integer output $F([a]_4)=a$. This is not well-defined. The same class can be written $[1]_4=[5]_4$, but the rule would output either $1$ or $5$. We may instead specify the unique representative in $\{0,1,2,3\}$; then the rule is well-defined because a selection convention has been explicitly supplied.

\pauseq{Does the rule $G([a]_4)=[a]_2$ depend on the representative?}{No. If $a'=a+4k$, then $a'-a=2(2k)$, so $a'$ and $a$ determine the same class modulo two. The rule is a legitimate function from classes modulo four to classes modulo two.}

The map in the pause question loses information: $[0]_4$ and $[2]_4$ both become $[0]_2$. Yet it preserves addition. This is a useful distinction. A map can preserve a structure without preserving every distinction between inputs.

\section{Structure-preserving maps}
A structure is a set together with specified operations or relations. The integers have addition, multiplication, and order. A graph has vertices and adjacency. A space with distance has points and a distance function. When learning a new subject, ask what its objects carry in addition to membership in a set.

For objects with addition, a map $T$ preserves addition when
\[
T(x+y)=T(x)+T(y).
\]
This is part of the definition of a linear map in Chapter~10. For a graph, an isomorphism is a bijective renaming of vertices that preserves adjacency in both directions. The general word \emph{isomorphism} means a reversible map preserving the specified structure; what must be preserved depends on the subject.

A symmetry is an isomorphism from an object to itself. Rotating four positions cyclically is a symmetry of their circular arrangement. Label the positions $0,1,2,3$ and set $r(i)=i+1\pmod4$. Applying $r$ four times returns to the original position. Its inverse is rotation by three positions, because $1+3\equiv0\pmod4$.

Symmetries can be composed, and composing two of them gives another symmetry of the same object. This last requirement is \emph{closure}: the operation stays within the chosen collection. Composition is \emph{associative}: $(f\circ g)\circ h=f\circ(g\circ h)$, because both apply $h$, then $g$, then $f$. The identity changes nothing, and every symmetry has an inverse in the collection.

A \emph{group} is a nonempty set with an operation that combines any two of its elements to give another element of the set, is associative, has an identity, and gives each element an inverse. These requirements name the pattern just seen. They do not require that reversing the order of two operations gives the same result; that is an additional property called commutativity.

For the four rotations, write $r_0,r_1,r_2,r_3$, where $r_j$ rotates by $j$ positions. Applying $r_3$ and then $r_2$ rotates by five positions, the same as one position, so $r_2\circ r_3=r_1$. The identity is $r_0$; the inverse of $r_1$ is $r_3$; and $r_2$ is its own inverse. The result of any composition is one of these four rotations, verifying closure in this example. No general group theory is needed to use these particular transformations.

\section{What symmetry permits us to assume}
Consider a nonempty sequence of signs $a_0,\ldots,a_{n-1}$, each equal to $1$ or $-1$. Negating every sign does not change any pair product, because $(-a_i)(-a_j)=a_i a_j$. Therefore a question expressed entirely through pair products can be studied after setting the first sign to $1$: if the first sign is $-1$, negate the sequence.

This reduction has two obligations. First, the transformation must preserve the property under study. Second, every case must be transformable to one satisfying the selected convention. Both hold here. The reduction would fail for a question about the sum of signs equaling a particular positive value, because negation changes the sign of that sum.

Mathematical taste includes noticing when a symmetry reduces irrelevant choices. It also includes refusing a convenient normalization until its invariance has been established.

\subsection*{What ``without loss of generality'' must justify}
The phrase \emph{without loss of generality} sometimes announces a reduction like making the first sign positive. Read it as a promise that no relevant case has been lost, not as permission to impose a convenient extra condition.

Here the proof of that promise has two cases. A sequence already beginning with $1$ requires no change. A sequence beginning with $-1$ can be negated and then begins with $1$. Negation preserves every pair product and can be undone by negating again. Thus proving a property expressed only through those products for first-positive sequences proves it for first-negative sequences too, by translating there and back. If the target property is not preserved, the phrase would be unjustified.

\section{Order is a different structure}
A partial order is a relation that is reflexive, antisymmetric, and transitive. Antisymmetry means that $x\le y$ and $y\le x$ together force $x=y$. This differs from the symmetry required for equivalence relations.

Set inclusion is a partial order. It need not compare every pair: neither $\{1\}$ nor $\{2\}$ contains the other. In a finite collection, an element is \emph{maximal} if no strictly larger element of that collection contains it. A \emph{maximum} contains every element of the collection. For the collection $\{\{1\},\{2\}\}$, both members are maximal and there is no maximum. The distinction will reappear when a matching cannot be enlarged by one obvious step but is not of largest possible size.

\begin{principle}{Habit 6: ask what a change of description preserves}
Before quotienting, normalizing, or identifying objects, specify the relation and the preserved property. Before applying a formula to equivalence classes, check independence of representatives. A short well-definedness check can prevent an entire argument from being about the wrong object.
\end{principle}

\begin{lab}{A definition that secretly depends on a choice}
\noindent\textbf{Deliberately flawed proposal.} ``Define the parity of a class modulo three as the parity of any integer representing it.''

\noindent\textbf{Learner.} ``Test two representatives of the same class. Do $0$ and $3$ give the same answer?''

\noindent\textbf{Diagnosis.} They represent the same class modulo three but have different parity. The proposed rule is not well-defined.

\noindent\textbf{Repair options.} Use a fixed representative from $\{0,1,2\}$ and acknowledge that the resulting function depends on that convention; or use classes modulo an even number, for which all representatives have the same parity. These are different mathematical constructions. Ask the assistant which one matches the original purpose rather than accepting either repair automatically.
\end{lab}

\section*{Exercises}
\begin{exercise}\label{ex:6.1}List the elements of $[-2]_5$ between $-12$ and $12$. Which representative in $\{0,1,2,3,4\}$ names the same class?\end{exercise}
\begin{exercise}\label{ex:6.2}Prove multiplication of residue classes is well-defined: $[a]_m[b]_m=[ab]_m$.\end{exercise}
\begin{exercise}\label{ex:6.3}For four positions, show rotation by two is its own inverse. Explain why that does not make it the identity.\end{exercise}
\begin{exercise}\label{ex:6.4}Is ``has the same number of elements'' an equivalence relation on finite sets? Explain all three conditions.\end{exercise}
\begin{exercise}\label{ex:6.5}For positive integers $m,d$, prove the rule $[a]_m\mapsto[a]_d$ is well-defined exactly when $d$ divides $m$.\end{exercise}
\begin{exercise}\label{ex:6.6}Give a finite family of sets with two maximal members and no maximum. Then add one set to the family so that it has a maximum.\end{exercise}


% ===== Week 7 =====

\chapter{Graphs: counting, paths, and extremal choices}\label{ch:graphs}
\goal{Model connections without coordinates. Count one collection in two ways. Use shortest or longest objects to force useful structure. Prepare the language needed for matching and for an actual repository theorem.}

\section{A graph is a pattern of connections}
A finite simple undirected graph has a finite set $V$ of vertices and a set $E$ of edges. Each edge is an unordered pair of distinct vertices. ``Simple'' excludes loops and repeated edges. Coordinates in a drawing are not part of this definition: moving a vertex on the page does not change which pairs are joined.

Two vertices are neighbors when an edge joins them. The degree $\deg(v)$ is the number of neighbors of $v$. A graph with vertices $a,b,c,d$ and edges $\{a,b\},\{b,c\},\{c,a\},\{c,d\}$ has degrees $2,2,3,1$.
\begin{center}
\begin{tikzpicture}[every node/.style={circle,draw=ink,minimum size=6mm,inner sep=1pt},line width=.7pt]
\node(a)at(0,0){$a$};\node(b)at(2,0){$b$};\node(c)at(1,1.5){$c$};\node(d)at(3,1.5){$d$};
\draw(a)--(b)--(c)--(a);\draw(c)--(d);
\end{tikzpicture}
\end{center}

\section{Double counting begins with one precise collection}
Let us first count the displayed graph in a deliberately redundant way. Edge $ab$ contributes the two records $(a,ab)$ and $(b,ab)$. Edge $bc$ contributes $(b,bc)$ and $(c,bc)$. Edges $ca$ and $cd$ contribute two records each as well. There are eight records altogether.

Now reorganize those same records by their vertex. There are two beginning with $a$, two with $b$, three with $c$, and one with $d$. Their total is $2+2+3+1=8$. Nothing has been counted differently; the records have only been grouped differently. The general theorem names this common collection explicitly so that we can check there are neither omissions nor duplicate records within either grouping.

The notation $\sum_{v\in V}$ means add one term for every vertex in $V$, rather than for a consecutive integer index. Finite sums may be grouped and reordered by the arithmetic rules from Chapter~1.

\par\noindent\begin{minipage}{\linewidth}
\begin{theorem}\label{thm:handshake}
In a finite simple undirected graph, $\sum_{v\in V}\deg(v)=2|E|$, where $|E|$ is the number of edges.
\end{theorem}
\end{minipage}\par

\begin{proof}
Count pairs $(v,e)$ in which $v$ is an endpoint of the edge $e$. Fixing a vertex $v$ gives $\deg(v)$ such pairs, so counting by vertices gives the sum on the left. Fixing an edge gives exactly two pairs, one for each endpoint, so counting by edges gives $2|E|$. Both counts refer to the same collection, so they are equal.
\end{proof}

As a consequence, the number of vertices of odd degree is even. Split the degree sum into its even terms and odd terms. The even terms sum to an even integer. The entire sum is even, so the sum of the odd terms is even. A sum of $k$ odd integers has the same parity as $k$, because it can be written as twice an integer plus $k$. Hence $k$ is even.

This is more informative than a numerical observation about one drawing. No rearrangement of the graph can produce exactly one odd-degree vertex under the stated rules.

\section{The pigeonhole principle and the average}
If $n$ objects occupy $m$ boxes and every box contains at most $k$ objects, then $n\le mk$: add the bounds for all boxes. Taking the contrapositive, if $n>mk$, some box contains more than $k$ objects. This is the pigeonhole principle in a quantitative form.

For real numbers $a_1,\ldots,a_n$ with $n>0$, at least one is at least their average $\bar a=(a_1+\cdots+a_n)/n$. Otherwise each would be strictly less than $\bar a$, and adding those strict inequalities would give a total smaller than $n\bar a$, contradicting the definition. Similarly, at least one is at most the average.

The word ``some'' is essential. An average does not control every entry. A graph with $m$ edges and $n>0$ vertices therefore has a vertex of degree at least $2m/n$, but not necessarily all vertices have that degree. Chapter~13 will turn the same observation into a method for proving existence through random choices.

\section{Walks, paths, and shortest objects}
A walk is a sequence of vertices in which consecutive vertices are joined by edges. It may repeat vertices. A path is a walk with no repeated vertices. Its length is the number of edges used, one less than the number of listed vertices.

If any walk joins two vertices, a shortest walk between them has no repeated vertex. Indeed, if the same vertex occurs twice, delete the part of the walk between those occurrences. The remaining sequence is still a walk with the same endpoints and has fewer edges, contradicting shortestness. Thus connected vertices admit a path.

For a concrete deletion in the displayed graph, consider the walk $a,c,b,c,d$. The vertex $c$ appears twice. Remove the segment that goes from the first $c$ through $b$ back to the second $c$, and retain $a,c,d$. The edge from $c$ to $d$ was already present after the second visit, so the shortened walk is legitimate.

In a general walk, equal visits at positions $i<j$ permit the same deletion. The edge that followed position $j$ can now follow position $i$ because those two vertices are identical. If position $j$ was the endpoint, simply stop at position $i$. In either case the endpoints are unchanged and at least one edge is removed. A shortest walk cannot admit this operation. This explains both why a repeated vertex is the obstruction and why removing it preserves the required kind of object.

The distance between two connected vertices is the length of a shortest path. The shortest choice is useful not because short objects are aesthetically preferable, but because a removable detour would contradict that choice. This is a recurring way in which an extremal object supplies additional structure for free.

\section{Directed edges and a guaranteed sink}
An oriented graph uses directed edges, written $u\to v$. There are no loops and no pair of opposite edges. Some vertex pairs may have no edge in either direction. The out-neighbors of $v$ are the vertices reached by one outgoing edge. A sink has no outgoing edges.

A directed cycle is a directed closed walk with no repeated vertices except its equal starting and ending vertex. An oriented graph is acyclic when it has no directed cycle.

Here \emph{nonempty} means there is at least one vertex, and \emph{finite} means there are only finitely many vertices. \emph{Acyclic} does not mean there are no edges; a chain $a\to b\to c$ has edges but no route returning to its start. A directed cycle such as $a\to b\to c\to a$ does return and is forbidden by the hypothesis.

The conclusion asks for one vertex with no outgoing edge. It does not say every vertex has this property. In the chain, $c$ is a sink, whereas $a$ and $b$ are not. In a graph containing a single vertex and no edges, that vertex is a sink. This smallest case also explains why we exclude the empty graph: without vertices, there is no candidate to supply.

\par\noindent\begin{minipage}{\linewidth}
\begin{theorem}\label{thm:sink}
Every nonempty finite acyclic oriented graph has a sink.
\end{theorem}
\end{minipage}\par

\noindent\textit{Plan.} Choose a directed path that cannot be made longer. Its last vertex cannot point to any vertex without either extending the path or creating a cycle.
\begin{proof}
There are finitely many directed paths, and at least one consisting of a single vertex, so choose one with greatest length. Write it $v_0\to v_1\to\cdots\to v_k$. Suppose its final vertex $v_k$ had an outgoing edge to a vertex $w$. If $w$ were not on the path, appending it would produce a longer directed path. If $w=v_j$ were on the path, then $j<k$ because loops are forbidden, and $v_j\to\cdots\to v_k\to v_j$ would be a directed cycle. Both possibilities contradict our assumptions. Hence $v_k$ has no outgoing edge and is a sink.
\end{proof}

\section{A small theorem with an extremal hypothesis}
A tournament is an oriented graph in which every pair of distinct vertices has exactly one directed edge between them. Think of each pair having one winner. The outdegree of a vertex counts the vertices it beats.

\par\noindent\begin{minipage}{\linewidth}
\begin{proposition}\label{thm:king}
In a nonempty finite tournament, a vertex of maximum outdegree can reach every other vertex by a directed path of length at most two.
\end{proposition}
\end{minipage}\par

\begin{proof}
Choose a vertex $v$ of maximum outdegree, and let $A$ be its out-neighbors. Consider any other vertex $u$. If $v\to u$, it is reached in one step. If some $a\in A$ satisfies $a\to u$, it is reached by $v\to a\to u$.

Suppose neither happens. The tournament condition then gives $u\to v$, and also $u\to a$ for every $a\in A$, because the alternative $a\to u$ was excluded. Thus $u$ has at least $|A|+1$ out-neighbors, whereas $v$ has exactly $|A|$. This contradicts maximum outdegree. Therefore one of the first two possibilities must occur.
\end{proof}

This theorem is not the second-neighborhood inequality considered below.

In the tournament proof, maximum outdegree is a comparison with \emph{every} vertex. The assumed failure for a particular $u$ forces it to beat all $|A|$ members of $A$, plus $v$, which is not in $A$ because loops are forbidden. That explains the extra one in $|A|+1$. We need not know how $u$ behaves against any remaining vertices: these already counted neighbors suffice for a contradiction. Reaching all vertices within two steps does not by itself compare the number reached in the first step with the number reached for the first time in the second step. Related vocabulary does not make two claims equivalent.

\section{A precise entry into a research statement}
Let $N_1^+(v)$ be the out-neighbors of $v$. Define
\[
N_2^+(v)=\{w\ne v:\ w\notin N_1^+(v),\ \text{and some }u\text{ satisfies }v\to u\to w\}.
\]
These are vertices at directed distance exactly two, not all endpoints of two-edge walks. They are counted as vertices, not as paths. In the triangle $a\to b\to c$ with $a\to c$, the vertex $c$ is not in $N_2^+(a)$ because it is already a first neighbor.

The statement documented for Family~173 of the OpenAI mathematics repository is that every nonempty finite oriented graph has some $v$ with $|N_2^+(v)|\ge|N_1^+(v)|$, with exactly this exclusion convention and with sinks allowed~\cite{oai173}. Our sink theorem proves this statement for the acyclic subclass: choose a sink, for which both sets are empty. This is a complete special-case proof, not a proof for all oriented graphs. Chapter~17 returns to the distinction between such a bridge result and the repository's full claim.

\begin{lab}{When ``two steps away'' is ambiguous}
Construct a directed triangle with edges $a\to b$, $b\to c$, and $a\to c$. Ask an assistant to compute both neighborhoods of $a$. If it counts $c$ as a second neighbor, ask whether ``second'' means a two-edge walk exists or the shortest directed distance is exactly two. Then test a directed three-cycle, where each vertex has one first and one second neighbor. Your output should state the definition before giving the counts.
\end{lab}

\section*{Exercises}
\begin{exercise}\label{ex:7.1}Find the degree sum and number of odd-degree vertices in the four-vertex graph displayed above.\end{exercise}
\begin{exercise}\label{ex:7.2}Prove that among five integers, at least two have the same remainder modulo four.\end{exercise}
\begin{exercise}\label{ex:7.3}Explain why the sink theorem can fail without acyclicity. Give the smallest oriented-graph counterexample.\end{exercise}
\begin{exercise}\label{ex:7.4}For edges $a\to b$, $a\to c$, $b\to d$, $c\to d$, list both neighborhoods of every vertex. Do not count $d$ twice.\end{exercise}
\begin{exercise}\label{ex:7.5}In a graph with $n>0$ vertices and $m$ edges, prove that some vertex has degree at most $2m/n$. Explain why this gives no guarantee that all degrees are near the average.\end{exercise}
\begin{exercise}\label{ex:7.6}A tournament vertex reaches every other vertex in at most two steps. Does that fact alone imply its second out-neighborhood is at least its first? Give a three-vertex example that answers the question.\end{exercise}


% ===== Week 8 =====

\chapter{Assignments, repairs, and obstructions}\label{ch:matching}
\goal{Turn an assignment problem into a precise graph model. Understand why a locally reasonable choice can fail. Derive an algorithm whose failure provides a mathematical certificate, not merely a report that a search stopped.}

\section{Why one acceptable task per person is not enough}
Two students, Ada and Ben, each accept task $1$. Neither accepts any other task. Individually both have an acceptable choice. Together they cannot receive distinct tasks. The problem is collective competition for the same resource.

Now suppose Ada accepts tasks $1,2$, while Ben accepts only $1$. A solution exists: Ada receives $2$ and Ben receives $1$. But a greedy procedure that gives Ada task $1$ immediately gets stuck. The failure of that procedure is not a proof of impossibility. It only says the current assignment cannot be extended without revising an earlier choice.

A bipartite graph models the problem. Its vertices are split into disjoint sets $L$ and $R$, and every edge joins one vertex in $L$ to one in $R$. Let $L$ be students, $R$ tasks, and an edge indicate acceptability. The word bipartite means that edges never connect two students or two tasks in this model.

A matching is a collection of edges with no shared endpoints. It assigns distinct tasks to the students that it covers. A matching \emph{saturates} $L$ when every student is covered. It need not use every task. This distinction matters when there are more tasks than students.

\section{Maximal does not mean maximum}
A matching is maximal if no single new edge can be added without creating a shared endpoint. It is maximum if no matching in the graph has more edges. The greedy choice Ada--$1$ above is maximal: both remaining edges touch an already used vertex. But it is not maximum, because the matching Ada--$2$, Ben--$1$ has two edges.

Inclusion-maximal uses the partial order from Chapter~6: no strictly larger matching contains the current one. In matching terminology, however, \emph{maximum} means largest \emph{number of edges}, not a matching that contains every other matching. Two maximum matchings can have different edges and neither need contain the other. Searching for maximum size may therefore require replacing earlier choices. A proof using an extremal object must specify which comparison is intended.

\subsection*{Read a matching as both a list and a function}
For the current matching Ada--$1$, the list of selected edges has one member. It determines a partial assignment: Ada has a task, Ben does not. Adding Ben--$1$ would give task $1$ two owners. Adding Ada--$2$ would give Ada two tasks. Either addition violates the definition of matching.

A full assignment would be an injective function from students to tasks, with the additional requirement that each assigned pair is an edge. This connects the graph problem to Chapter~3. Injectivity prevents two students from receiving the same task. Acceptability is a separate restriction; an injective function using unacceptable tasks would not solve this problem.

The graph's edge set lists \emph{allowed} assignments. The matching lists \emph{currently chosen} assignments. Confusing these two sets would make the repair operation difficult to follow: changing the matching does not change which tasks are acceptable.

\section{An alternating repair}
Relative to a current matching $M$, an alternating path alternates between edges outside $M$ and edges inside $M$. An augmenting path begins and ends at unmatched vertices and starts and ends with edges outside $M$. In a bipartite assignment problem, it begins at an unmatched student and ends at an unused task.

For Ada and Ben, start with $M=\{\text{Ada--}1\}$. The path
\[
\text{Ben}\;--\;1\;--\;\text{Ada}\;--\;2
\]
uses an unmatched edge, a matched edge, then an unmatched edge. Replace the matching edges along the path by the previously unmatched ones. Ben receives $1$, and Ada moves to $2$.

Write the change as a small ledger. Ben initially has no task and finally has task $1$. Ada initially has task $1$ and finally has task $2$. Task $1$ changes owner rather than being assigned twice. Task $2$ changes from unused to used. Thus one more student is matched without leaving a previously matched student unmatched.

For a longer path, the same exchange repeats along a chain. Internal students each move to the next task; internal tasks each change owner. The unmatched student at the start gains a task, and the unused task at the end makes the chain possible. The following lemma checks this description at every kind of vertex. It is called a lemma because the algorithm will use its conclusion as a supporting fact.

\par\noindent\begin{minipage}{\linewidth}
\begin{lemma}\label{thm:augment}
Flipping the membership of every edge along an augmenting path produces a matching with one more edge.
\end{lemma}
\end{minipage}\par

\begin{proof}
An augmenting path has distinct vertices. Each internal vertex has two path edges, exactly one originally in the matching. After the flip, exactly the other path edge belongs to the new matching, so that vertex still has one assigned edge. Each endpoint was unmatched and gains its single incident path edge. Vertices outside the path are unchanged. No internal vertex can have another matching edge outside the path, because it already had a matching edge on the path. Hence no vertex receives two matching edges.

The alternating path starts and ends with edges outside $M$, so it contains one more outside edge than inside edge. The flip adds the outside edges and removes the inside edges. The total matching size therefore increases by one.
\end{proof}

The argument has both an invariant and progress. The invariant is ``the selected edges form a matching.'' The progress measure is the number of selected edges. Since it can never exceed $|L|$, there can be at most $|L|$ augmentations.

\section{Searching while preserving a trace}
Choose an unmatched student $x$. Explore from a student to a task along an edge not currently matched; from a matched task, continue to the unique student occupying it. If an unused task is reached, the recorded predecessor links give an augmenting path. Repeated vertices need not be explored again. Since there are finitely many vertices, this exploration terminates.

One may use breadth-first search, which explores by increasing path length, but the proof only needs a complete search of the reachable vertices under these alternating rules. Keeping predecessor links is essential for reconstruction: a statement that ``a task was reached'' is less useful than a path showing how assignments can be repaired.

A beginner's pseudocode is:
\par\noindent\begin{minipage}{\linewidth}
\begin{lstlisting}
Start with the empty matching.
Choose a student who is still unmatched.
Explore alternating paths from that student.
If an unused task is reached, flip along the path.
If no unused task is reachable, return an obstruction.
Repeat until every student is matched.
\end{lstlisting}
\end{minipage}\par

The last line of the failure branch is deliberately stronger than ``return failure.'' We now derive the promised obstruction.

\subsection*{What ``the search is complete'' means}
Reachability is defined by the allowed search moves, not by a guess about a promising path. From the root student, inspect every unmatched acceptable edge. For each reached occupied task, record its owner and inspect that owner's unmatched acceptable edges. Continue until an unused task is found or there is nothing newly reachable to inspect. Do not ignore an available branch merely because another branch looks more promising.

Recording a predecessor means remembering which earlier vertex first led to each new vertex. Following these records backward reconstructs the route. A vertex already recorded need not be recorded again; this prevents endless exploration around a cycle. There are finitely many vertices to record. If no unused task is found after all recorded vertices have been inspected, no allowed search move has been left unexplored. The proof of impossibility below needs this last fact, not merely the fact that one attempted route stopped.

\section{A failed search exposes a shortage}
For $S\subseteq L$, let $N(S)$ be the tasks acceptable to at least one student in $S$. Suppose a complete alternating search from an unmatched student $x$ reaches no unused task. Let $A$ be the reached students, including $x$, and let $B$ be the reached tasks.

Before the general counting argument, consider three students $a,b,c$ who can all use only tasks $1,2$. Suppose $a$ is matched to $1$ and $b$ to $2$, with root $x=c$. The search reaches both tasks and then their owners $a,b$. It finds no unused task. Thus $A=\{c,a,b\}$ and $B=\{1,2\}$.

The existing matching pairs task $1$ with $a$ and task $2$ with $b$. Every reached student \emph{except the unmatched root} appears in exactly one such pair. There is no pair for $c$, which is precisely the shortage. We now show that this one-extra-student pattern holds whenever a complete search fails, not just in this example.

Every task in $B$ is matched; otherwise the search would have succeeded. Its matched student is reached by the search rule and lies in $A\setminus\{x\}$. Conversely, every reached student other than $x$ was reached from a matched task. Because matching partners are unique, the matching pairs give a bijection between $B$ and $A\setminus\{x\}$. Therefore
\[
|A|=|B|+1.
\]

A bijection here is not an abstract new object that we must discover. It is the existing partner rule: send each task in $B$ to its owner. Two tasks cannot have the same owner because $M$ is a matching. Every student in $A\setminus\{x\}$ is an owner of a reached task because that is the only rule that can reach a new student. Thus the partner rule is injective and surjective onto that set. Equal finite sizes follow, and adding back the single root gives the displayed equation.

The separate identity $N(A)=B$ is needed because the shortage condition counts \emph{all acceptable tasks} of $A$, not merely whichever tasks the search happened to record. A proof that reached tasks are too few would be inconclusive if some acceptable tasks had been missed. The next two-containment argument rules that out.

We also have $N(A)=B$. Every reached task was entered from a reached student, so $B\subseteq N(A)$. For the reverse containment, take a task neighboring a reached student. If their connecting edge is unmatched, the exploration follows it, so the task lies in $B$. If their edge is matched, the student is not the unmatched root $x$, and its matched task is the task from which it was reached; that task already lies in $B$. Thus all neighbors lie in $B$.

Combining the two identities gives
\[
|N(A)|=|B|=|A|-1<|A|.
\]
There are fewer available tasks than students in $A$. No assignment can give all those students distinct acceptable tasks. The search has produced an independently checkable certificate of impossibility.

\begin{principle}{Habit 7: make failure informative}
A failed attempt is not automatically a counterexample. Ask whether the failure can be turned into a concrete object that violates a necessary condition. In this problem, the object is a set of students competing for too few tasks. In another problem it may be two distinct inputs with the same output, refuting a claim of injectivity, or a particular value violating a required inequality.
\end{principle}

\section{A complete trace}
Let the acceptable tasks be
\[
N(a)=\{1,2\},\qquad N(b)=\{1\},\qquad N(c)=\{2,3\}.
\]
A sequence of choices could first match $a$ to $1$ and $c$ to $2$. Student $b$ is unmatched. From $b$ the search reaches task $1$, then its matched student $a$, then task $2$, then its matched student $c$, then the unused task $3$. The augmenting path is
\[
b--1--a--2--c--3.
\]

\begin{center}
\begin{tikzpicture}[x=0.88in,y=0.5in,every node/.style={font=\small}]
\foreach \x/\name in {0/b,1/1,2/a,3/2,4/c,5/3}
 {\node[circle,draw=ink,minimum size=7mm,inner sep=1pt] (v\x) at (\x,0) {$\name$};}
\draw[accent,dashed,thick] (v0)--(v1) (v2)--(v3) (v4)--(v5);
\draw[ink,line width=1.2pt] (v1)--(v2) (v3)--(v4);
\node[align=center] at (2.5,-.8) {Solid: remove the old matched edges.\\Dashed: add the new matched edges.};
\end{tikzpicture}
\end{center}

After the flip, the matching is $b$--$1$, $a$--$2$, $c$--$3$. An early choice was not irrevocable; the alternating structure described exactly how to revise it without creating a new conflict.

Now remove task $3$ from $c$'s options. The same search reaches all three students and only tasks $1,2$. It cannot augment, and $A=\{a,b,c\}$ satisfies $|N(A)|=2<3$. This is not merely a bad search order: it certifies that no full assignment exists.

\begin{lab}{Interrogating an algorithmic failure}
Ask an assistant to solve the three-student instance above greedily, then to explain any failure. Require it to distinguish ``my current matching is stuck'' from ``no full matching exists.'' For a feasible instance, demand an alternating repair path. For the modified infeasible instance, demand a specific set $A$ and the computed set $N(A)$. Verify both certificates by hand before running the companion code.
\end{lab}

\section*{Exercises}
\begin{exercise}\label{ex:8.1}Give a maximal matching that is not maximum in the Ada--Ben example, and prove both assertions.\end{exercise}
\begin{exercise}\label{ex:8.2}For $N(a)=\{1,2\}$, $N(b)=\{1,3\}$, $N(c)=\{2\}$, start from $a$--$2$, $b$--$1$ and find an augmenting path.\end{exercise}
\begin{exercise}\label{ex:8.3}For three students all accepting only tasks $1,2$, give a shortage certificate and explain why each student's having two options is insufficient.\end{exercise}
\begin{exercise}\label{ex:8.4}Why must an augmenting path have no repeated vertex in the flip argument? Explain how a shortest alternating route supplies such a path.\end{exercise}
\begin{exercise}\label{ex:8.5}Explain why a proof that matching size increases by one proves termination only when the left set is finite. State the upper bound on the number of augmentations.\end{exercise}
\begin{exercise}\label{ex:8.6}Suppose a failed alternating search gives reached sets $A,B$. Prove $N(A)=B$ carefully, treating matched and unmatched edges separately.\end{exercise}


% ===== Week 9 =====

\chapter{Hall's theorem: a complete proof apprenticeship}\label{ch:hall}
\goal{Read and reconstruct a multi-stage theorem proof. Learn why equality in a necessary bound can determine the right decomposition. Compare a constructive proof with an inductive proof and identify their shared obstruction.}

\section{A condition that sees every competing group}
The assignment examples suggest the condition
\[
|N(S)|\ge |S|\qquad\text{for every subset }S\subseteq L.
\tag{H}\label{eq:hall}
\]
It says that every group of students collectively has at least as many available tasks as students. Checking only one student at a time misses shared bottlenecks. Checking only the entire class misses a smaller group trapped among too few tasks.

For example, let $N(a)=N(b)=\{1\}$ and $N(c)=\{2,3\}$. The whole class has three tasks available and every student has an option. Yet the group $\{a,b\}$ has only one task. The failure is local to a subset, even though the desired assignment is global.

Read condition (H) as a repeated test: choose any group $S$ of students, gather its distinct acceptable tasks into $N(S)$, and compare the two counts. A task accepted by three members of $S$ is still one task in $N(S)$, not three units of capacity.

The empty group has no acceptable tasks to gather, so its test is $0\ge0$. A group consisting of one student tests whether that student has an option. A group consisting of all students tests overall capacity. The intermediate groups test competition that neither of those extreme tests sees. The word ``every'' is what includes all of these sizes.

The theorem below is an \emph{if and only if} statement. One direction says an assignment implies these inequalities. The other says the inequalities imply an assignment. There is no circularity in using the inequalities in the second direction: there they are the hypothesis, not something inferred from a matching we have already assumed exists.

\par\noindent\begin{minipage}{\linewidth}
\begin{theorem}[Hall's matching theorem]\label{thm:hall}
A finite bipartite graph with parts $L,R$ has a matching saturating $L$ if and only if condition \eqref{eq:hall} holds.
\end{theorem}
\end{minipage}\par

This classical theorem is a standard bridge from elementary counting to structural combinatorics; the exposition here is self-contained. Further developments appear in the graph and matching material of the MIT text~\cite{mcs}.

\section{First prove the necessary direction}
Assume a matching saturates $L$. Take any $S\subseteq L$. The matching assigns one acceptable task to each student in $S$, and these tasks are distinct. All are in $N(S)$. Therefore $N(S)$ contains at least $|S|$ elements, proving \eqref{eq:hall}.

Notice how little was needed. We did not inspect the rest of the graph. A candidate matching immediately provides the required distinct tasks. The harder direction must reverse this reasoning: a family of inequalities must force a coherent assignment.

\section{The constructive proof is already within reach}
Start from the empty matching and repeatedly apply the search procedure from Chapter~8. If a student remains unmatched, search from that student. Under condition \eqref{eq:hall}, the search cannot fail: failure would produce a subset $A$ with $|N(A)|<|A|$, contradicting \eqref{eq:hall}. Thus it finds an augmenting path and increases the matching size by one.

At most $|L|$ such increases are possible. The procedure therefore reaches a matching saturating $L$. This proves sufficiency. The case $L=\varnothing$ is covered by the empty matching, and the condition is then true because the only subset is empty.

This argument proves both existence and a method of construction. It also tells us how to distinguish a true impossibility from an unlucky sequence of choices. The algorithm returns either an assignment or a shortage set, each checkable without trusting the search process that found it.

\section{A second proof: equality tells us where to split}
An inductive proof reveals a different useful idea. A subset $S$ is called tight when $|N(S)|=|S|$. Such a subset has just enough tasks, with no spare capacity. It is a natural block to solve first. If no nonempty proper subset is tight, every such subset has at least one spare task, and that spare capacity lets us make an arbitrary first assignment.

\noindent\textit{Proof plan.} Induct on $|L|$. If a proper tight block exists, solve inside the block and outside it. Otherwise assign one student any acceptable task and show the remaining graph still satisfies Hall's condition.

\subsection*{A concrete tight block before the general proof}
Let four students have options
\[
 N(a)=N(b)=\{1,2\},\qquad N(c)=\{2,3\},\qquad N(d)=\{3,4\}.
\]
The group $S=\{a,b\}$ has exactly two available tasks, so it is tight. Assigning all of $S$ must use both tasks $1,2$. The remaining students therefore have options $c:\{3\}$ and $d:\{3,4\}$ after those two tasks are removed. They can use $3$ and $4$ respectively.

Notice what happened to $T=\{c\}$. Originally $N(T)=\{2,3\}$, which overlaps $N(S)=\{1,2\}$. In the remaining problem its neighborhood is $N'(T)=\{3\}$. The prime in $N'$ labels a \emph{new neighborhood function}; it is not differentiation. The combined group $S\cup T$ had three distinct tasks. Subtracting the two tasks committed to $S$ left one for $T$.

The general proof must establish this remaining capacity for every group $T$ outside $S$, not just for a single student. That is why it applies the original condition to $S\cup T$. Applying it to $T$ alone would give no control over how many of $T$'s options disappear when $N(S)$ is removed.

The notation $S\subsetneq L$ means $S\subseteq L$ but $S\ne L$; such a subset is called \emph{proper}. Nonempty and proper ensure $0<|S|<|L|$. Both the block $S$ and the remaining set $L\setminus S$ are then smaller than the original left set, making the induction hypothesis available.

\begin{proof}[Inductive proof of sufficiency]
We use strong induction on $n=|L|$. For $n=0$, the empty matching works. Assume $n>0$ and the result holds for all smaller left sets. Suppose our graph satisfies \eqref{eq:hall}.

\textbf{Case 1: a nonempty proper tight subset exists.} Choose $S$ with $\varnothing\ne S\subsetneq L$ and $|N(S)|=|S|$. 

\emph{First smaller problem: match the tight block.} Restrict the graph to left part $S$ and right part $N(S)$. For every $U\subseteq S$, all neighbors of $U$ already lie in $N(S)$, so the restricted neighborhood is exactly $N(U)$. The original condition gives $|N(U)|\ge|U|$. Thus the restricted graph satisfies Hall's condition. Since $|S|<n$, induction gives a matching of all students in $S$ into $N(S)$.

\emph{Second smaller problem: keep enough tasks for the other students.} Remove all tasks in $N(S)$ and consider parts $L\setminus S$ and $R\setminus N(S)$. Take any $T\subseteq L\setminus S$. Its available tasks in this remaining graph form
\[
N'(T)=N(T)\setminus N(S).
\]
Because $S$ and $T$ are disjoint, $|S\cup T|=|S|+|T|$. Applying Hall's condition to $S\cup T$ gives
\[
|N(S)\cup N(T)|\ge |S|+|T|.
\]
The union on the left is the disjoint union of $N(S)$ and $N(T)\setminus N(S)$. Consequently,
\[
|N(S)|+|N'(T)|\ge |S|+|T|.
\]
Substitute $|N(S)|=|S|$ and subtract $|S|$ from both sides to obtain $|N'(T)|\ge|T|$. This holds for every $T$, including the empty set. The remaining graph satisfies Hall's condition and has fewer than $n$ students. Induction therefore gives a matching of all remaining students using only tasks outside $N(S)$. \emph{Recombine the two answers.} The first matching uses students in $S$ and tasks in $N(S)$. The second uses students outside $S$ and tasks outside $N(S)$. No student or task can therefore occur in both matchings. Their union is a matching, and every original student belongs to one of the two matched groups. Hence it saturates $L$.

\textbf{Case 2: no nonempty proper subset is tight.} Hall's condition and the integrality of cardinalities imply
\[
|N(S)|\ge |S|+1
\]
for every nonempty proper subset $S\subsetneq L$. Choose a student $x\in L$. Hall's condition for $\{x\}$ ensures at least one acceptable task $y\in N(\{x\})$. Assign $x$ to $y$, then remove both vertices.

For any nonempty subset $T\subseteq L\setminus\{x\}$, the set $T$ is a proper subset of the original $L$. Before removal it has at least $|T|+1$ neighbors. Removing the single task $y$ deletes at most one of them. Hence its new neighborhood has size at least $|T|$. The empty subset satisfies the condition automatically. Thus the smaller graph satisfies Hall's condition and induction supplies a matching of its students. Adding the edge $x$--$y$ gives a matching saturating $L$.

The added edge cannot conflict with that smaller matching because its two endpoints were removed before the smaller problem was solved. When $n=1$, there are no remaining students. The smaller matching is empty, and the chosen edge alone solves the problem. We have therefore not relied on a nonexistent nonempty remaining group in this boundary case.

The two cases exhaust the possibilities, completing the induction.
\end{proof}

\section{Why every part of the case split is needed}
The adjective \emph{proper} ensures both induced left sets in Case~1 are smaller. The adjective \emph{nonempty} prevents a vacuous split that leaves the original problem unchanged. Tightness is what lets us subtract $|S|$ without losing capacity in the remaining graph. In Case~2, the extra integer unit is exactly what pays for removing one task.

An assistant that says ``split into a subset and its complement and apply induction'' has not yet proved the theorem. It must show that both smaller problems satisfy the condition being used as an induction hypothesis. That verification is the heart of the proof, not a bookkeeping detail.

\pauseq{In Case~1, why not remove only the tasks actually chosen for $S$ instead of all of $N(S)$?}{The matching for $S$ uses $|S|$ distinct tasks, and tightness says $N(S)$ has exactly $|S|$ tasks. It therefore uses all of $N(S)$. Removing that set is both natural and necessary to prevent conflicts.}

\subsection*{Using a theorem requires matching its entire input specification}
The induction hypothesis is not simply ``smaller graphs have matchings.'' That statement is false. It says \emph{smaller finite bipartite graphs satisfying Hall's condition} have the required matchings. Every application must therefore supply three pieces of information: the smaller left and right sets, a smaller left size, and the inequalities for every subset of that smaller left set.

The neighborhood calculations in the proof supply the third piece. Removing vertices preserves bipartiteness and finiteness, but it does not automatically preserve Hall's condition. This is a general reading habit: when a proof says ``apply the theorem,'' expand its hypotheses and check the correspondence. A matching name is not enough; the mathematical input must fit.

\section{Mathematical taste from equality cases}
Many arguments begin with an inequality. Equality may identify a rigid block; strict inequality may provide spare capacity. Hall's proof uses this distinction systematically. It is not an isolated trick attached to one theorem.

The useful question is not just ``Can the bound be sharpened?'' It is also ``What happens when the bound is attained?'' If equality forces a recognizable structure, that structure may allow a decomposition. If strictness supplies a quantitative margin, the margin may absorb a perturbation or a removal. We will see this margin reappear as $1-q>0$ in the contraction theorem.

\begin{lab}{A proof map with an explicit missing obligation}
Give an assistant only the inductive plan, not the completed proof. Ask it to state the exact condition required for each smaller graph. Require it to expand the neighborhood calculation in Case~1 as a disjoint union. Then ask it why the second case still works when $|L|=1$. The answer is that the chosen edge exists by Hall's condition, and the remaining left set is empty; there are no nonempty residual subsets to check.

Your final record should contain a one-paragraph plan, the full inequality chain for $N'(T)$, and an explanation of the spare unit in Case~2. Reconstruct one of the two proofs without the assistant. Understanding two plans is useful; being unable to reproduce either is a signal to revisit the dependency you have not yet mastered.
\end{lab}

\section*{Exercises and midterm portfolio}
\begin{exercise}\label{ex:9.1}Give a graph with three students and three tasks for which every student has a neighbor and $|N(L)|=|L|$, but no saturating matching exists. Identify the violating subset.\end{exercise}
\begin{exercise}\label{ex:9.2}For $N(a)=\{1\}$, $N(b)=\{1,2\}$, $N(c)=\{2,3\}$, verify Hall's condition for every subset and find a nonempty proper tight subset.\end{exercise}
\begin{exercise}\label{ex:9.3}In Case~1 of the proof, justify the equality $|N(S)\cup N(T)|=|N(S)|+|N(T)\setminus N(S)|$ without assuming the original neighborhoods are disjoint.\end{exercise}
\begin{exercise}\label{ex:9.4}Explain why strict inequality becomes $|N(S)|\ge|S|+1$. Would the same inference hold for arbitrary real quantities?\end{exercise}
\begin{exercise}\label{ex:9.5}Suppose each student accepts exactly $d>0$ tasks and each task is acceptable to at most $d$ students. Use double counting and Hall's theorem to prove a full assignment exists.\end{exercise}
\begin{exercise}\label{ex:9.6}Complete the midterm portfolio using the graph of Exercise~9.2 and an infeasible comparison graph. Supply an assignment certificate, a shortage certificate, one complete sufficiency proof, and a paragraph distinguishing maximal from maximum matchings. The solution part provides a model portfolio structure.\end{exercise}


% ===== Week 10 =====

\chapter{Linearity: changing the language of a problem}\label{ch:linear}
\goal{Learn vectors and matrices from their actions on lists of numbers. Recognize linearity as preservation of addition and scaling. Use a change of representation to explain uniqueness and orthogonality rather than merely perform calculations.}

\section{Vectors are ordered measurements}
Throughout the coordinate examples, $n$ is a positive integer. A vector in $\RR^n$ is an ordered list of $n$ real numbers. For $n=2$, write $x=(x_1,x_2)$. The coordinates might record two resource amounts or the location of a point. Order matters: $(1,3)$ and $(3,1)$ usually describe different data.

Add vectors coordinate by coordinate and multiply them by a real scalar $c$:
\[
(x_1,x_2)+(y_1,y_2)=(x_1+y_1,x_2+y_2),\qquad c(x_1,x_2)=(cx_1,cx_2).
\]
The zero vector has all coordinates zero. Subtraction is addition of the negative vector. An expression $a_1v_1+\cdots+a_kv_k$ is a linear combination of vectors $v_1,\ldots,v_k$ with scalar coefficients $a_1,\ldots,a_k$.

The vectors $e_1=(1,0)$ and $e_2=(0,1)$ are building blocks for the plane:
\[
(x_1,x_2)=x_1e_1+x_2e_2.
\]
This expression is unique, since equality of ordered pairs forces equality in each coordinate. A list that expresses every vector uniquely by linear combinations is a basis. The point of a basis is not that its vectors are intrinsically special; it provides a coordinate description that can be transformed or chosen to suit a problem.

\subsection*{Why the two words in ``unique representation'' matter}
With $e_1=(1,0)$ and $e_2=(0,1)$, the vector $(3,-2)$ is $3e_1-2e_2$. Existence of such a representation for every vector means the building blocks \emph{span} the plane. Uniqueness means the same vector cannot be made from two different coefficient lists.

The single vector $(1,1)$ does not span the plane: every multiple has equal coordinates, so it cannot make $(1,0)$. The list $(1,0),(0,1),(1,1)$ does span it, but does not give unique coefficients. The vector $(1,1)$ can be made as one copy of the third vector or as one copy each of the first two. A basis requires both spanning and uniqueness.

An equivalent test for uniqueness is that a combination equaling zero must have all coefficients zero; this property is called \emph{linear independence}. To see the connection, subtract two representations of the same vector. The differences of their coefficients make a combination equal to zero. If only zero coefficients can do this, the two original coefficient lists agree. Conversely, a nonzero coefficient list making zero would give a second representation of the zero vector besides the all-zero list.

\section{A linear map is determined by its building blocks}
A map $T:\RR^n\to\RR^m$ is linear when it preserves addition and scalar multiplication. Equivalently, for all vectors $x,y$ and real scalars $a,b$,
\[
T(ax+by)=aT(x)+bT(y).
\]
For a linear map on the plane,
\[
T(x_1,x_2)=x_1T(e_1)+x_2T(e_2).
\]
Knowing the two outputs $T(e_1),T(e_2)$ therefore tells us every output. This is a mathematical explanation of why a finite array can describe a transformation acting on infinitely many inputs.

Suppose $T(e_1)=(1,1)$ and $T(e_2)=(1,-1)$. Then
\[
T(x_1,x_2)=(x_1+x_2,x_1-x_2).
\]
We record the two output vectors as columns of a matrix:
\[
A=\begin{pmatrix}1&1\\1&-1\end{pmatrix},\qquad
Ax=\begin{pmatrix}x_1+x_2\\x_1-x_2\end{pmatrix}.
\]
An $m$-by-$n$ matrix has $m$ rows and $n$ columns. Its action is defined by
\[
(Ax)_i=\sum_{j=1}^n A_{ij}x_j\qquad(1\le i\le m).
\]
The subscript $ij$ identifies the entry in row $i$, column $j$. The formula says to multiply entries in matching positions of a row and the input list, and then add those products. We will call this operation a dot product below.

For example, if a row is $(2,-1,3)$ and the input is $(4,5,0)$, the corresponding output coordinate is
\[
 2\cdot4+(-1)\cdot5+3\cdot0=3.
\]
The input must have three coordinates because that row has three entries. An $m$-by-$n$ matrix takes $n$-coordinate inputs and produces $m$-coordinate outputs. The two letters indicate possibly different sizes; a matrix need not be square.

We sometimes write a vector horizontally as $(x_1,x_2)$ and sometimes vertically as a column. They name the same ordered data here. In a product $Ax$, we use the column convention so the displayed array agrees with the row-by-input formula.

\pauseq{Is $F(x_1,x_2)=(x_1+1,x_2)$ linear?}{No. Every linear map must send zero to zero: $T(0)=T(0+0)=T(0)+T(0)$ implies $T(0)=0$. Here $F(0,0)=(1,0)$. A translation adds a fixed vector to every input. It preserves differences between inputs but generally moves the zero vector. Such a rule is an example of an affine map; we need only the concrete distinction from a linear map here.}

\section{Kernel: the information lost at zero}
The kernel of a linear map is the set of inputs sent to zero:
\[
\ker T=\{x:T(x)=0\}.
\]
If two inputs have the same output, their difference belongs to the kernel because
\[
T(x-y)=T(x)-T(y)=0.
\]
Thus every collision can be translated into an invisible nonzero input.

\par\noindent\begin{minipage}{\linewidth}
\begin{theorem}\label{thm:kernel}
A linear map is injective if and only if its kernel consists only of the zero vector.
\end{theorem}
\end{minipage}\par

\begin{proof}
If $T$ is injective and $T(x)=0$, then $T(x)=T(0)$. Injectivity gives $x=0$, so the kernel contains no other vector.

Conversely, suppose the only vector mapped to zero is zero. If $T(x)=T(y)$, linearity gives $T(x-y)=0$. The kernel assumption forces $x-y=0$, hence $x=y$. Therefore $T$ is injective.
\end{proof}

The kernel theorem is a two-direction statement about a \emph{linear} map. In its first direction, linearity ensures $T(0)=0$, so an input with output zero can be compared with the zero input. In its second direction, linearity is needed to turn a collision into $T(x-y)=0$. For a nonlinear map, inspecting only inputs sent to zero may miss collisions elsewhere. For instance, $F(x)=x^2$ sends only zero to zero, yet $F(1)=F(-1)$.

This also separates solving a particular equation $T(x)=b$ from proving all targets attainable. A zero kernel says there cannot be two different solutions for any one target. It does not assert that a solution exists for every target. The inclusion $t\mapsto(t,0)$ from $\RR$ to $\RR^2$ is linear and injective, but no input maps to $(0,1)$. To see why counting equations alone cannot establish uniqueness, consider $T(x_1,x_2)=(x_1+x_2,2x_1+2x_2)$. It has two outputs but loses the nonzero vector $(1,-1)$. The second equation contains no new distinction beyond the first.

For the earlier matrix $A$, apply its sum-and-difference rule twice:
\begin{align*}
 A(Ax)&=A(x_1+x_2,x_1-x_2)\\
 &=\big((x_1+x_2)+(x_1-x_2),\ (x_1+x_2)-(x_1-x_2)\big)\\
 &=(2x_1,2x_2)=2x.
\end{align*}
Therefore $Ax=0$ implies $2x=0$, hence $x=0$. The map is injective. Moreover, for any target $b$, taking $x=Ab/2$ gives $Ax=A(Ab)/2=b$. This proves surjectivity and explicitly constructs the inverse. We have proved reversibility without introducing determinants.

\section{The geometry of an inner product}
For vectors $u,v\in\RR^n$, define their dot product by
\[
u\cdot v=\sum_{j=1}^n u_jv_j.
\]
This dot product is the real coordinate \emph{inner product} used in this chapter. The name describes the same operation, not a second calculation.

Symmetry follows because each product $u_jv_j$ equals $v_ju_j$. The phrase ``linear in each argument'' means that one argument can be held fixed while addition and scaling in the other distribute through the dot product. For example,
\[
 (a u+b w)\cdot v
 =\sum_j(au_j+bw_j)v_j
 =a\sum_j u_jv_j+b\sum_j w_jv_j
 =a(u\cdot v)+b(w\cdot v).
\]
The same expansion works in the second argument. All sums here contain exactly $n$ terms. The Euclidean length is
\[
\norm{u}_2=\sqrt{u\cdot u}=\sqrt{\sum_{j=1}^n u_j^2}.
\]
Each square is nonnegative, so the square root is defined. Length is zero exactly when every coordinate is zero.

The subscript $2$ on $\norm{u}_2$ identifies this square-based length. It is not the second coordinate of $u$ and does not restrict the vector to two coordinates. Vectors are orthogonal when their dot product is zero. In the plane this agrees with perpendicularity, but the algebraic definition works in any finite dimension.

For $u=(1,1)$ and $v=(1,-1)$, $u\cdot v=1-1=0$. Orthogonality says the positive and negative contributions cancel. This same cancellation will express zero correlation between sign sequences in Chapter~14.

\section{Deriving a bound by removing a component}
We want to compare the size of a dot product with the lengths of its two inputs. For $u=(3,1)$ and $v=(1,1)$, the dot product is four and their squared lengths are ten and two. The proposed bound becomes $4\le\sqrt{10}\sqrt2$, or, after squaring nonnegative quantities, $16\le20$.

How could we find a general reason? Subtract a multiple $t v$ from $u$. Its squared length cannot be negative, no matter which real number $t$ we choose. The useful choice will remove the component that has a nonzero dot product with $v$. Solve the equation
\[
 (u-tv)\cdot v=0.
\]
Distribution gives $u\cdot v-t(v\cdot v)=0$. When $v\ne0$, its squared length $v\cdot v$ is positive, so solving for $t$ gives
\[
 t=\frac{u\cdot v}{v\cdot v}.
\]
This explains the coefficient that appears in the proof: it is the solution of a simple equation, not an unexplained guess. In our example $t=4/2=2$, so $u-tv=(1,-1)$, whose dot product with $(1,1)$ is indeed zero.

The residual vector $u-tv$ is what remains after subtracting the chosen multiple. Its squared length is two in the example. The proof below turns the nonnegativity of this residual into the desired bound. Absolute value around $u\cdot v$ is needed because a dot product can be negative, whereas the product of two lengths is nonnegative.

\par\noindent\begin{minipage}{\linewidth}
\begin{theorem}[Cauchy--Schwarz inequality]\label{thm:cs}
For all $u,v\in\RR^n$, $|u\cdot v|\le\norm{u}_2\norm{v}_2$.
\end{theorem}
\end{minipage}\par

\noindent\textit{Plan.} The squared length of $u-tv$ is nonnegative for every real $t$. Choose $t$ to remove as much of the component along $v$ as possible.
\begin{proof}
If $v=0$, both sides involving the dot product are zero and the inequality holds. Assume $v\ne0$, so $v\cdot v>0$. For any real $t$, expand one argument at a time:
\begin{align*}
 (u-tv)\cdot(u-tv)
 &=u\cdot(u-tv)-t\,v\cdot(u-tv)\\
 &=u\cdot u-t(u\cdot v)-t(v\cdot u)+t^2(v\cdot v)\\
 &=u\cdot u-2t(u\cdot v)+t^2(v\cdot v).
\end{align*}
The last step uses symmetry of the dot product. The left side is a squared length, so the resulting expression is nonnegative.
Choose $t=(u\cdot v)/(v\cdot v)$, which is defined because the denominator is positive. Substitution yields
\begin{align*}
0&\le u\cdot u-2\frac{(u\cdot v)^2}{v\cdot v}
 +\frac{(u\cdot v)^2}{v\cdot v}\\
&=u\cdot u-\frac{(u\cdot v)^2}{v\cdot v}.
\end{align*}
Move the fraction to the left and multiply by the positive quantity $v\cdot v$:
\[
(u\cdot v)^2\le(u\cdot u)(v\cdot v).
\]
Both sides are nonnegative. Taking nonnegative square roots gives the stated inequality, because $\sqrt{(u\cdot v)^2}=|u\cdot v|$.
\end{proof}

For $v\ne0$, equality holds exactly when the chosen vector $u-tv$ has zero length, that is, when $u$ is a scalar multiple of $v$.

Indeed, after choosing $t$, the difference between the squared bound's right and left sides is
\[
 (u\cdot u)(v\cdot v)-(u\cdot v)^2
 =(v\cdot v)\,\norm{u-tv}_2^2.
\]
Since $v\cdot v>0$, the difference vanishes exactly when the residual vector is zero. This is the equality case, not an additional assumption imposed to prove the inequality. If $v=0$, equality holds for every $u$, as the proof's separate first case already showed. This equality condition explains when the bound loses no information. It is often worth studying equality before trying to improve a bound.

\subsection*{From the dot-product bound to a distance bound}
The next chapter needs the fact that traveling through an intermediate vector cannot be shorter in total than the direct route. We establish its algebraic form here rather than leave a prerequisite hidden in an exercise. Expand
\begin{align*}
 \norm{u+v}_2^2
 &=(u+v)\cdot(u+v)\\
 &=\norm{u}_2^2+2(u\cdot v)+\norm{v}_2^2\\
 &\le\norm{u}_2^2+2|u\cdot v|+\norm{v}_2^2\\
 &\le\norm{u}_2^2+2\norm{u}_2\norm{v}_2+\norm{v}_2^2\\
 &=(\norm{u}_2+\norm{v}_2)^2.
\end{align*}
The first inequality uses $a\le|a|$ for a real number; the second uses Cauchy--Schwarz. Both lengths in the final comparison are nonnegative, so taking square roots preserves the inequality. We obtain the \emph{triangle inequality}
\[
 \norm{u+v}_2\le\norm{u}_2+\norm{v}_2.
\]
Exercise~10.5 asks you to reconstruct this argument, not to supply an unexplained prerequisite before you can continue.

\section{Matrix multiplication records composition}
If $A$ maps $\RR^n$ to $\RR^m$ and $B$ maps $\RR^m$ to $\RR^k$, their product $BA$ records ``$B$ after $A$.'' Expanding one coordinate gives
\[
(B(Ax))_i=\sum_{r=1}^m B_{ir}\sum_{j=1}^n A_{rj}x_j
=\sum_{j=1}^n\left(\sum_{r=1}^m B_{ir}A_{rj}\right)x_j.
\]
The sums are finite, so regrouping is legitimate. This derives $(BA)_{ij}=\sum_rB_{ir}A_{rj}$. Matrix multiplication is not an arbitrary convention; it is forced by composition of the actions.

To inspect the regrouping without nested sums, take both $A$ and $B$ to be two-by-two matrices. The first coordinate after applying $B$ to $Ax$ has the form
\begin{align*}
 B_{11}(A_{11}x_1+A_{12}x_2)
 &+B_{12}(A_{21}x_1+A_{22}x_2)\\
 &=(B_{11}A_{11}+B_{12}A_{21})x_1\\
 &\quad +(B_{11}A_{12}+B_{12}A_{22})x_2.
\end{align*}
The coefficients multiplying $x_1$ and $x_2$ are the first row of the composed matrix. In the general formula, the index $r$ runs through the intermediate coordinates, while $j$ runs through the original inputs. These are different jobs. The summation notation is a compact record of this same distributive calculation, not a new operation.

The transpose $A^{\mathsf T}$ exchanges rows and columns. Consequently, $(AA^{\mathsf T})_{ij}$ is the dot product of row $i$ with row $j$. A matrix identity $AA^{\mathsf T}=nI$ says its rows have squared length $n$ and different rows are orthogonal. Here $I$ is the identity matrix, with ones on the diagonal and zeros elsewhere. This is the exact identity we will later read in a formalized research statement.

\begin{lab}{Replacing equation counting by a structural test}
\noindent\textbf{Deliberately flawed reply.} ``There are two linear equations and two unknowns, so the solution is unique.''

\noindent\textbf{Learner.} ``Test the matrix with rows $(1,1)$ and $(2,2)$. Find a nonzero vector it sends to zero. Then state a correct criterion for uniqueness.''

\noindent\textbf{Model answer.} The vector $(1,-1)$ is in the kernel. Any existing solution can be changed by a multiple of this vector without changing the outputs. Uniqueness of solutions, when they exist, follows from a zero kernel; existence for every target is a separate surjectivity claim. Ask the assistant to give one target that has solutions and one that does not for this matrix.
\end{lab}

\section*{Exercises}
\begin{exercise}\label{ex:10.1}For $A=\begin{pmatrix}2&1\\0&3\end{pmatrix}$ and $x=(1,-2)$, compute $Ax$. Explain the answer using the columns of $A$.\end{exercise}
\begin{exercise}\label{ex:10.2}Find the kernel of $T(x,y)=x+y$. Describe all solutions of $T(x,y)=5$.\end{exercise}
\begin{exercise}\label{ex:10.3}Show the map $T(x,y)=(x+y,x-y)$ has inverse $(s,t)\mapsto((s+t)/2,(s-t)/2)$. Check both compositions.\end{exercise}
\begin{exercise}\label{ex:10.4}Compute the dot product and lengths of $u=(1,2)$ and $v=(2,-1)$. Check Cauchy--Schwarz for this pair.\end{exercise}
\begin{exercise}\label{ex:10.5}Using Cauchy--Schwarz, prove the triangle inequality $\norm{u+v}_2\le\norm{u}_2+\norm{v}_2$ by first squaring both sides.\end{exercise}
\begin{exercise}\label{ex:10.6}If $u_1,\ldots,u_k$ are nonzero pairwise orthogonal vectors, show that $a_1u_1+\cdots+a_ku_k=0$ forces every coefficient $a_j=0$.\end{exercise}


% ===== Week 11 =====

\chapter{Distance, error, and the meaning of a limit}\label{ch:limits}
\goal{Replace a visual impression of convergence by an explicit tolerance guarantee. Understand why small successive changes do not suffice. Identify completeness as the assumption that prevents a sequence from approaching a missing point.}

\section{Distance begins with absolute value}
Recall the magnitude notation from Chapter~2. The absolute value $|x|$ is $x$ when $x\ge0$ and $-x$ when $x<0$. In particular, $|-3|=3$: the minus sign in the definition makes a negative input positive. Distance between real numbers is $|x-y|$. For $x=-2$ and $y=3$, it is $|-5|=5$, the same as $|y-x|$.

We now use this familiar distance to understand a more general idea. The plan is not to learn a new kind of arithmetic. It is to isolate the properties of distance that an argument needs, so the argument can work on objects other than single real numbers.

The triangle inequality is
\[
|x+y|\le|x|+|y|.
\]
To derive it, note that $-|x|\le x\le|x|$ and $-|y|\le y\le|y|$. Adding gives $-(|x|+|y|)\le x+y\le|x|+|y|$, which is equivalent to the inequality. Replacing $x$ by $a-b$ and $y$ by $b-c$ gives
\[
|a-c|\le|a-b|+|b-c|.
\]
The sum of the two leg lengths is an upper bound on the direct distance, not generally a lower bound. This is why triangle inequalities appear so often in error estimates: introduce an intermediate quantity that is easier to compare with each endpoint.

If $|x-a|\le\eps$ and $|y-b|\le\delta$, then
\[
|(x+y)-(a+b)|\le |x-a|+|y-b|\le\eps+\delta.
\]
Errors add in this estimate. We have not assumed they actually point in the same direction, so the bound may be conservative. A valid upper bound need not be attained in every example.

\section{A metric retains only the properties we use}
A function on $X\times X$ takes a pair of points as input. We write its output as $d(x,y)$ and want that output to behave as a distance. The following definition specifies what this means.
\par\noindent\begin{minipage}{\linewidth}
\begin{definition}[Metric]\label{def:metric}
A metric on a set $X$ assigns a real number $d(x,y)$ to every pair $x,y\in X$ and satisfies these requirements for all $x,y,z\in X$:
\begin{enumerate}
\item $d(x,y)\ge0$, and $d(x,y)=0$ holds exactly when $x=y$.
\item $d(x,y)=d(y,x)$.
\item $d(x,z)\le d(x,y)+d(y,z)$.
\end{enumerate}
A set together with a specified metric is called a metric space and is written $(X,d)$.
\end{definition}
\end{minipage}\par

The first requirement rules out negative distances and prevents different points from having zero distance. The second says that reversing the comparison does not change its value. The third is the triangle inequality: a detour through $y$ bounds the direct distance from $x$ to $z$.

To \emph{use} a metric already given in a theorem, we may use all three requirements. To \emph{propose} a new metric, we must check all three. These are the two directions of using a definition from Chapter~1. Merely naming a formula ``distance'' does not establish that it meets them.

The real line with $d(x,y)=|x-y|$ is an example. In $\RR^n$, the Euclidean distance is $\norm{x-y}_2$, whose triangle inequality was derived in Chapter~10 and revisited in Exercise~10.5. Another useful choice is
\[
d_\infty(x,y)=\max_{1\le j\le n}|x_j-y_j|.
\]
It records the worst coordinate error. For every coordinate $j$,
\[
|x_j-z_j|\le|x_j-y_j|+|y_j-z_j|\le d_\infty(x,y)+d_\infty(y,z).
\]
Taking the maximum over $j$ proves its triangle inequality. This example shows that a distance can be selected to match the question being controlled.

A norm measures the size of a vector, is nonnegative and zero exactly at the zero vector, and satisfies $\norm{cx}=|c|\norm{x}$, and the triangle inequality. It produces a metric by $d(x,y)=\norm{x-y}$. A general metric space need not allow addition of points. Our later proof will be stronger if it uses only distances when vector subtraction is unnecessary.

\section{Convergence promises eventual control}
A sequence is a list $x_0,x_1,x_2,\ldots$ indexed by natural numbers. It converges to $x_*$ when, for every tolerance $\eps>0$, all sufficiently late terms lie within that tolerance:
\[
\forall\eps>0\ \exists N\in\NN\ \forall n\ge N:\quad d(x_n,x_*)<\eps.
\]
The threshold $N$ may depend on $\eps$. It must work for every later $n$, not just for one convenient term.

Read the order as a challenge and response. First someone asks for a positive tolerance, such as $0.01$. We must produce a threshold $N$. Then they may choose \emph{any} index at or beyond $N$, and its distance from the same proposed limit must be below $0.01$. They may later ask for a smaller tolerance; we may then choose a larger threshold.

The point $x_*$ is fixed before these tolerance requests. The star is a label for the proposed destination, not multiplication and not a special numerical value. The notation $x_n\to x_*$ records convergence. It does not say the sequence ever equals its limit. For $x_n=1/(n+1)$, every term is positive even though the limit is zero.

Choosing one late term near the destination is not enough. The sequence $0,1,0,1,\ldots$ repeatedly returns to zero, but has no tail in which \emph{every} term is within $1/4$ of zero. The quantifier ``every $n\ge N$'' rules out that behavior.

For $x_n=1/(n+1)$ on the real line, take a positive integer $N>1/\eps$. Then for $n\ge N$,
\[
0<x_n=\frac1{n+1}\le\frac1{N+1}<\eps.
\]
This proves $x_n\to0$. A finite prefix, however irregular, does not affect the limit because the definition only controls sufficiently late terms.

For $x_n=1-2^{-n}$, the error relative to $1$ is $2^{-n}$. To show it tends to zero without logarithms, observe by induction that $2^n\ge n+1$. Hence $2^{-n}\le1/(n+1)$, and the previous tolerance argument applies.

\section{Limits are unique}
\par\noindent\begin{minipage}{\linewidth}
\begin{proposition}\label{thm:limitunique}
A sequence in a metric space has at most one limit.
\end{proposition}
\end{minipage}\par

\begin{proof}
Suppose $x_n\to a$ and $x_n\to b$, and assume $a\ne b$. Set $D=d(a,b)>0$ and choose the tolerance $D/3$. The first convergence supplies an index beyond which $d(x_n,a)<D/3$; the second supplies an index beyond which $d(x_n,b)<D/3$. Take $n$ beyond both indices. By symmetry and the triangle inequality,
\[
D=d(a,b)\le d(a,x_n)+d(x_n,b)<2D/3.
\]
This contradicts $D>0$. Therefore $a=b$.
\end{proof}

Notice why we used one index beyond both thresholds. Two separate late terms would not give the same two-step path from $a$ through $x_n$ to $b$. Quantifier control is doing mathematical work here.

\section{Cauchy sequences do not need a known destination}
A sequence is Cauchy when its sufficiently late terms are mutually close:
\[
\forall\eps>0\ \exists N\ \forall m,n\ge N:\quad d(x_m,x_n)<\eps.
\]
The destination does not appear. This is useful when we are trying to prove that an iterative construction has a limit rather than starting with a known candidate.

Every convergent sequence is Cauchy. Here is the argument with both late indices visible. Given $\eps>0$, ask convergence for the smaller tolerance $\eps/2$. It supplies one threshold $N$ such that every term with index at least $N$ is within $\eps/2$ of $x_*$. Now choose any $m,n\ge N$. Both comparisons are available, so
\[
 d(x_m,x_n)\le d(x_m,x_*)+d(x_*,x_n)
 <\frac\eps2+\frac\eps2=\eps.
\]
The same threshold works for \emph{every pair} of later terms, proving the Cauchy condition. We split the tolerance in half because the triangle inequality will add two errors. The converse requires an additional property of the space: mutual closeness does not by itself put a destination inside the allowed set.

Small consecutive changes alone do not imply the Cauchy condition. Let $s_n=\sum_{j=1}^n1/j$, with $s_0=0$. Then $s_{n+1}-s_n=1/(n+1)\to0$. But for $n\ge1$,
\[
s_{2n}-s_n=\sum_{j=n+1}^{2n}\frac1j\ge n\frac1{2n}=\frac12.
\]
There are $n$ terms, each at least $1/(2n)$. No threshold can make all tail differences smaller than $1/4$. Thus the sequence is not Cauchy. We need control of accumulated changes, not just individual steps.

\subsection*{One step versus any two late steps}
For the harmonic sums, inspect $n=4$: the difference $s_8-s_4$ contains the four terms $1/5,1/6,1/7,1/8$, each at least $1/8$, so the difference is at least $1/2$. A single step after index four is smaller, but many small steps accumulate.

To refute the Cauchy condition, fix the tolerance $1/4$ once and for all. For \emph{any} proposed threshold $N$, choose an integer $n\ge\max\{N,1\}$ and set $m=2n$. Both indices are at least $N$, yet their distance is at least $1/2$, not below $1/4$. This is the exact negation of eventual control over every late pair. We did not need to prove a formula for the infinite sum or even discuss such a sum.

\section{Completeness excludes missing limits}
A metric space is complete when every Cauchy sequence in it converges to a point in the same space. We take completeness of the real numbers as a standard starting property of the real-number system. It is not something inferred from a numerical plot.

Completeness is a property of the chosen space, not an extra adjective for one sequence. To use it, first prove the sequence is Cauchy; then the definition guarantees a limit belonging to that space. To show a space is \emph{not} complete, it is enough to give one Cauchy sequence whose only possible destination has been excluded.

We do not prove completeness of the real numbers here. It is an explicit starting fact about the number line, just as the existence of nonnegative square roots was an explicit starting fact earlier. We do prove the particular consequences of completeness that we use. This keeps the foundational boundary visible instead of hiding it behind an unsupported assertion that a limit exists.

The interval $(0,1)$ with ordinary distance is not complete. The sequence $x_n=2^{-(n+1)}$ lies in that interval and is Cauchy: for $m,n\ge N$, both terms lie between zero and $2^{-(N+1)}$, so their distance is at most $2^{-(N+1)}$, which tends to zero. Its real limit is zero, but zero is excluded from $(0,1)$. It cannot converge to a different point of the interval because real limits are unique.

A subset $C$ of a metric space is closed when it contains the limit of every convergent sequence of its points. We use this sequential definition here. If the ambient space is complete and $C$ is closed, then $C$ is complete with the inherited distance: a Cauchy sequence in $C$ converges in the ambient space, and closedness puts its limit back in $C$.

In particular, the interval $[0,1]$ is complete. To check closedness rather than assume it, suppose a sequence of points in $[0,1]$ converges in the real line to a number $L$.

If $L<0$, choose the positive tolerance $-L/2$. For a sufficiently late term, $|x_n-L|<-L/2$, so $x_n<L-L/2=L/2<0$. This contradicts $x_n\in[0,1]$. If $L>1$, instead choose tolerance $(L-1)/2$. A sufficiently late term then satisfies
\[
 x_n>L-(L-1)/2=(L+1)/2>1,
\]
again contradicting membership. Both excluded cases are impossible, so $0\le L\le1$. The interval is closed in the complete real line, and hence is complete with the same distance.

\section{A geometric bound controls an entire tail}
For $0<q<1$, the finite geometric sum satisfies
\[
1+q+\cdots+q^{r-1}=\frac{1-q^r}{1-q}.
\]
To see the cancellation, first take $r=4$ and give the sum a name, $S=1+q+q^2+q^3$. Then
\begin{align*}
 S&=1+q+q^2+q^3,\\
 qS&=\phantom{1+{}}q+q^2+q^3+q^4,\\
 S-qS&=1-q^4.
\end{align*}
The left side is $(1-q)S$, so $S=(1-q^4)/(1-q)$. For general positive $r$, the common terms are $q$ through $q^{r-1}$ and the unmatched end terms are $1$ and $q^r$. The same subtraction gives $(1-q)S=1-q^r$. Since $1-q>0$, division is permitted. Therefore
\[
\sum_{j=n}^{m-1}q^j=q^n\frac{1-q^{m-n}}{1-q}\le\frac{q^n}{1-q}\qquad(m>n).
\]
The bound must get small as the starting index $n$ grows. We therefore need to justify $q^n\to0$ for every fixed $0<q<1$, not just for $q=1/2$.

Put $c=1/q-1>0$, so $1/q=1+c$. We prove $(1+c)^n\ge1+nc$ by induction. At $n=0$ both sides are one. If it holds at $n$, multiplication by $1+c>0$ gives
\[
 (1+c)^{n+1}\ge(1+nc)(1+c)
 =1+(n+1)c+nc^2\ge1+(n+1)c.
\]
Thus the bound holds for every $n$. Taking reciprocals of positive quantities reverses order, giving
\[
 0<q^n\le\frac1{1+nc}.
\]
For a requested tolerance $\eps>0$, choose a positive integer $N>1/(c\eps)$. If $n\ge N$, then $1+nc>nc\ge Nc>1/\eps$, so $q^n<\eps$. This is the required convergence proof. The threshold may depend on the fixed factor $q$ through $c$, and on the requested tolerance.

If a sequence satisfies $d(x_{j+1},x_j)\le Dq^j$, repeated triangle inequalities now give
\[
d(x_m,x_n)\le\sum_{j=n}^{m-1}d(x_{j+1},x_j)
\le\frac{Dq^n}{1-q}.
\]
The right side tends to zero independently of how late $m$ is. This proves the Cauchy condition.

There are two thresholds we must not confuse. The integer $n$ is the starting point of a tail comparison, and $m>n$ can be arbitrarily far beyond it. Our bound contains $n$ but not $m$, so one choice of $n$ controls even very long tails.

For $D>0$, a given tolerance $\eps$ asks for $q^n<\eps(1-q)/D$. The preceding geometric-limit argument supplies a threshold $N$ for that positive number. For any two indices at least $N$, put the smaller one in the role of $n$; symmetry handles their order, and equal indices have zero distance. If $D=0$, every step has zero length, so the sequence is constant. This also covers the case where dividing by $D$ would not have been permitted. We have found exactly the stronger control missing from the harmonic-sum example.

\begin{lab}{Asking for the quantifier that failed}
\noindent\textbf{Deliberately flawed reply.} ``The updates get arbitrarily small, so the iterates converge.''

\noindent\textbf{Learner.} ``Compare consecutive differences with the Cauchy condition. Test the partial sums of $1/j$. Then give a stronger hypothesis that controls every tail, not just one step.''

\noindent\textbf{Model repair.} A geometric bound on successive distances gives the displayed tail estimate. Completeness then supplies a limit in the space. Record these as separate obligations: tail control and existence of an included destination.
\end{lab}

\section*{Exercises}
\begin{exercise}\label{ex:11.1}If $|x-a|\le0.02$ and $|y-b|\le0.03$, bound $|(2x-y)-(2a-b)|$.\end{exercise}
\begin{exercise}\label{ex:11.2}Prove $3/(n+1)\to0$ with a stated choice of threshold for each positive tolerance.\end{exercise}
\begin{exercise}\label{ex:11.3}Prove $d_\infty$ vanishes exactly for equal vectors and is symmetric. Together with the main-text argument, this completes the metric verification.\end{exercise}
\begin{exercise}\label{ex:11.4}For $d(x_{j+1},x_j)\le2(1/3)^j$, give a bound for $d(x_m,x_n)$ valid for every $m>n$.\end{exercise}
\begin{exercise}\label{ex:11.5}Give a Cauchy sequence in $(2,3)$ that has no limit in that space. Prove both assertions.\end{exercise}
\begin{exercise}\label{ex:11.6}A set is bounded above if some real number exceeds or equals every member. Its supremum is its least upper bound. Show that $(0,1)$ has supremum $1$ but no maximum. Explain why replacing ``supremum'' by ``maximum'' would change the claim.\end{exercise}


% ===== Week 12 =====

\chapter{The contraction theorem: an unfamiliar proof made readable}\label{ch:contraction}
\goal{Assemble a complete theorem from a small prerequisite chain. Prove existence, uniqueness, convergence, and an error certificate as separate conclusions. Test each hypothesis by removing it.}

\section{From solving an equation to repeating a map}
A fixed point of a map $T:X\to X$ is a point $x_*$ satisfying $T(x_*)=x_*$. It is unchanged by the update. The equation $x=(x+1)/2$ from Chapter~4 is a fixed-point equation. Starting from $x_0=0$ and repeatedly applying the right side gives a sequence approaching its solution.

This suggests an algorithm, but not yet a theorem. Some update rules oscillate or diverge. Some have several fixed points, and some approach a point excluded from the domain. We need assumptions that explain why the iteration works, rather than evidence from one favorable plot.

A map is a contraction with factor $q$ when $0\le q<1$ and
\[
d(Tx,Ty)\le q\,d(x,y)\qquad\text{for every }x,y\in X.
\]
It uniformly reduces the distance between any two inputs. The word uniformly is encoded by one fixed factor $q$ that works for all pairs. The factor cannot secretly vary toward one as the inputs change unless another argument supplies the needed control.

\subsection*{Read the two inputs in the contraction inequality}
The inputs $x$ and $y$ are arbitrary points of the space. They are not assumed to be successive iterates and not assumed to be close to the fixed point. The same $q$ must work before either input is chosen. For $T(x)=(x+1)/2$,
\[
 |T(x)-T(y)|=\left|\frac{x+1-y-1}{2}\right|
 =\frac12|x-y|,
\]
so $q=1/2$ works for all real pairs. The constant terms cancel when two outputs are compared.

The notation $T:X\to X$ is a \emph{self-map} requirement: each output is a permitted input for the next application. It is not automatic from an algebraic formula. Restricting the input set may improve an inequality while causing outputs to leave that set. We must check both the domain condition and the distance inequality.

\section{The theorem and its dependency map}
\par\noindent\begin{minipage}{\linewidth}
\begin{theorem}[Contraction theorem]\label{thm:banach}
Let $(X,d)$ be a nonempty complete metric space, and let $T:X\to X$ satisfy the contraction inequality with $0\le q<1$. Then $T$ has exactly one fixed point $x_*$. For every initial point $x_0\in X$, the iterates $x_{n+1}=T(x_n)$ converge to $x_*$. If $0<q<1$, they satisfy
\[
d(x_n,x_*)\le\frac{q^n}{1-q}\,d(x_1,x_0).
\tag{12.1}\label{eq:apriori}
\]
For every point $z\in X$, whether or not it is an iterate,
\[
d(z,x_*)\le\frac{d(z,Tz)}{1-q}.
\tag{12.2}\label{eq:residual}
\]
\end{theorem}
\end{minipage}\par


Before starting, separate the theorem's promises. Existence asks for a point that is unchanged by $T$. Uniqueness asks why two such points cannot differ. Convergence asks whether the sequence produced from \emph{any} initial point approaches that point. The error bounds ask how far a finite approximation can be from it. Proving one of these does not automatically prove the others.

The first bound uses $x_0$ and its first update $x_1$, so it describes the iterated sequence. The second uses an arbitrary $z\in X$, so it can certify an approximation produced by another procedure. In both inequalities, $1-q$ is positive. That positive margin is a hypothesis of the theorem; it is what makes division legitimate and the error bound finite.

Here is the whole proof before its details. The contraction inequality makes successive step lengths decrease geometrically. The geometric-sum bound makes the iterates Cauchy. Completeness supplies a limit in the space. A distance comparison shows that this limit is fixed. Another comparison rules out two fixed points. Finally, the same estimates quantify the error.

Only a few ingredients are needed: induction, the triangle inequality, geometric sums, the definition of Cauchy, and completeness. We do not need calculus or an entire course in analysis. This is what a sufficient background plan looks like when it is tied to an exact proof.

\section{The proof, one obligation at a time}
\begin{proof}
First handle $q=0$. The contraction inequality gives $d(Tx,Ty)=0$ for every pair, so $T$ is constant. Choose $x_0\in X$ and set $c=T(x_0)$. Since $c\in X$, constancy gives $T(c)=T(x_0)=c$. Thus $c$ is fixed, every iterate after the first equals $c$, and no other fixed point is possible. The residual estimate is exact because $d(z,Tz)=d(z,c)$. We now assume $0<q<1$.

\textbf{1. Construct the sequence and control one step.} Choose any $x_0\in X$ and define $x_{n+1}=T(x_n)$. The self-map assumption ensures that every iterate remains in $X$. Set $D=d(x_1,x_0)$. We claim
\[
d(x_{n+1},x_n)\le q^nD\qquad(n\ge0).
\]
At $n=0$, this is equality. If the bound holds at $n$, the contraction inequality gives
\[
d(x_{n+2},x_{n+1})=d(Tx_{n+1},Tx_n)
\le q\,d(x_{n+1},x_n)\le q^{n+1}D.
\]
Induction proves the claim.

\textbf{2. Control every tail.} For $m>n$, apply the triangle inequality through the intermediate iterates and use the geometric-sum estimate from Chapter~11:
\[
d(x_m,x_n)\le\sum_{j=n}^{m-1}d(x_{j+1},x_j)
\le D\sum_{j=n}^{m-1}q^j
\le\frac{Dq^n}{1-q}.
\]
If $D=0$, all iterates are the same and the Cauchy conclusion is immediate. Otherwise, because $q^n\to0$, for each $\eps>0$ there is $N$ with $Dq^N/(1-q)<\eps$. For any $m,n\ge N$, assume $m>n$ by symmetry, or handle $m=n$ with distance zero. The bound gives $d(x_m,x_n)<\eps$. Thus the sequence is Cauchy.

\textbf{3. Obtain a point in the space.} Completeness supplies $x_*\in X$ with $x_n\to x_*$. This is precisely where completeness is used. The preceding bounds alone would not ensure that a limit belongs to an incomplete space.

\textbf{4. Prove that the limit is fixed.} For every $n$,
\begin{align*}
d(Tx_*,x_*)&\le d(Tx_*,Tx_n)+d(Tx_n,x_*)\\
&\le q\,d(x_*,x_n)+d(x_{n+1},x_*).
\end{align*}
For any $\eps>0$, choose $n$ so late that both distances on the final line are smaller than $\eps/(q+1)$. This is possible because both the sequence and its one-step shift converge to $x_*$. The displayed bound is then strictly smaller than $\eps$. A nonnegative number smaller than every positive $\eps$ must be zero; otherwise choose $\eps$ equal to that number.

More explicitly, the quantity on the left, $d(Tx_*,x_*)$, does not depend on $n$. We have shown that for \emph{each} positive tolerance we can choose a late enough $n$ making an upper bound smaller than that tolerance. If the left side had some positive value $D$, requesting tolerance $D/2$ would give $D<D/2$, which is impossible. We are not replacing an approximate equality by an exact equality after a finite number of steps. Hence $d(Tx_*,x_*)=0$ and $Tx_*=x_*$.

\textbf{5. Prove uniqueness.} Suppose $y_*$ is another fixed point. Then
\[
d(x_*,y_*)=d(Tx_*,Ty_*)\le q\,d(x_*,y_*).
\]
Subtracting the rightmost quantity gives $(1-q)d(x_*,y_*)\le0$. Both factors are nonnegative and $1-q>0$, so the distance must be zero. Thus $x_*=y_*$.

\textbf{6. Certify the error.} Fix $n$. For any later $m$,
\[
d(x_n,x_*)\le d(x_n,x_m)+d(x_m,x_*)
\le\frac{Dq^n}{1-q}+d(x_m,x_*).
\]
The last term can be made smaller than any positive tolerance. Therefore $d(x_n,x_*)\le Dq^n/(1-q)$: if the left side exceeded that bound, choose a tolerance smaller than the positive excess. This proves \eqref{eq:apriori}.

The logical form is worth recognizing. We have a fixed quantity $A=d(x_n,x_*)$ and a proposed bound $B=Dq^n/(1-q)$. For every $\eta>0$, we can show $A<B+\eta$ by choosing a later $m$. If $A>B$, choose $\eta=(A-B)/2$; then $B+\eta<A$, contradicting that estimate. Therefore $A\le B$. An arbitrarily small leftover error is removed by this argument, not by pretending it was exactly zero at a finite step.

For any $z\in X$, insert $Tz$ as an intermediate point and use $Tx_*=x_*$:
\[
d(z,x_*)\le d(z,Tz)+d(Tz,Tx_*)
\le d(z,Tz)+q\,d(z,x_*).
\]
Subtract the final distance term and divide by $1-q>0$ to obtain \eqref{eq:residual}. All conclusions now follow.
\end{proof}

\section{What the two error estimates are for}
The a priori estimate predicts how many steps may be needed before computing all of them. For $T(x)=(x+1)/2$, $x_0=0$, we have $q=1/2$ and $D=1/2$, so the bound is $2^{-n}$. At $n=10$, it is $1/1024<0.001$. The exact formula from Chapter~5 shows that this bound is attained.

The residual $d(z,Tz)$ measures failure to satisfy the fixed-point equation. The residual estimate can certify a point obtained by any method, not just the prescribed iteration. It is especially useful when the exact solution is unknown but one update can be evaluated reliably. A small residual is meaningful only together with the stability factor $1/(1-q)$. As $q$ approaches one, the same residual permits a larger solution error.

For an iterate $x_n$ with $n\ge1$, one also has
\[
d(x_n,x_*)\le\frac{q}{1-q}\,d(x_n,x_{n-1}).
\]
Indeed, its residual is $d(x_n,x_{n+1})\le qd(x_{n-1},x_n)$, and then \eqref{eq:residual} applies. This is an a posteriori estimate: it uses the observed size of the latest step.

\section{A nonlinear example without derivatives}
Consider $T(x)=(x^2+1)/3$ on $X=[0,1]$. First check the space and the map, not just the inequality. The interval is complete by Chapter~11. For $x\in[0,1]$, $x^2\in[0,1]$, so $T(x)\in[1/3,2/3]\subseteq X$. Thus $T$ is a self-map.

For $x,y\in[0,1]$,
\[
|T(x)-T(y)|=\frac{|x^2-y^2|}{3}
=\frac{|x-y|\,|x+y|}{3}\le\frac23|x-y|.
\]
The factor $q=2/3$ works for all inputs in the interval. The theorem supplies a unique fixed point in $[0,1]$.

Starting at zero gives $x_1=1/3$, $x_2=10/27$, and $x_3=829/2187$. For $z=x_2$, the residual is
\[
\left|\frac{10}{27}-\frac{829}{2187}\right|=\frac{19}{2187}.
\]
The certified error is at most $3\cdot19/2187=19/729$. This conclusion does not require knowing the exact solution. As a separate algebraic check, the fixed-point equation gives $x^2-3x+1=0$, hence $(2x-3)^2=5$ and $x=(3\pm\sqrt5)/2$. Only the minus sign lies in $[0,1]$.

For the fraction calculation, start from $x_1=1/3$:
\[
 x_2=\frac{(1/3)^2+1}{3}
 =\frac{1/9+9/9}{3}=\frac{10}{27}.
\]
The next update squares $10/27$, not its rounded decimal value:
\[
 x_3=\frac{100/729+729/729}{3}=\frac{829}{2187}.
\]
To subtract $x_2$ from $x_3$, use denominator $2187=27\cdot81$: $x_2=810/2187$. The residual is $(829-810)/2187=19/2187$. Finally $1/(1-q)=1/(1-2/3)=3$, giving the certificate $19/729$.

The algebraic check has a different purpose. Multiplying the fixed-point equation by three gives $x^2-3x+1=0$. Multiply by four and add five to complete the square: $(2x-3)^2=5$. The square-root definition gives $2x-3=\sqrt5$ or $-\sqrt5$. Since $2<\sqrt5<3$, the minus choice produces a value between zero and $1/2$, whereas the plus choice exceeds two. This proves which candidate belongs to the interval without requiring a memorized quadratic formula.

\section{Remove the hypotheses one at a time}
If completeness is omitted, the map $T(x)=x/2$ on $(0,1)$ is a contraction and a self-map but has no fixed point there: the only real solution is zero, which is excluded.

If strict contraction is replaced by $q=1$, the identity map on $\RR$ has every point fixed, while $T(x)=x+1$ has none. Both preserve distances. The boundary $q=1$ therefore does not merely worsen the numerical rate; it can destroy existence or uniqueness.

If the self-map condition is omitted, $T(x)=x/2+1$ on inputs $[0,1]$ has factor $1/2$ but leaves the interval. Its real fixed point is $2$, outside the intended space. A correct inequality cannot repair a missing domain condition.

\begin{lab}{Demanding the minimum sufficient background}
Give the theorem statement to the assistant and declare that you know Chapters~1--10 but not completeness. Request a dependency plan for this proof, not a survey of analysis. The required additions are metrics, convergence, Cauchy sequences, completeness, and a finite geometric-tail estimate. Ask for one example, one nonexample, and one diagnostic question for each. Then reconstruct Steps~2--4 above. Reject a reply that simply says ``by continuity'' without either defining continuity or deriving the displayed distance bound.
\end{lab}

\section*{Exercises}
\begin{exercise}\label{ex:12.1}For $T(x)=(3x+1)/4$ on $\RR$, find its fixed point, contraction factor, and exact iterates from $x_0=0$.\end{exercise}
\begin{exercise}\label{ex:12.2}For $q=3/4$ and $d(z,Tz)\le0.002$, give a certified upper bound on $d(z,x_*)$.\end{exercise}
\begin{exercise}\label{ex:12.3}Show that $T(x)=(x^2+1)/3$ maps $[1/3,2/3]$ into itself and has contraction factor $4/9$ there. Explain why this can improve a certificate after the first step.\end{exercise}
\begin{exercise}\label{ex:12.4}Identify the failed theorem hypothesis for each of the three groups of counterexamples in the preceding section. Explain why the other relevant assumptions still hold.\end{exercise}
\begin{exercise}\label{ex:12.5}For $T(x)=ax+b$ on $\RR$ with $|a|<1$, find the fixed point and prove the exact formula $x_n-x_*=a^n(x_0-x_*)$. What changes when $a<0$?\end{exercise}
\begin{exercise}\label{ex:12.6}Suppose two self-maps $T,S$ of the same nonempty complete metric space are contractions with factor at most $q<1$, and $d(Tx,Sx)\le\eta$ for every $x$. Prove their fixed points satisfy $d(x_T,x_S)\le\eta/(1-q)$.\end{exercise}


% ===== Week 13 =====

\chapter{Averaging, randomness, and existence}\label{ch:probability}
\goal{Understand finite probability as weighted counting. Prove existence by averaging without confusing it with a guarantee for every outcome. Learn when an existence argument can be turned into a deterministic construction.}

\section{A finite probability model}
A finite sample space $\Omega$ is a set of possible outcomes. Assign each outcome $\omega$ a probability $p(\omega)\ge0$ with $\sum_{\omega\in\Omega}p(\omega)=1$. For an event $A\subseteq\Omega$, its probability is
\[
\Prob(A)=\sum_{\omega\in A}p(\omega).
\]
In a uniform model, all outcomes have equal probability $1/|\Omega|$, so probability is simply the fraction of outcomes in the event. A probability model describes a stipulated mathematical experiment. It is not automatically a justified model of real-world uncertainty.

For two independently labeled coins, the outcomes $00,01,10,11$ are equally likely. The event that the labels differ has probability $2/4=1/2$. Independence means the joint probability factors into the individual probabilities; here each ordered pair has probability $(1/2)(1/2)$. We will use independent labels to construct a graph partition, not assume independence of every event arising from that partition.

\subsection*{The outcomes must be specified before they are counted}
For two labeled coins, $01$ and $10$ are different outcomes: in one, the first coin is zero and the second is one; in the other their values are reversed. Calling both outcomes ``one of each'' would merge two equally likely possibilities into one description. A uniform probability assignment on the three descriptions ``two zeros,'' ``one of each,'' and ``two ones'' would therefore describe a different experiment.

An event is a set of outcomes. The complement of $A$ is the set of outcomes not in $A$, written $\Omega\setminus A$. The two sets are disjoint and together contain every outcome, so
\[
 \Prob(A)+\Prob(\Omega\setminus A)=\sum_{\omega\in\Omega}p(\omega)=1.
\]
Consequently, if $\Prob(A)<1$, the complementary event has positive probability. In a finite model a positive-probability event contains at least one outcome: an empty event would have probability zero. We will use exactly this implication in an existence argument below.

The symbol $\Omega$, read ``capital omega,'' is just a name for the sample space. A lowercase $\omega$ names one of its members. No further meaning should be inferred from the Greek letters.

\section{Expectation is a weighted average}
A random variable $X$ is simply a real-valued function on $\Omega$. Its expectation is
\[
\E X=\sum_{\omega\in\Omega}p(\omega)X(\omega).
\]
The word random describes how an input outcome is chosen; once the outcome is specified, $X(\omega)$ is an ordinary number.

For the two-coin experiment, let $X$ count how many labels are one. Its values are $X(00)=0$, $X(01)=1$, $X(10)=1$, and $X(11)=2$. The expectation is
\[
 \E X=\tfrac14\cdot0+\tfrac14\cdot1+\tfrac14\cdot1+\tfrac14\cdot2=1.
\]
This is the average of values, not an additional outcome. To see that the average need not itself occur, take one fair coin and a variable that is zero for label zero and two for label one. Its expectation is still one, although it never takes value one.

The sum $X+Y$ of two random variables means the function whose value at an outcome $\omega$ is $X(\omega)+Y(\omega)$. The same outcome is used for both evaluations. In the expectation calculation that follows, distributing the sum of these two ordinary values is therefore legitimate even when the variables are dependent.

For random variables $X,Y$, finite-sum algebra gives
\begin{align*}
\E(X+Y)&=\sum_\omega p(\omega)(X(\omega)+Y(\omega))\\
&=\sum_\omega p(\omega)X(\omega)+\sum_\omega p(\omega)Y(\omega)\\
&=\E X+\E Y.
\end{align*}
Likewise $\E(cX)=c\E X$. These identities require no independence between $X$ and $Y$. Their values may be tightly related; distributing a finite sum is still valid.

There is one more property of finite averages that we need. If $X(\omega)\le Y(\omega)$ for every outcome, then $\E X\le\E Y$. Each weight $p(\omega)$ is nonnegative, so multiplying the pointwise inequality by that weight preserves its direction. Adding all these inequalities gives the expectation comparison. This is called \emph{monotonicity of expectation}. It is the justification whenever we take expectations of an inequality.

The word \emph{pointwise} means ``for each fixed input outcome separately.'' It does not mean only on average. A pointwise inequality implies an average inequality by this argument, but the reverse implication is generally false.

At least one outcome with positive probability has $X(\omega)\ge\E X$. Otherwise every positive-probability value would be smaller than the expectation, and their weighted average would also be smaller than it, a contradiction. This is the same average argument used for graph degrees, now with weights.

\section{Indicators convert events into quantities}
For an event $A$, define its indicator $I_A$ to be $1$ on $A$ and $0$ outside. Then
\[
\E I_A=\sum_{\omega\in A}p(\omega)=\Prob(A).
\]
If $A_1,\ldots,A_m$ are events, the variable $I_{A_1}+\cdots+I_{A_m}$ counts how many occur. Its expectation is $\sum_j\Prob(A_j)$, even if the events overlap.

For example, in the two-coin experiment let $A$ say the first label is one and let $B$ say the second label is one. At outcome $11$, both indicators are one, so $I_A+I_B=2$. The indicator of their union is only one: it records whether at least one event occurs, not how many occur. At $00$, all three quantities are zero. At $01$ or $10$, the sum and union indicator are both one. Thus the union indicator is no larger than the sum at every outcome. For several events the same reasoning gives the next inequality.

Pointwise, $I_{A_1\cup\cdots\cup A_m}\le\sum_j I_{A_j}$: if no event occurs both sides are zero, and if at least one occurs the right side is at least one. Taking expectations gives the union bound
\[
\Prob(A_1\cup\cdots\cup A_m)\le\sum_{j=1}^m\Prob(A_j).
\]
If the sum on the right is strictly less than one, the union does not include every positive-probability outcome. Thus some outcome avoids every event. Thinking of the $A_j$ as bad events turns an upper bound into an existence theorem.

\pauseq{Does a sum of bad-event probabilities equal to one guarantee an outcome avoiding them all?}{No. For a fair bit, the events ``the bit is zero'' and ``the bit is one'' each have probability $1/2$, but their union is the entire sample space. The strict inequality supplies the spare probability needed by the argument.}

\section{Every graph has a large cut}
A cut of an undirected graph partitions its vertices into two sets. An edge crosses the cut when its endpoints are in different sets. We want many crossing edges.

\par\noindent\begin{minipage}{\linewidth}
\begin{theorem}\label{thm:cut}
Every finite simple undirected graph with $m$ edges has a cut containing at least $m/2$ crossing edges.
\end{theorem}
\end{minipage}\par

\noindent\textit{Plan.} Assign each vertex one of two sides independently with equal probabilities. Count crossing edges by indicators, then use the fact that some outcome is at least the average.
\begin{proof}
The outcomes are all assignments of the two labels $0,1$ to the vertices, each equally likely. For each edge $e=\{u,v\}$, let $I_e$ indicate that the endpoint labels differ. Because $u,v$ are distinct and their labels are independent, their four ordered label pairs are equally likely. Two pairs differ, so $\E I_e=1/2$.

The total number of crossing edges is $X=\sum_{e\in E}I_e$. Linearity of expectation yields
\[
\E X=\sum_{e\in E}\E I_e=\frac m2.
\]
Some labeling has $X\ge\E X=m/2$, and its two label classes form the desired cut. If the graph has no edges, any labeling gives zero and the same conclusion holds.
\end{proof}

The indicators for different edges need not be jointly independent, and the proof does not claim they are. For a triangle, the three crossing indicators cannot all be one, because three vertices cannot each have a label different from the other two. Yet their expectations still add.

When $m$ is odd, the theorem's integer-valued conclusion is at least the next integer above $m/2$. A triangle has $m=3$, so some cut has at least two edges; indeed, putting one vertex on one side and two on the other gives exactly two.

\section{Turning the average argument into a construction}
The random proof shows that a favorable labeling exists. We can also construct one by deciding the labels sequentially while keeping the expected final cut size from decreasing.

At a partial assignment, let $C$ be the number of crossing edges whose endpoints are both already assigned. Let $U$ be the number of edges with at least one unassigned endpoint. If all remaining labels were chosen independently and fairly, every unresolved edge would cross with probability $1/2$, whether one or both endpoints are unassigned. The expected final count is therefore
\[
F=C+\frac U2.
\]
At the beginning, $F=m/2$.

Now assign the next vertex $v$. Suppose it has $\ell$ already assigned neighbors on the left and $r$ on the right. Placing it on the left resolves $\ell+r$ previously unresolved edges and makes $r$ of them cross. The change in $F$ is
\[
r-\frac{\ell+r}{2}=\frac{r-\ell}{2}.
\]
Placing it on the right gives change $(\ell-r)/2$. At least one is nonnegative. Choose the side opposite the larger number of already assigned neighbors. Edges from $v$ to still unassigned vertices remain unresolved and contribute the same $1/2$ as before.

After all vertices are assigned, $U=0$ and $F$ is the actual cut size. Since $F$ never decreased from $m/2$, the final cut has at least $m/2$ edges. This is a deterministic construction with a maintained lower bound. The technique is an instance of conditional expectation, but its entire calculation is visible here; no general probability theory is needed.

\subsection*{A full trace of the maintained average}
Take a triangle with vertices $a,b,c$ and three edges. Initially no labels are assigned, so $C=0$, $U=3$, and $F=3/2$. Assign $a$ to the left. No edge has both endpoints assigned yet, so $F$ stays $3/2$.

Vertex $b$ has one assigned neighbor on the left and none on the right: $\ell=1,r=0$. Put $b$ on the right. Edge $ab$ is now a crossing edge, while $ac$ and $bc$ remain unresolved. Thus $C=1,U=2$, and $F=2$. The increase is $(\ell-r)/2=1/2$, as predicted.

Vertex $c$ has one assigned neighbor on each side, so either choice changes $F$ by zero. Put it on the left. Now $ab$ and $bc$ cross, while $ac$ does not. We have $C=2,U=0,F=2$. The final value is an actual cut size, not an expectation about unfinished choices.

This trace distinguishes three quantities. $C$ counts crossings already determined. $U$ counts edges still awaiting at least one label. $F$ estimates the final count under stipulated fair choices for the unassigned labels. The algorithm need not actually toss coins. At each step it chooses a side whose computed $F$ is at least the current value. When nothing remains unassigned, the maintained estimate becomes the actual answer.

\section{A bound on rare large values}
For a nonnegative random variable $X$ and a real threshold $a>0$, we have
\[
\Prob(X\ge a)\le\frac{\E X}{a}.
\]
To prove it, split the expectation into outcomes where $X\ge a$ and the rest. The first part is at least $a\Prob(X\ge a)$, while the second is nonnegative. Hence $\E X\ge a\Prob(X\ge a)$, and division by $a$ gives the claim. This is Markov's inequality in a finite setting.

The nonnegativity assumption is not decorative. Large negative values could offset large positive ones in the expectation. For example, a variable equal to $100$ or $-100$ with equal probability has expectation zero but is at least $100$ with probability $1/2$.

\begin{principle}{Habit 8: choose a quantity whose average is easier than its best value}
When direct construction is difficult, define a finite experiment and count the desired or undesired features by indicators. Ask which averages can be computed exactly. Then state the logical conclusion carefully: one suitable outcome, a probability bound, or an actual efficient construction are different results.
\end{principle}

\begin{lab}{A proof that asks for an unnecessary assumption}
\noindent\textbf{Deliberately flawed criticism.} ``The cut proof is invalid because the edge indicators are not independent.''

\noindent\textbf{Learner.} ``Which identity requires independence? Expand the expectation of the sum directly. Separately identify where independence of vertex labels was used.''

\noindent\textbf{Model distinction.} Independence of the two endpoint labels gives each edge crossing probability $1/2$. Linearity of expectation then sums these probabilities without requiring independence between the edge events. Write a triangle example to demonstrate the difference. Finally, run the deterministic construction by hand on a four-cycle.
\end{lab}

\section*{Exercises}
\begin{exercise}\label{ex:13.1}A variable takes values $0,2,5$ with probabilities $1/2,1/3,1/6$. Compute its expectation and identify a value at least that large.\end{exercise}
\begin{exercise}\label{ex:13.2}For a triangle, list the possible cut sizes and compute the average over all eight vertex labelings.\end{exercise}
\begin{exercise}\label{ex:13.3}Run the deterministic cut construction in vertex order $1,2,3,4$ on the cycle with edges $12,23,34,41$. Break ties by choosing the left side.\end{exercise}
\begin{exercise}\label{ex:13.4}Four bad events each have probability at most $0.2$. Prove some outcome avoids all four. State what can be concluded if each bound is instead $0.25$.\end{exercise}
\begin{exercise}\label{ex:13.5}Give an example showing that a large expected value does not imply every outcome is large. Then state the existence conclusion that does follow.\end{exercise}
\begin{exercise}\label{ex:13.6}For a nonnegative random variable with expectation $3$, bound $\Prob(X\ge12)$. Give a finite example that attains the bound.\end{exercise}


% ===== Week 14 =====

\chapter{Sign sequences: from experiments to a research boundary}\label{ch:signs}
\goal{Combine symmetry, sums, linear algebra, and exact computation in one investigation. Derive a necessary condition rather than infer a classification from a finite search. Read a research statement whose objects are now familiar.}

\section{A sequence with small echoes}
A sign sequence of positive length $n$ is a list $a=(a_0,\ldots,a_{n-1})$ with each entry equal to $1$ or $-1$. Its aperiodic autocorrelation at a shift $s$ with $1\le s<n$ is
\[
A_s(a)=\sum_{j=0}^{n-1-s}a_j a_{j+s}.
\]
Slide the sequence by $s$ positions, multiply overlapping entries, and add the products. ``Aperiodic'' means entries beyond the endpoints are not wrapped around. At shift zero, the analogous sum is $n$, because every sign has square one.

For $a=(1,1,1,-1)$,
\begin{align*}
A_1&=(1)(1)+(1)(1)+(1)(-1)=1,\\
A_2&=(1)(1)+(1)(-1)=0,\\
A_3&=(1)(-1)=-1.
\end{align*}
A Barker sequence is a sign sequence with $|A_s|\le1$ at every nonzero aperiodic shift. Thus this example qualifies. The definition makes sense at length one as well: there are no nonzero shifts to check, so the condition holds vacuously.

Every $A_s$ is a sum of $n-s$ signs and therefore has the same parity as $n-s$. To see this, if $r$ summands are negative, the sum is $(n-s-r)-r=n-s-2r$. Hence a Barker sequence must have $A_s=0$ when $n-s$ is even, and $A_s=1$ or $-1$ when $n-s$ is odd. The definition has immediately produced an arithmetic restriction.

\subsection*{Decode the summation limits before computing}
At shift $s$, a term $a_j a_{j+s}$ is available only while both indices belong to $0,1,\ldots,n-1$. Starting at $j=0$, the last allowed $j$ satisfies $j+s=n-1$, or $j=n-1-s$. There are therefore $n-s$ terms, including both endpoints of that index range.

For a length-five sequence and shift two, the sum is
\[
 A_2=a_0a_2+a_1a_3+a_2a_4.
\]
There is no term $a_3a_5$ in an aperiodic correlation, because $a_5$ is outside the list. At the largest allowed shift $s=n-1$, only $a_0a_{n-1}$ remains. Checking these endpoint cases is an effective way to detect an incorrect summation limit.

The subscript on $A_s$ records the shift, whereas subscripts on $a_j$ record positions. The uppercase quantity is a sum built from the lowercase entries. Keeping these roles distinct helps when several indices appear in one calculation.

\section{Periodic correlation is a different quantity}
For periodic correlation, wrap the indices around modulo $n$:
\[
C_s(a)=\sum_{j=0}^{n-1}a_j a_{(j+s)\bmod n},\qquad0\le s<n.
\]
There are always $n$ summands. For $a=(1,1,1,-1)$, direct multiplication gives $C_0=4$ and $C_1=C_2=C_3=0$.

For $1\le s<n$, the split occurs precisely at $j=n-s$. When $0\le j\le n-1-s$, the index $j+s$ is still below $n$, and the contribution is $A_s$. When $n-s\le j\le n-1$, the index wraps once and becomes $j+s-n$. Thus
\[
 C_s=\sum_{j=0}^{n-1-s}a_j a_{j+s}
     +\sum_{j=n-s}^{n-1}a_j a_{j+s-n}.
\]
In the second sum, set $r=j+s-n$. At its first term $j=n-s$, we have $r=0$; at its last term $j=n-1$, we have $r=s-1$. Solving the same equation for $j$ gives $j=r+n-s$. Substitution therefore yields
\begin{align*}
 \sum_{j=n-s}^{n-1}a_j a_{j+s-n}
 &=\sum_{r=0}^{s-1}a_{r+n-s}a_r\\
 &=\sum_{r=0}^{s-1}a_r a_{r+n-s}\\
 &=A_{n-s}.
\end{align*}
The middle equality commutes the two real factors. The final equality uses the definition of $A_{n-s}$: its last index is $n-1-(n-s)=s-1$. This is a change of indexing, not an approximation.

For a concrete check, take $n=5,s=2$. The periodic sum has five terms:
\[
 C_2=\underbrace{a_0a_2+a_1a_3+a_2a_4}_{A_2}
       +\underbrace{a_3a_0+a_4a_1}_{A_3}.
\]
The last two become $a_0a_3+a_1a_4$ on reversing the factors. The example displays exactly the general split. Therefore
\[
C_s=A_s+A_{n-s}.
\tag{14.1}\label{eq:periodaperiod}
\]
The two definitions are related but not interchangeable. For $a=(1,1,-1)$, we have $A_1=0$, $A_2=-1$, so it is Barker; yet $C_1=-1$, not zero. A claim about one kind of correlation cannot be transferred to the other without a justified bridge.

\section{Circulant matrices package the same products}
Form an $n$-by-$n$ matrix by cyclically shifting the sequence in each row:
\[
H_{ij}=a_{(j-i)\bmod n},\qquad0\le i,j<n.
\]
Such a matrix is circulant. For our length-four example,
\[
H=\begin{pmatrix}
1&1&1&-1\\
-1&1&1&1\\
1&-1&1&1\\
1&1&-1&1
\end{pmatrix}.
\]
Every row has squared length $n$ because its $n$ entries have square one. To compare rows $i$ and $\ell$, write their dot product before changing notation:
\[
 \sum_{j=0}^{n-1}a_{(j-i)\bmod n}\,a_{(j-\ell)\bmod n}.
\]
Set $k=(j-i)\bmod n$. As $j$ visits all $n$ positions, $k$ does too: subtracting a fixed shift just rotates the index list. Modulo $n$, the second index is
\[
 j-\ell=(j-i)+(i-\ell)=k+(i-\ell).
\]
Hence the dot product becomes
\[
 \sum_{k=0}^{n-1}a_k a_{(k+i-\ell)\bmod n}=C_{(i-\ell)\bmod n}.
\]
When $i\ne\ell$ and both row indices lie between zero and $n-1$, their difference is not zero modulo $n$. Distinct rows therefore correspond to nonzero shifts. Conversely, for any nonzero shift $s$, the pair of rows $i=s,\ell=0$ produces that correlation. Both directions are needed to conclude that different rows are orthogonal \emph{exactly when} all nonzero periodic correlations vanish.

A real sign Hadamard matrix is a square matrix with entries $\pm1$ and mutually orthogonal rows. In the matrix language of Chapter~10, the condition is
\[
HH^{\mathsf T}=nI.
\]
A circulant sign Hadamard matrix therefore exists exactly when a sign sequence has $C_s=0$ for all nonzero periodic shifts. We have changed the language of the same question: sequence correlations became row orthogonality.

\subsection*{Read the matrix identity entry by entry}
The identity $HH^{\mathsf T}=nI$ contains $n^2$ scalar equalities. At entry $(i,\ell)$, the left side is the dot product of row $i$ with row $\ell$. The right side equals $n$ if $i=\ell$ and zero otherwise. Thus the diagonal equations record row lengths, and the off-diagonal equations record pairwise orthogonality.

For the displayed four-row matrix, the first and second rows have dot product $-1+1+1-1=0$. The first row dotted with itself gives four. The general reindexing argument handles all the remaining pairs together. A matrix formula is useful because it groups related equations, but understanding it requires knowing what one of its entries means.

\section{The sum of all correlations reveals a square}
Let $r=\sum_{j=0}^{n-1}a_j$ be the row sum. Add every periodic correlation, including shift zero:
\begin{align*}
\sum_{s=0}^{n-1}C_s
&=\sum_{j=0}^{n-1}a_j\left(\sum_{s=0}^{n-1}a_{(j+s)\bmod n}\right)\\
&=\sum_{j=0}^{n-1}a_j\left(\sum_{k=0}^{n-1}a_k\right)\\
&=r^2.
\end{align*}
In the second line, for each fixed $j$, the indices $(j+s)\bmod n$ run through every position exactly once. This is the specific symmetry that justifies the reindexing.

One way to see the whole count is to use three abstract entries $a_0,a_1,a_2$. The three periodic correlations list these products:
\begin{align*}
 C_0&=a_0a_0+a_1a_1+a_2a_2,\\
 C_1&=a_0a_1+a_1a_2+a_2a_0,\\
 C_2&=a_0a_2+a_1a_0+a_2a_1.
\end{align*}
Together they contain each ordered product $a_j a_k$ once. Expanding $(a_0+a_1+a_2)^2$ gives those same nine ordered products, so the totals agree. For general $n$, each fixed first index $j$ meets every second index exactly once as the shift varies. That is why the total is $(\sum_j a_j)^2$, rather than, for example, $\sum_j a_j^2$.

This identity itself does not require signs; the same finite rearrangement holds for real entries. The restriction to signs is used afterward to identify $C_0=n$ and to make the row sum an integer. Separating those roles tells us which part of the argument might generalize.

If all nonzero $C_s$ vanish, the left side is just $C_0=n$. Hence
\[
n=r^2.
\tag{14.2}\label{eq:squareorder}
\]
The integer $r$ is a sum of signs, so the order $n$ must be a square. This is a proved necessary condition, valid for every positive length, not a pattern extrapolated from an experiment.

\section{Orthogonality also forces divisibility by four}
\par\noindent\begin{minipage}{\linewidth}
\begin{lemma}\label{thm:hadamardorder}
If a real sign Hadamard matrix has order $n>2$, then $4$ divides $n$.
\end{lemma}
\end{minipage}\par

\begin{proof}
Multiply each column by the sign in its first row. This makes the first row all $1$'s and preserves every row dot product, because multiplying both factors in a column by the same sign multiplies their product by one.

The second row is orthogonal to the all-ones first row, so its signs sum to zero. It therefore contains equally many plus and minus signs, implying $n$ is even. Permute columns so the second row has its $n/2$ plus signs first and its $n/2$ minus signs last. Column permutations also preserve row dot products.

These two transformations simplify the matrix for the proof; they do not assert that its original first row was already all positive or that its original second row was already sorted. Multiplying a column by $-1$ changes that coordinate in every row, so each term in a row dot product gets two minus signs and is unchanged. Reordering columns applies the same permutation to both factors in every row comparison, merely reordering the sum.

We do not need these transformations to preserve the circulant pattern. The lemma applies to \emph{every} sign Hadamard matrix, so its proof needs only to preserve the sign entries and row orthogonality. Circulant structure is used separately in the square-order identity. Requiring every intermediate representation to retain more structure than the lemma uses would obscure the argument.

Since $n>2$, a third row exists. Let $x$ be the sum of its entries in the first half of the columns and $y$ the sum in the second half. Orthogonality with the first row gives $x+y=0$. Orthogonality with the second row gives $x-y=0$. Adding and subtracting yield $x=y=0$.

In detail, adding the equations gives $2x=0$, hence $x=0$. Substituting into $x+y=0$ gives $y=0$. Let $h=n/2$ be the number of entries in the first half, and suppose exactly $p$ of those entries are positive. The other $h-p$ are negative, so their sum is $p-(h-p)=2p-h$. Since this sum is $x=0$, we obtain $h=2p$, proving that the half-length is even. This explains the second divisibility-by-two step rather than merely repeating the whole-row argument.

The first half of the third row therefore consists of $n/2$ signs summing to zero. It has equally many plus and minus entries, so $n/2$ is even. Thus $n$ is divisible by four.
\end{proof}

\par\noindent\begin{minipage}{\linewidth}
\begin{proposition}\label{thm:circulantnecessary}
A real circulant sign Hadamard matrix of positive order has $n=1$ or $n=4u^2$ for some positive integer $u$.
\end{proposition}
\end{minipage}\par

\begin{proof}
Equation~\eqref{eq:squareorder} gives $n=r^2$ for an integer $r$. If $n>1$, it is not $2$, because $2$ is not an integer square. Hence $n>2$, and the lemma gives $4\mid n$. The integer $r$ must then be even: an odd square is odd. Write $|r|=2u$ with $u>0$. It follows that $n=4u^2$.
\end{proof}

This leaves candidate orders $4,16,36,\ldots$. It does not construct matrices at those orders and does not rule them out. Necessary conditions define a research boundary: the argument has eliminated many cases and has stated exactly what remains.

\section{An exact finite experiment}
There are $2^n$ labeled sign sequences of length $n$. The companion program enumerates all of them for $1\le n\le13$, using integer arithmetic. It counts Barker sequences and sequences with zero nonzero periodic correlations. No symmetry quotient is taken, so rotations, reversals, and global sign changes are counted separately whenever they produce different lists.

\begin{center}
\begin{tabular}{@{}r rr@{}}\toprule
Length & Barker sequences & Zero nonzero periodic correlations\\\midrule
1&2&2\\2&4&0\\3&4&0\\4&8&8\\5&4&0\\6&0&0\\7&4&0\\8&0&0\\9&0&0\\10&0&0\\11&4&0\\12&0&0\\13&4&0\\\bottomrule
\end{tabular}
\end{center}
The program's enumeration was run for this edition. The table is evidence about exactly this finite range. It cannot rule out a sequence at an untested larger length. The number of possible cases doubles at each increase of length, which also explains why an unqualified instruction to ``test many more cases'' is not a mathematical strategy by itself.

Global sign negation preserves every pair product, as established in Chapter~6. We could restrict the first sign to $1$ and double the counts afterward, because for a nonempty sequence the negated sequence is distinct and has the opposite first sign. The companion's primary enumeration intentionally does not rely on this reduction; it uses the symmetry as a separate cross-check.

\section{What the repository statement adds}
The scope note for Family~179 documents a formalized classification of positive circulant sign Hadamard orders as exactly $1$ and $4$. It separately documents the even-length Barker consequence as lengths $2$ or $4$, and states that the paper's odd-length classification is outside that selected additional formal statement~\cite{oai179}. These are claims and scope descriptions from the repository, not results established by our finite experiment or by Proposition~\ref{thm:circulantnecessary}.

At this point you can understand the objects and the exact strength of the circulant statement. You also know a complete elementary supporting argument and its limit: excluding $4u^2$ for $u>1$ requires additional mathematics. Chapter~17 shows how to request that additional background without pretending that a short elementary argument has already supplied it.

\begin{lab}{Preventing an experiment from becoming a false proof}
\noindent\textbf{Deliberately flawed conclusion.} ``We found circulant examples only at lengths $1$ and $4$, so those are the only possible lengths.''

\noindent\textbf{Learner.} ``Separate the finite search statement, the necessary-order theorem, and the external classification. Which of these rules out order $16$? Which has actually been proved here?''

\noindent\textbf{Model audit.} The finite search stops at $13$, so says nothing about $16$. The necessary condition permits $16=4\cdot2^2$. The external classification excludes it, but its proof must be read or independently checked. The course has established neither that exclusion by enumeration nor the full classification by its elementary argument.
\end{lab}

\section*{Exercises}
\begin{exercise}\label{ex:14.1}Compute every nonzero aperiodic and periodic correlation of $(1,1,-1)$.\end{exercise}
\begin{exercise}\label{ex:14.2}Prove that negating every sign preserves both kinds of correlation.\end{exercise}
\begin{exercise}\label{ex:14.3}For $(1,1,1,-1)$, verify one off-diagonal entry and one diagonal entry of $HH^{\mathsf T}=4I$ directly.\end{exercise}
\begin{exercise}\label{ex:14.4}Prove $C_{n-s}=C_s$ for a real sign sequence by reindexing. State the allowed range of $s$.\end{exercise}
\begin{exercise}\label{ex:14.5}Use the two necessary conditions to decide which of $8,9,16,25,36$ remain candidates for a circulant sign Hadamard order. Do not call a remaining candidate an existing example.\end{exercise}
\begin{exercise}\label{ex:14.6}For a finite search of all length-$n$ sign sequences, explain why a single example certifies existence while zero examples certifies nonexistence only at that length. State what a proof of nonexistence for every length would additionally require.\end{exercise}


% ===== Week 15 =====

\chapter{Lean I: reading the state of a proof}\label{ch:lean1}
\goal{Connect the proof states you already understand to a formal language. Read a theorem's signature before its tactics. Write small proofs whose mathematical meaning remains visible.}

\section{A checker is not the same as an explanation}
Lean is a programming language and proof assistant. A proposition is represented by a type of proofs, and a completed proof is checked against its stated type. Tactics are instructions that construct proof terms; they are not substitutes for the underlying statement. The official introductions explain these ideas in more detail~\cite{leanprop,leantactics}.

Our purpose is not to memorize a catalogue of tactics. It is to make the discipline of Chapter~4 explicit: which objects and hypotheses are available, and what exactly remains to be established? The examples below reproduce mathematics already learned, so that a syntax difficulty can be distinguished from a conceptual difficulty.

The official installation route uses VS Code and the Lean 4 extension, whose setup guide installs the needed toolchain~\cite{leaninstall}. The companion files use only Lean's core and standard library, not Mathlib. The file \code{lean-toolchain} requests version 4.19.0 as a fixed target, not as a claim about the latest release. The reproducibility note records which execution checks were and were not possible for this edition.

\subsection*{From a mathematical sentence to a file}
A Lean file is an ordinary text file whose name ends in \code{.lean}. It contains declarations: some define objects or meanings, and others state and prove results. In the companion, \code{import Std} makes the standard library available. The line \code{namespace MathCourse} groups our names so they do not collide with names from another file. A closing \code{end MathCourse} ends that group; it does not end a mathematical proof.

Open \code{proofs.lean} as a file, not as text to paste into the terminal. The terminal command \code{lean proofs.lean} asks the Lean executable to check that file. The name after \code{lean} is the filename. A terminal is a text window for issuing commands to installed programs; it does not interpret Lean proof lines as shell commands. The practical appendix explains how to run the separate companions.

You can read the mathematics in this chapter before installing anything. Later, the editor's displayed proof state lets you compare the stated hypotheses and goal after each line with the explanations here. These teaching descriptions are not a claim that this edition's files were compiler-checked; the execution status is stated in Appendix~\ref{app:practical}.

\section{A few symbols with exact roles}
In Lean, \code{x : Nat} says $x$ is a natural number. The name \code{Nat} corresponds to $\NN$ in this book. The expression \code{P : Prop} says $P$ is a proposition. A hypothesis \code{h : P} is evidence for $P$, not a Boolean flag merely reporting that a test was true.

Function application is written with a space: \code{f x} means $f(x)$. The function \code{fun x => f x} sends an input $x$ to $f(x)$. The colon introduces a type or a theorem statement. The symbol \code{:=} supplies a definition or proof. Ordinary \code{=} inside a statement still means equality.

It helps to read a declaration by its parts before reading an entire proof. In
\par\noindent\begin{minipage}{\linewidth}
\begin{lstlisting}
 def twice (n : Nat) : Nat := n + n
\end{lstlisting}
\end{minipage}\par

\code{def} announces a definition, \code{twice} is its chosen name, \code{(n : Nat)} declares the input, and the final \code{: Nat} declares the output's type. The text after \code{:=} gives the rule. A \emph{type} specifies what kind of object an expression is: here both the input and output are natural numbers.

A theorem declaration uses the same basic separation. Its name and parameters come first, its claimed proposition follows a colon, and a proof follows \code{:=}. Writing \code{by} begins a sequence of proof-building instructions. The command does not mean ``assume the statement is true.'' It begins the work required to justify it.

The theorem below has no input parameter. It claims one concrete equality about the definition just supplied.

For example:
\par\noindent\begin{minipage}{\linewidth}
\begin{lstlisting}
def twice (n : Nat) : Nat := n + n

theorem twice_two : twice 2 = 4 := by
  rfl
\end{lstlisting}
\end{minipage}\par

The first declaration defines a function. The second claims an equality and gives a proof. The tactic \code{rfl} closes an equality whose two sides reduce to the same expression by definitions and computation. Here \code{twice 2} unfolds to \code{2 + 2}, which computes to \code{4}. It would not prove an arbitrary true equality that requires an independent theorem.

The theorem's equality and the definition's assignment have different jobs. The definition supplies what \code{twice} means. The theorem says that evaluating that meaning at input two gives four. A different name would not change the mathematics. A different definition body could change the claim's truth.

The word \emph{reduce} means replace a defined expression by its body and perform the permitted computation. In this example, reduction is the sequence \code{twice 2}, then \code{2 + 2}, then \code{4}. When both sides reduce to the same expression, \code{rfl} can establish equality. For a general equality whose two sides do not reduce to the same expression, a mathematical lemma may still be needed even if numerical instances are easy.

Names are not mathematical content. A theorem called \code{all\_problems\_solved} may have a trivial statement. Always read the expression after its parameters and colon.

\subsection*{A proposition and a proof of it are different objects}
In ordinary writing, $P$ might be the sentence ``this integer is even.'' A proof of $P$ is an argument establishing that sentence, not another name for the integer. Lean writes \code{P : Prop} to declare the proposition and \code{hp : P} to declare evidence for it.

An implication $P\Rightarrow Q$ can be used as a rule: give it a proof of $P$, and it supplies a proof of $Q$. Lean expresses this as \code{P -> Q}. If \code{h : P -> Q} and \code{hp : P}, the expression \code{h hp} applies the rule and has type \code{Q}. This is function application, not multiplication. The familiar inference ``$P$; $P$ implies $Q$; therefore $Q$'' has been made explicit.

This is the part of the propositions-as-types viewpoint we need first. It does not require understanding the implementation of a proof assistant before reading the example. We learn how the notation records an inference already understood in ordinary mathematics.

\section{An implication as a change of goal}
Suppose $P,Q,R$ are propositions, $h_{PQ}$ proves $P\Rightarrow Q$, and $h_{QR}$ proves $Q\Rightarrow R$. We want $P\Rightarrow R$.
\par\noindent\begin{minipage}{\linewidth}
\begin{lstlisting}
theorem implication_chain (P Q R : Prop)
    (hpq : P -> Q) (hqr : Q -> R) : P -> R := by
  intro hp
  have hq : Q := hpq hp
  exact hqr hq
\end{lstlisting}
\end{minipage}\par

Read the proof line by line. Initially the goal is \code{P -> R}. The command \code{intro hp} says: assume $P$, call its proof \code{hp}, and now prove $R$. The command \code{have hq : Q := hpq hp} applies the implication $P\Rightarrow Q$ to the available proof of $P$, producing a proof of $Q$. The last line applies $Q\Rightarrow R$ to that proof and supplies exactly what the goal needs.

The corresponding states are:
\begin{center}
\begin{tabularx}{\textwidth}{@{}p{1.3in}XX@{}}\toprule
After & Available new fact & Remaining goal\\\midrule
Before the proof & $P\Rightarrow Q$, $Q\Rightarrow R$ & $P\Rightarrow R$\\
\code{intro hp} & $P$ & $R$\\
\code{have hq} & $Q$ & $R$\\
\code{exact hqr hq} & a proof of $R$ & no remaining goal\\\bottomrule
\end{tabularx}
\end{center}
These are instructional descriptions of the states, not a saved compiler transcript. In an editor, inspect Lean's actual state after each line and compare it with this reasoning.

\section{Conjunction means two obligations}
To prove $Q\land P$ from $P\land Q$, extract the two components and put them in the opposite order:
\par\noindent\begin{minipage}{\linewidth}
\begin{lstlisting}
theorem and_swap (P Q : Prop) : And P Q -> And Q P := by
  intro h
  exact And.intro h.right h.left
\end{lstlisting}
\end{minipage}\par

Here \code{And P Q} means the conjunction $P\land Q$. The expressions \code{h.right} and \code{h.left} extract its second and first components. The constructor \code{And.intro} packages two proofs into a conjunction; its argument order must match the desired conclusion. Thus the proof of $Q$ comes first and the proof of $P$ second.

This is the formal counterpart of proving two set containments or two directions of an equivalence. Do not let one completed branch obscure an unfinished second obligation.

\subsection*{Parameters, arrows, and implicit inputs}
The type \code{A -> B} describes functions taking an input of type \code{A} and returning an output of type \code{B}. The type \code{A -> B -> C} groups as \code{A -> (B -> C)}: after receiving an \code{A}, the function receives a \code{B}, then returns a \code{C}. This lets Lean write multiple inputs by spaces, as in \code{f x y}. Application groups to the left: first \code{f x}, then its result applied to \code{y}.

For a hypothesis saying injectivity holds, the inputs will be two objects and evidence that their outputs agree. Its result will be evidence that the objects agree. Counting these inputs will make the application in the next proof readable.

The word \code{Type} in the next declaration indicates a type of objects rather than a proposition to prove. The notation \code{\{A B : Type\}} introduces types as implicit parameters: Lean normally infers them from the supplied function. Parentheses around a parameter instead request that it be supplied explicitly. The mathematical domains are still present when braces are used; they have not been removed from the statement.

Finally, \code{fun x => g (f x)} is the rule sending $x$ to $g(f(x))$. It names a function without first giving it a separate declaration. In ordinary mathematical notation it is $x\mapsto g(f(x))$.

\section{The function-composition proof returns}
We define injectivity explicitly, rather than assume knowledge of a library name:
\par\noindent\begin{minipage}{\linewidth}
\begin{lstlisting}
def Inj {A B : Type} (f : A -> B) : Prop :=
  forall x y : A, f x = f y -> x = y

theorem inj_comp {A B C : Type}
    (f : A -> B) (g : B -> C)
    (hf : Inj f) (hg : Inj g) :
    Inj (fun x => g (f x)) := by
  intro x y h
  have hxy : f x = f y := hg (f x) (f y) h
  exact hf x y hxy
\end{lstlisting}
\end{minipage}\par

The word \code{forall} is ASCII notation for the universal quantifier. \code{Type} denotes a type of mathematical objects. Curly braces mark parameters Lean can usually infer from other inputs: if it knows the type of $f$, it can infer its input and output types. The definition \code{Inj} says exactly what the earlier mathematical definition said.

Read its body in order. \code{forall x y : A} means take any two inputs of the domain. The remaining expression \code{f x = f y -> x = y} says that equality of their outputs implies equality of those inputs. Thus the declared hypothesis \code{hg : Inj g} can be applied to two elements of \code{B} and a proof that their \code{g} outputs agree.

In the proof, those two elements are \code{f x} and \code{f y}, not \code{x} and \code{y}. The evidence \code{h} concerns their final outputs. This accounts for all three arguments in \code{hg (f x) (f y) h}. Its result is the intermediate equality that \code{hf} can then use.

The first proof line introduces arbitrary inputs and an equality of final outputs. The expression \code{hg (f x) (f y) h} supplies two inputs of $g$ and the evidence that their outputs agree. It returns equality of the intermediate values. The final application of \code{hf} then returns equality of the original inputs. These are the same two steps as Theorem~\ref{thm:injcomp}.

If Lean rejects the application, inspect the actual types of \code{h}, \code{hf}, and \code{hg}. A common mistake is to apply \code{hf} directly to equality after $g$, skipping the intermediate equality that injectivity of $g$ must establish.

\section{Existence requires a witness and evidence}
To prove $\exists n\in\NN:n+1=4$, supply the candidate $3$ and evidence that it satisfies the property:
\par\noindent\begin{minipage}{\linewidth}
\begin{lstlisting}
theorem a_witness : Exists (fun n : Nat => n + 1 = 4) := by
  exact Exists.intro 3 rfl
\end{lstlisting}
\end{minipage}\par

The expression \code{Exists (fun n : Nat => n + 1 = 4)} states that some natural number has the displayed property. This is the typed form of the existential quantifier. \code{Exists.intro} packages a witness with its proof. The equality is computational in this case, so \code{rfl} suffices. For a more difficult witness, the second argument would be a substantive proof.

Conversely, using an existential hypothesis requires unpacking both components. Knowing merely that ``something exists'' is not enough to apply a formula to it without introducing a name and its guaranteed property. Chapter~16 uses this operation in a complete surjectivity proof.

\subsection*{The type of a number changes available operations}
A natural number has no negative values. Lean's natural-number subtraction is truncated at zero: \code{(7 - 10 : Nat)} is zero. Integer subtraction is different: \code{(7 - 10 : Int)} is minus three. The type annotation inside the parentheses specifies which operation is meant.

In ordinary mathematics we also need this distinction, although informal notation sometimes leaves the domain implicit. An argument that subtracts a larger integer from a smaller one and then treats the result as a negative number cannot be translated unchanged into natural-number subtraction. Before asking an assistant to repair a formal proof, check whether it has silently changed the kind of number being used.

\section{Equality can be transported through a map}
If $h$ proves $x=y$, then \code{congrArg f h} proves $f(x)=f(y)$. This is not injectivity: it applies the same function to equal inputs. Injectivity goes in the opposite direction and needs an extra hypothesis.

The command \code{rw [h]} rewrites an occurrence of one side of an equality using the other. A \code{calc} block organizes a chain of equalities, with a justification for each link~\cite{leanquant}. For instance:
\par\noindent\begin{minipage}{\linewidth}
\begin{lstlisting}
theorem left_inverse_inj {A B : Type}
    (f : A -> B) (g : B -> A)
    (h : forall x, g (f x) = x) : Inj f := by
  intro x y hxy
  calc
    x = g (f x) := (h x).symm
    _ = g (f y) := congrArg g hxy
    _ = y := h y
\end{lstlisting}
\end{minipage}\par

The first step reverses the identity $g(f(x))=x$; \code{.symm} reverses an equality proof. The middle step applies $g$ to the assumed equality $f(x)=f(y)$. The final step uses the inverse identity at $y$. The underscores reuse the previous expression, keeping the chain readable.

\begin{lab}{One repair at a time}
Give the assistant the composition theorem with the middle \code{have} line removed. Ask for the missing mathematical fact first, not a replacement script. Then request exactly one Lean line expressing that fact. After inserting it, inspect the remaining goal yourself. Finally change the theorem to the false converse from Chapter~3. Before requesting tactics, supply the finite counterexample and explain why no sound proof of that unrestricted statement should exist.
\end{lab}

\section*{Exercises}
\begin{exercise}\label{ex:15.1}In \code{h : P -> Q}, distinguish the role of \code{h}, \code{P}, and \code{Q}. What does \code{h hp} provide when \code{hp : P}?\end{exercise}
\begin{exercise}\label{ex:15.2}Write a Lean proof of $P\land Q\Rightarrow P$. Explain its only mathematical inference.\end{exercise}
\begin{exercise}\label{ex:15.3}Write a Lean theorem giving a natural-number witness for $n+2=7$.\end{exercise}
\begin{exercise}\label{ex:15.4}Explain why \code{congrArg g h} cannot in general recover $x=y$ from \code{h : g x = g y}. Which hypothesis is missing?\end{exercise}
\begin{exercise}\label{ex:15.5}Using \code{Inj}, formalize the claim that if $g\circ f$ is injective, then $f$ is injective. Do not assume $g$ itself is injective.\end{exercise}
\begin{exercise}\label{ex:15.6}Explain why \code{(7 - 10 : Nat)} and \code{(7 - 10 : Int)} describe different computations. What lesson does this give for translating a theorem's domain?\end{exercise}


% ===== Week 16 =====

\chapter{Lean II: definitions, dependencies, and trust}\label{ch:lean2}
\goal{Use formalization to expose hidden choices in a statement. Read the hypothesis list and definition bodies, not only the accepted proof. Practice a complete debugging cycle that preserves the intended claim.}

\section{The statement must be right before the proof can help}
A formal proof establishes the encoded proposition from its assumptions. It does not, by itself, establish that this proposition matches a sentence in a manuscript, a user's intention, or a physical model. This distinction is part of formal-proof validation, not a criticism of the proof assistant~\cite{leanvalidate}.

Consider the difference between ``an endpoint of a two-edge walk'' and ``a vertex at shortest directed distance exactly two.'' Both can be defined precisely and used in accepted proofs. Only one matches the second-neighborhood convention in Chapter~7. Formalization makes the distinction explicit, but a reader must still inspect it.

Before looking at tactics, make a correspondence table: each mathematical object and assumption on one side, its formal representation on the other. Check domains, strict inequalities, nonemptiness, and excluded cases. Then check that the theorem's conclusion has the intended quantifier order.

\subsection*{What to inspect when the result is unfamiliar}
First locate the entire theorem declaration. Its parameters contain objects and hypotheses; its final proposition contains the conclusion. Then follow definitions appearing in that proposition. A definition may contain conditions that a short English name does not reveal. Only after the statement is understood should you inspect how the proof constructs its evidence.

For example, the name \code{Second} below cannot tell us whether first neighbors are excluded. The definition body can. Similarly, the name \code{Surj} cannot by itself tell us its domain, codomain, or order of witnesses. A two-column translation makes these visible before any tactics are used.

\section{A complete existence proof in Lean}
Recall that surjectivity means every target has a preimage. Define it explicitly:
\par\noindent\begin{minipage}{\linewidth}
\begin{lstlisting}
def Surj {A B : Type} (f : A -> B) : Prop :=
  forall y : B, Exists (fun x : A => f x = y)

theorem surj_comp {A B C : Type}
    (f : A -> B) (g : B -> C)
    (hf : Surj f) (hg : Surj g) :
    Surj (fun x => g (f x)) := by
  intro c
  cases hg c with
  | intro b hb =>
    cases hf b with
    | intro a ha =>
      exact Exists.intro a (by rw [ha, hb])
\end{lstlisting}
\end{minipage}\par

Start with an arbitrary target $c$.

At this point the goal is an existence statement: find some $a$ in the input type $A$ with $g(f(a))=c$. We do not yet know which $a$ to use. First use the hypothesis about $g$, because it is the map whose output has type $C$ and can reach our target $c$.

The proof of existence supplied by \code{hg c} contains a witness and a property. The lines
\par\noindent\begin{minipage}{\linewidth}
\begin{lstlisting}
 cases hg c with
 | intro b hb =>
\end{lstlisting}
\end{minipage}\par

mean: unpack that existence evidence, name the witness \code{b}, name its property \code{hb}, and continue inside the indented branch. Here \code{intro} following the vertical bar names the constructor case of an existential proof; it is not the earlier standalone tactic for assuming a premise. Both uses are related to introducing information, but their positions in the syntax differ.

The second unpacking does the same with \code{hf b}. Inside its branch, the available information is
\[
 b\in B,\quad g(b)=c,\qquad a\in A,\quad f(a)=b.
\]
The target still asks for an input to the composition. We now have the candidate $a$. It remains to verify its property, which follows by substituting $f(a)=b$ and then $g(b)=c$. The command \code{rw [ha, hb]} performs these two substitutions in that order, establishing $g(f(a))=c$.

The witnesses are chosen in reverse order from function execution: first find an input to $g$ that reaches the final target, then find an input to $f$ that reaches that intermediate value. This is the same dependency logic from Chapter~2, now visible in a formal script.

\subsection*{An ordinary proof beside the syntax}
The surjectivity argument can be written without Lean as a three-line dependency chain:
\[
 c\ \longleftarrow\ g(b)=c,\qquad
 b\ \longleftarrow\ f(a)=b,\qquad
 g(f(a))=g(b)=c.
\]
The first two statements use hypotheses to obtain witnesses. The last checks the proposed final witness. Lean's \code{Exists.intro a} packages that witness with its evidence. The small nested \code{by} expression supplies only the equality proof needed for that package; it does not start a new theorem with different assumptions.

Indentation records which named witnesses are in scope. A witness unpacked inside a branch is available in that branch. The proof does not get to use an unnamed existential object as though its value had already been chosen. This is the formal version of the distinction between ``there is a solution'' and ``let this named object be a solution'' from Chapter~4.

\section{Sets as membership predicates}
A subset of a type $A$ can be represented by a predicate \code{s : A -> Prop}, where \code{s x} means $x$ belongs to it.

A \emph{predicate} is a rule producing a proposition from an input. For example, ``$n$ is even'' is a predicate on integers: for each particular integer it gives a statement that can be proved or refuted. The predicate describes the set of inputs for which it holds.

There is no contradiction between representing a set as a collection of members in Chapter~3 and as a membership predicate here. The question ``does $x$ belong?'' connects the two descriptions. A predicate is not necessarily a Boolean program that computes an answer immediately. A proposition may require substantial reasoning to decide or prove. In these examples we use predicates to express meanings, not to promise a fast membership algorithm. Define intersection and preimage using their meanings:
\par\noindent\begin{minipage}{\linewidth}
\begin{lstlisting}
def meet {A : Type} (s t : A -> Prop) : A -> Prop :=
  fun x => And (s x) (t x)

def preimage {A B : Type} (f : A -> B)
    (s : B -> Prop) : A -> Prop :=
  fun x => s (f x)

theorem preimage_meet {A B : Type} (f : A -> B)
    (s t : B -> Prop) (x : A) :
    Iff (preimage f (meet s t) x)
      (meet (preimage f s) (preimage f t) x) := by
  rfl
\end{lstlisting}
\end{minipage}\par

The expression \code{Iff P Q} means logical equivalence of $P$ and $Q$. Unfolding the left side gives \code{And (s (f x)) (t (f x))}. Unfolding the right side gives the very same expression. The theorem is short because the definitions already encode the reasoning. It is not short because an opaque automation step skipped an unexplained mathematical argument.

This example also shows why definitions deserve close attention. Replacing \code{meet} by a differently defined operation could preserve the theorem name but change the statement. Formal names are handles for expressions, not warranties about their intended meaning.

\section{Encoding the exclusion in a second neighborhood}
We can represent an oriented edge relation by \code{r : V -> V -> Prop}; the proposition \code{r v w} means $v\to w$. The local membership condition for a genuine second neighbor can be written
\par\noindent\begin{minipage}{\linewidth}
\begin{lstlisting}
def Second {V : Type} (r : V -> V -> Prop)
    (v w : V) : Prop :=
  And (Not (w = v)) (And (Not (r v w))
    (Exists (fun u : V => And (r v u) (r u w))))
\end{lstlisting}
\end{minipage}\par

The three components say the endpoint is not the initial vertex, it is not already a first neighbor, and a two-edge route exists. Here \code{Not P} means logical negation, so \code{Not (w = v)} expresses $w\ne v$. The explicit \code{And} and \code{Exists} names make the grouping visible and avoid confusing propositional inequality with a Boolean comparison.

To read the nesting, temporarily give the three conditions names:
\[
 P:\ w\ne v,\qquad Q:\ \text{not }v\to w,\qquad
 R:\ \text{some }u\text{ gives }v\to u\to w.
\]
Then the definition has the shape \code{And P (And Q R)}. A proof of it contains a proof of $P$ and, as its second component, a pair containing proofs of $Q$ and $R$. Accordingly, \code{h.left} returns $P$, \code{h.right.left} returns $Q$, and \code{h.right.right} returns $R$.

The existential statement inside $R$ does not require the intermediate vertex to be unique. It supplies at least one valid two-edge route. Multiple choices of $u$ still describe membership of the \emph{same endpoint} $w$ in a neighborhood, rather than several copies of that endpoint.

A small but useful consequence is that a second neighbor is not a first neighbor:
\par\noindent\begin{minipage}{\linewidth}
\begin{lstlisting}
theorem second_not_first {V : Type}
    (r : V -> V -> Prop) (v w : V)
    (h : Second r v w) : Not (r v w) := by
  exact h.right.left
\end{lstlisting}
\end{minipage}\par

The nested conjunction is grouped to the right. Thus \code{h.right} contains the last two conditions, and its \code{left} component is the required exclusion. Reading this projection correctly is more important than memorizing its name.

The full repository theorem also needs a finite vertex type, nonemptiness, and an oriented relation, and counts the relevant finite sets. Our local definition does not formalize that whole theorem. It isolates a piece of statement meaning that can be checked before studying a large proof.

\section{A successful proof can still be irrelevant}
The following theorem is logically valid:
\par\noindent\begin{minipage}{\linewidth}
\begin{lstlisting}
theorem impossible_context (h : False) : (0 : Nat) = 1 := by
  exact False.elim h
\end{lstlisting}
\end{minipage}\par

It does not establish that zero equals one unconditionally. It says that from the impossible assumption \code{False}, any conclusion follows. The assumption is part of the statement. A reader who looks only at the conclusion would misunderstand what was proved.

This matters when an assistant ``repairs'' a failing proof by adding a hypothesis. Adding injectivity may make a false cancellation theorem true. Adding existence of the desired solution may trivialize an existence theorem. The code may then check, but the mathematical task has changed. Review the statement difference before reviewing the proof difference.

\section{Placeholders and axiom dependencies}
Lean permits unfinished developments containing \code{sorry}, which uses a placeholder axiom. An accepted declaration that depends on such a placeholder is not a completed proof. The command \code{\#print axioms theoremName} displays axiom dependencies and is useful for detecting this situation. Standard foundational axioms and evaluation-related trust assumptions must also be interpreted rather than treated as interchangeable errors~\cite{leanaxioms}.

For the companion, inspect commands such as
\par\noindent\begin{minipage}{\linewidth}
\begin{lstlisting}
#print axioms MathCourse.inj_comp
#print axioms MathCourse.surj_comp
#print axioms MathCourse.preimage_meet
\end{lstlisting}
\end{minipage}\par

An axiom list is not a substitute for the theorem signature. The impossible-context theorem can have no additional axiom dependency while still assuming \code{False} as a parameter. Conversely, a theorem may use ordinary foundational principles without assuming its desired conclusion. Both the assumptions in the signature and the dependencies of the proof must be understood.

The command-line check \code{lean proofs.lean} should be run from the companion directory with its pinned toolchain. Check the process result, warnings, and requested dependency output. Do not interpret absence of visible editor errors before checking has finished as a completed verification. Larger or untrusted projects call for the stronger validation procedures described in the official documentation~\cite{leanvalidate}.

\section{A debugging sequence that does not lose the mathematics}
When a line fails, first freeze the statement. Read the goal and the types of the available facts. Translate both into ordinary mathematics. Decide whether the issue is a false claim, a missing lemma, a mismatch of domains, or merely an expression written in a different form. Only then ask for a tactic repair.

For a syntax issue, a one-line correction may be enough. For a missing theorem, request its exact statement and hypotheses. For a false target, stop trying to automate the proof and construct a counterexample. For a representation mismatch, justify the translation before rewriting. This sequence uses Lean as feedback without letting its error messages dictate an unprincipled search for any statement that will compile.

\begin{lab}{An assumption audit in two columns}
Take the impossible-context example and the injective-composition theorem. For each, make one column listing explicit parameters and another listing imported or foundational proof dependencies. Ask the assistant why an empty dependency list would not make the first theorem a proof of $0=1$. Then compare the \code{Second} definition with a definition containing only the existential two-step condition. Use the directed triangle from Chapter~7 to explain the semantic difference before discussing formal syntax.
\end{lab}

\section*{Exercises}
\begin{exercise}\label{ex:16.1}In the surjectivity proof, explain what \code{hb} and \code{ha} state and why the final witness belongs to the domain of $f$.\end{exercise}
\begin{exercise}\label{ex:16.2}Replace \code{meet} by a union predicate and prove the corresponding preimage-union membership equivalence.\end{exercise}
\begin{exercise}\label{ex:16.3}From \code{h : Second r v w}, write a proof that $w\ne v$ and a proof that a two-edge route exists.\end{exercise}
\begin{exercise}\label{ex:16.4}An assistant adds \code{h : False} to a theorem that it could not prove. Explain precisely what has changed and why accepting the new file does not answer the original question.\end{exercise}
\begin{exercise}\label{ex:16.5}Design a three-step review that detects a changed domain, a missing exclusion, and an unfinished proof placeholder. State which artifact you inspect at each step.\end{exercise}
\begin{exercise}\label{ex:16.6}Explain the difference between evaluating a finite mathematical test in Python, proving a finite proposition in Lean, and proving a universal theorem in Lean. Give an example of each scope.\end{exercise}


% ===== Week 17 =====

\chapter{Entering an unfamiliar proof without starting over}\label{ch:bridge}
\goal{Build a prerequisite map for a particular argument, learn only its missing nodes, and then reconstruct the whole proof. Use the same method to read a research repository without mistaking an entry-point file for all of its mathematics.}

\section{A new theorem, with a familiar foundation}
Imagine encountering the following sentence in a paper:
\begin{quote}
Every sequence in a complete, totally bounded metric space has a convergent subsequence.
\end{quote}
You know what a metric, completeness, and convergence mean from Chapters~11--12. ``Totally bounded'' and ``subsequence'' are new. The unhelpful response is to request a complete course in topology. A more useful response asks what these two ideas do in this particular argument.

\textbf{Reader:} State the proof's mechanism without using any new theorem names.

\textbf{Assistant:} Repeatedly keep infinitely many terms inside a smaller region. Choose later and later terms from those nested selections. They become Cauchy; completeness supplies their limit.

\textbf{Reader:} Which missing facts make those steps legitimate?

\textbf{Assistant:} We need finite coverings by small balls, the fact that one of finitely many classes containing infinitely many indices must itself be infinite, and a way to choose strictly increasing indices.

This is a proposed dependency map, not yet a proof. We will now supply every missing mathematical link. No general theorem about compactness is needed. The minimum background is relative to this chosen proof; another proof may require a different collection of concepts.

\section{Capsule one: small regions that cover the space}
For a metric space $(X,d)$, a \emph{ball} with center $c\in X$ and radius $r>0$ is
\[
 B(c,r)=\set{x\in X:d(x,c)<r}.
\]
It need not look round in a drawing. The definition only uses distance. If $x,y$ lie in the same ball, the triangle inequality gives
\[
 d(x,y)\le d(x,c)+d(c,y)<r+r=2r.
\]
This is the property of balls that the proof will use: one small ball controls distances between all of its points. A bound on all pairwise distances within a set is called a diameter bound.

\par\noindent\begin{minipage}{\linewidth}
\begin{definition}
A metric space is \emph{totally bounded} if, for every $r>0$, finitely many balls of radius $r$, with centers in the space, cover the entire space.
\end{definition}
\end{minipage}\par

The centers and their number may depend on $r$. We are not claiming that one finite list works at all resolutions.

For example, $[0,1]$ is totally bounded with ordinary distance. Given $r>0$, choose a positive integer $N$ with $1/N<r$. Use centers $0,1/N,\ldots,1$. Every point between two consecutive centers is within $1/N$ of at least one of them, and hence within $r$. The endpoints are centers themselves. This gives a finite cover explicitly.

A metric space is \emph{bounded} when some fixed finite number is an upper bound for the distance between every pair of its points. This is weaker than total boundedness. On $\NN$, define $d(m,n)=0$ if $m=n$ and $d(m,n)=1$ otherwise. This is a metric: for distinct endpoints, an intermediate point cannot make both distances zero, so their sum is at least one. All distances are at most one, yet a ball of radius $1/2$ contains only its center. Finitely many such balls cannot cover infinitely many natural numbers. Thus this space is bounded but not totally bounded.

\section{Capsule two: retaining infinitely many indices}
A sequence $(x_n)$ has infinitely many \emph{indices}, even if some or all of its values are repeated. We must track indices, not just distinct points.

Suppose infinitely many indices are distributed among finitely many classes. If each class had only finitely many indices, their union would be finite: add their finite sizes to obtain a finite upper bound. That contradicts the assumed infinitely many indices. Therefore at least one class contains infinitely many indices. This is the infinite version of the pigeonhole mechanism from Chapter~7.

If covering balls overlap, assign each index to the first ball in a fixed finite list that contains its point. The classes then form a partition. Alternatively, use the fact that a finite union of finite sets is finite directly on the sets of indices in the balls. Either route justifies keeping one ball containing infinitely many of the current indices.

\section{Capsule three: a subsequence must preserve order}
A \emph{subsequence} is a sequence of the form
\[
 x_{n_1},x_{n_2},x_{n_3},\ldots,
 \qquad n_1<n_2<n_3<\cdots.
\]
We may skip terms but may not reorder them or reuse a single index indefinitely. Every infinite subset of $\NN$ contains an index larger than any specified integer $M$: otherwise it would be contained in the finite set $\set{0,\ldots,M}$.

Consequently, if an infinite set of permissible indices is available at each stage, we can choose a new index larger than the previous one. This small fact is exactly what turns nested infinite selections into a subsequence rather than an unordered collection of points.

\subsection*{Plan the construction with three distinct kinds of object}
There will be points of the metric space, balls containing points, and sets of sequence indices. They are not interchangeable. A ball contains values such as $x_7$. A selected index set contains the natural number seven because its value $x_7$ lies in the selected ball. Repeated equal values can therefore contribute many different indices.

Start with every index permissible. At the first stage, cover the space by finitely many balls of radius $1/4$ and keep the indices whose points lie in one ball containing infinitely many of them. Any two retained values are within $1/2$ of one another. At the second stage, cover the space by balls of radius $1/8$, but consider only indices retained at the first stage. Keep infinitely many of those inside a single new ball. Their values are now within $1/4$ of one another.

Continuing, the radii at stage $k$ are $2^{-(k+1)}$, so twice a radius is $2^{-k}$. The factor of two comes from the two-leg distance estimate within a ball. We shrink the radius by an extra factor of two precisely to get this convenient bound.

The \emph{sets of indices} are nested because each stage discards some previously permissible indices and introduces none. The chosen balls themselves need not be nested as subsets of the space. We only need all later chosen values to retain the earlier distance guarantees. The proof will use nestedness of the index sets to establish that fact.

\section{The complete proof}
\par\noindent\begin{minipage}{\linewidth}
\begin{theorem}\label{thm:subsequence}
Let $(X,d)$ be a nonempty, complete, totally bounded metric space. Every sequence $(x_n)_{n\ge0}$ in $X$ has a subsequence converging to a point of $X$.
\end{theorem}
\end{minipage}\par

\begin{proof}
We first construct nested infinite sets of indices. Set $S_0=\NN$. For each integer $k\ge1$, cover $X$ by finitely many balls of radius $2^{-(k+1)}$. The indices in $S_{k-1}$ are infinite in number. By the preceding pigeonhole argument, one of these balls contains $x_n$ for infinitely many $n\in S_{k-1}$. Let $S_k$ be that infinite set of indices. Then
\[
 S_k\subseteq S_{k-1},
 \qquad d(x_i,x_j)<2\cdot2^{-(k+1)}=2^{-k}
 \quad(i,j\in S_k).
\]
The inclusion matters: the selections at later stages retain the distance control obtained at earlier stages.

Choose $n_1\in S_1$. Once $n_{k-1}$ is chosen, choose $n_k\in S_k$ with $n_k>n_{k-1}$. Such an index exists because $S_k$ is infinite. This constructs a subsequence $(x_{n_k})$.

Here $k$ counts the extraction stage, while $n_k$ is an index in the original sequence. For example, the first few chosen indices might be $n_1=4,n_2=9,n_3=15$. The extracted terms would be $x_4,x_9,x_{15}$, in that order. The index values need not follow a simple formula. What matters is that each $n_k$ is chosen from $S_k$ and is strictly larger than the preceding one.

For $\ell\ge k$, nestedness means $S_\ell\subseteq S_k$. Since $n_\ell\in S_\ell$, it follows that $n_\ell\in S_k$. Thus not only the $k$th chosen term but \emph{every later chosen term} obeys the distance control from stage $k$. This is the bridge from one-stage choices to a Cauchy tail.

We show it is Cauchy. Given $\eps>0$, choose $K$ with $2^{-K}<\eps$. For $\ell,m\ge K$, nestedness gives $S_\ell\subseteq S_K$ and $S_m\subseteq S_K$. Thus $n_\ell,n_m\in S_K$, and
\[
 d(x_{n_\ell},x_{n_m})<2^{-K}<\eps.
\]
This is precisely the Cauchy condition. Completeness now supplies a point $x_*\in X$ to which the subsequence converges. No further extraction is needed.

We did not choose an index belonging to every set $S_k$ at once. There might be no such index: the nested infinite sets $\{k,k+1,k+2,\ldots\}$ have empty intersection. Instead we chose a new later index from each successive set and proved that any sufficiently late pair belongs together to an earlier selected set. That is all the Cauchy definition asks for. Replacing this step by ``take an index in the intersection'' would introduce a false extra claim.
\end{proof}

\pauseq{Where is each hypothesis used?}{Total boundedness gives a finite cover at every smaller radius. Finiteness lets the pigeonhole argument retain infinitely many indices. Nestedness transfers earlier distance bounds to all later choices. Completeness places the limit inside $X$.}

\section{Why the missing assumptions cannot be ignored}
In the discrete metric on $\NN$, the sequence $0,1,2,\ldots$ has no Cauchy subsequence: distinct chosen terms stay distance one apart. The space is complete. To see why, apply the Cauchy condition with tolerance $1/2$. Any two terms beyond the resulting threshold must have distance less than $1/2$. Their distance can only be zero or one, so it must be zero: all those late terms are equal. An eventually constant sequence converges to that constant member of $\NN$. However, the space is not totally bounded. Completeness alone does not force a convergent subsequence.

The interval $(0,1)$ is totally bounded. Given $r>0$, partition $[0,1]$ into $N$ equal pieces with $1/(2N)<r$, and use the midpoints of these pieces as centers. All centers belong to $(0,1)$, and every point of $(0,1)$ is within $1/(2N)$ of a center. However, the sequence $x_n=1/(n+2)$ has no subsequence converging inside $(0,1)$: every subsequence still converges to zero in $\RR$, and limits are unique. This space lacks completeness.

These counterexamples do more than decorate the theorem. They identify two different mechanisms: having enough small finite descriptions and having a place for Cauchy limits to reside. A useful question about a new hypothesis is therefore not only ``What does it mean?'' but ``Which mechanism fails without it?''

\subsection*{Check what was learned, not only what was named}
For a new concept, first state the problem it solves. Total boundedness gives a finite list of small regions at any requested scale. Then give the definition with its quantifiers. The list and its length may change when the requested radius changes, but for one radius the same list must cover every point.

Next test the boundary. The discrete metric shows that an overall bound on distances does not provide those finite small-scale descriptions. The open interval shows that finite small-scale descriptions alone do not supply all missing limits. Finally return to the exact proof line where the concept is used. If you can name the definition but cannot explain that line, the prerequisite has not yet been learned sufficiently for this purpose.

For an unfamiliar research proof, an assistant's first prerequisite list is only a proposal. Discovering a further missing step is not a failure of the method. Add that step, learn it, and return to the argument. The list should become sufficient for the chosen proof through inspection and reconstruction; it should not be treated as automatically minimal or complete because an AI generated it.

\section{The prerequisite map after the work is done}
The map is now concrete:
\[
\begin{gathered}
 \text{triangle inequality}\ \longrightarrow\ \text{small-ball diameter bound},\\
 \text{finite cover + infinite indices}\ \longrightarrow\ \text{nested infinite selections},\\
 \text{unbounded index sets}\ \longrightarrow\ \text{increasing subsequence},\\
 \text{nested distance bounds}\ \longrightarrow\ \text{Cauchy subsequence},\\
 \text{completeness}\ \longrightarrow\ \text{limit in the space}.
\end{gathered}
\]
The first and last lines were already familiar. The middle lines were the minimal missing background for this proof. After learning each new node, we immediately used it in the argument. That return to the target is essential: otherwise just-in-time learning becomes another disconnected vocabulary list.

For a genuinely unfamiliar paper, mark a prerequisite as \emph{understood} only when you can explain its use and solve a small transfer problem. Mark an external theorem as \emph{borrowed} until its hypotheses and application have been checked; learning its full proof can be a separate, explicit task. A proof map is allowed to contain borrowed results, but it must not conceal them.

\section{A repository reading that follows the same method}
We now return to the repository that motivated the course. Its documentation distinguishes manuscripts and verification scopes; a repository listing is not, by itself, a guarantee about every claim in every paper~\cite{oairepo}. The appropriate unit of reading is an exact theorem together with its definitions and dependencies, not a title or an impressive final line.

For Family~173, the statement concerns a nonempty finite oriented graph and a vertex whose second out-neighborhood has at least as many vertices as its first. Chapter~7 supplies the mathematical vocabulary; Chapter~16 explains the exclusions in second-neighborhood membership~\cite{oai173}. In the formal signature, \code{Fintype V} supplies finiteness, \code{Nonempty V} rules out the empty graph, and \code{DecidableEq V} provides a way to decide equality for the finite constructions. The edge relation has type \code{V -> V -> Prop}.

The entry point imports modules for reduction, pruning, and an extremal contradiction. Its proof assumes a counterexample, obtains a transformed finite oriented relation from a reduction theorem, and applies the pruning and extremal results to obtain a contradiction~\cite{oaimain}. The reduction's output is a relation on pairs of finite indices with additional properties, including positive indegree and a strict subset-growth condition~\cite{oaireduction}. Understanding those property names requires opening their definitions; translating the names into plausible English is not enough.

A responsible reading notebook therefore has three different entries: the final statement is decoded; the logical role of the reduction is identified; the mathematical proof of the reduction and the meaning of its additional predicates remain tasks until their source has been read. The full repository argument is not reproduced or independently verified in this book. We have prepared the vocabulary and method for entering it, not substituted a short dependency sketch for that argument.

The comparison configuration identifies the challenge module, the solution module, the theorem name, and permitted axioms~\cite{oaiconfig}. A placeholder in a challenge template is not equivalent to a placeholder in the submitted solution's proof. Conversely, a theorem name in a solution file is not sufficient evidence that it matches the challenge. The configuration and the actual checking procedure matter.

For Family~179, Chapter~14 lets you decode the matrix and correlation conditions and prove some necessary restrictions yourself. The repository's scope note reports a stronger classification and distinguishes its even-length Barker statement from odd-length claims outside that selected formalization~\cite{oai179}. You can now ask a precise next question: which lemma rules out the candidate orders $4u^2$ with $u>1$ that our elementary argument still permits? That is a mathematically informed question, unlike asking an assistant simply to ``explain the whole repository.''

\begin{lab}{Request background, then test whether it is sufficient}
Begin with Theorem~\ref{thm:subsequence} while temporarily hiding its proof. Ask the assistant for a dependency map using only the knowledge from Chapter~12. Reject a response that merely lists ``compactness, topology, sequences.'' Require an example, a nonexample, and the exact role of every proposed concept. After reconstructing the proof, close the explanation and answer: why must the sets of indices be nested, why must indices increase, and why is boundedness not enough? A map that leaves one of these questions unanswered is not yet sufficient.

Then open one of the two repository scope notes and its named entry point. Produce a one-page reading record distinguishing what you can now reconstruct from what is still borrowed. Use the source file names and actual theorem identifiers. Do not report a repository build or formal verification unless you have performed it and retained the result.
\end{lab}

\section*{Exercises}
\begin{exercise}\label{ex:17.1}Explain why infinitely many sequence indices inside a ball do not imply infinitely many distinct points inside that ball. Does the proof need distinct points?\end{exercise}
\begin{exercise}\label{ex:17.2}An assistant chooses an arbitrary $n_k\in S_k$ without requiring $n_k>n_{k-1}$. Identify the gap and repair the choice. Explain why the repair is always possible.\end{exercise}
\begin{exercise}\label{ex:17.3}Prove that the discrete metric on $\NN$ is complete but not totally bounded. Use it to explain the difference between completeness and the conclusion of Theorem~\ref{thm:subsequence}.\end{exercise}
\begin{exercise}\label{ex:17.4}Give a full finite-ball-cover argument for total boundedness of $[a,b]\subset\RR$, where $a<b$. State exactly what the number of balls may depend on.\end{exercise}
\begin{exercise}\label{ex:17.5}A directed graph has edges $v\to u$, $u\to w$, and $v\to w$. Compare the two-edge-reachability and exact-second-neighborhood definitions at $v$. Explain the formalization error caused by replacing one with the other.\end{exercise}
\begin{exercise}\label{ex:17.6}Write a minimal-background request for the contraction theorem and one for the subsequence theorem. For each, identify a concept needed in that proof but not in the other. Then state what your repository reading record may legitimately claim without a completed repository build.\end{exercise}


% ===== Week 18 =====

\chapter{A complete AI-assisted investigation}\label{ch:capstone}
\goal{Turn a theorem into a small investigation: vary an assumption, derive a new target statement, test it, prove it, examine sharpness, and report its scope. Finish with a reproducible mathematical account rather than an unexamined AI transcript.}

\section{The question that starts the project}
The contraction theorem assumes exact updates $x_{n+1}=T(x_n)$. A computation, a measurement-driven update, or a numerical subroutine may instead produce an approximate update. Does the iteration still approach the fixed point? If not, can its error be controlled?

This is a useful research-style question because it changes one clearly identified feature while retaining the earlier structure. It does not require inventing a subject from nothing. We will complete the investigation as a worked capstone. The resulting bounds are standard mathematical consequences of contraction; this chapter does not claim them as new research.

Let $(X,d)$ be a nonempty complete metric space and let $T:X\to X$ satisfy
\[
 d(Tx,Ty)\le q\,d(x,y),\qquad 0<q<1.
\]
Theorem~\ref{thm:banach} gives a unique fixed point $x_*$. Suppose an approximate sequence $(y_n)$ in $X$ satisfies
\[
 d(y_{n+1},T(y_n))\le\eps_n,\qquad \eps_n\ge0.
\]
The errors may vary with $n$ and need not be random. Membership $y_n\in X$ remains an assumption: an error bound alone does not force an approximate point to stay in a constrained domain.

\section{A dialogue that changes the claim, not the evidence}
\textbf{Reader:} I expect small update errors still to give convergence to $x_*$. Try to disprove this before proving it.

\textbf{Assistant:} A constant bias can move the limiting value. For $T(x)=qx$ on $\RR$, the update $y_{n+1}=qy_n+\eta$ tends toward $\eta/(1-q)$, not zero, when $\eta>0$.

\textbf{Reader:} Then the target should be an error bound. Can the same idea also distinguish bounded errors from errors that tend to zero?

\textbf{Assistant:} First bound the distance at the next step by $q$ times the present distance plus the update error. Then unfold the recurrence.

The assistant has proposed a mechanism. We now derive it independently, prove the unfolded form, and test whether its constant can be improved. The revised question is more informative than the initial conjecture: it separates stability, convergence, and bias.

\section{The one-step estimate}
Write $e_n=d(y_n,x_*)$. Insert the exact update between the approximate point and the fixed point:
\begin{align*}
 e_{n+1}
 &=d(y_{n+1},x_*)\\
 &\le d(y_{n+1},T(y_n))+d(T(y_n),x_*)
 &&\text{triangle inequality}\\
 &=d(y_{n+1},T(y_n))+d(T(y_n),T(x_*))
 &&\text{$T(x_*)=x_*$}\\
 &\le\eps_n+q\,d(y_n,x_*)
 &&\text{update and contraction bounds}\\
 &=q e_n+\eps_n.
\end{align*}
Every assumption has a visible role. No independence, differentiability, or coordinate representation was used.

\subsection*{Unfold three steps before writing the general sum}
The one-step estimate does not yet display how old errors persist. Substitute it into itself, remembering that $q>0$ preserves inequalities on multiplication:
\begin{align*}
 e_1&\le qe_0+\eps_0,\\
 e_2&\le qe_1+\eps_1
       \le q^2e_0+q\eps_0+\eps_1,\\
 e_3&\le qe_2+\eps_2
       \le q^3e_0+q^2\eps_0+q\eps_1+\eps_2.
\end{align*}
At step three, the initial error has been multiplied by $q$ three times. The error $\eps_0$ introduced in the first update has been multiplied by $q$ twice afterward. The most recent update error $\eps_2$ has not yet been reduced by a later update, so its coefficient is one.

At general step $n$, error $\eps_j$ was introduced while producing $y_{j+1}$. There are $n-(j+1)=n-1-j$ subsequent updates before $y_n$. Its weight is therefore $q^{n-1-j}$. The latest index is $j=n-1$, giving $q^0=1$. The oldest is $j=0$, giving $q^{n-1}$. These endpoint checks explain the summation limits in the next statement.

The theorem separates two error hypotheses. A uniform bound $\eps_j\le\eta$ allows a persistent nonzero disturbance. The condition $\eps_j\to0$ requires the disturbances eventually to become arbitrarily small. The second can imply convergence to the original fixed point even when the accumulated sum of disturbances is not finite. We will prove the conclusions separately rather than interchange these assumptions.

\par\noindent\begin{minipage}{\linewidth}
\begin{theorem}\label{thm:inexact}
Under the preceding assumptions, for every integer $n\ge1$,
\[
 e_n\le q^n e_0+\sum_{j=0}^{n-1}q^{n-1-j}\eps_j.
\]
If $\eps_j\le\eta$ for every $j$, where $\eta\ge0$, then
\[
 e_n\le q^n e_0+\eta\frac{1-q^n}{1-q}.
\]
If instead $\eps_n\to0$, then $y_n\to x_*$.
\end{theorem}
\end{minipage}\par

\begin{proof}
For $n=1$, the first inequality is precisely $e_1\le qe_0+\eps_0$. Suppose the formula holds at $n$. The one-step bound gives
\begin{align*}
 e_{n+1}
 &\le q\left(q^n e_0+\sum_{j=0}^{n-1}q^{n-1-j}\eps_j\right)+\eps_n\\
 &=q^{n+1}e_0+\sum_{j=0}^{n-1}q^{n-j}\eps_j+q^0\eps_n\\
 &=q^{n+1}e_0+\sum_{j=0}^{n}q^{n-j}\eps_j.
\end{align*}
This is the claimed formula with $n+1$ in place of $n$, so induction proves it for all $n\ge1$.

When $\eps_j\le\eta$, all weights are nonnegative. Replace each $\eps_j$ by $\eta$ and reverse the order of the powers:
\[
 \sum_{j=0}^{n-1}q^{n-1-j}\eps_j
 \le\eta\sum_{k=0}^{n-1}q^k
 =\eta\frac{1-q^n}{1-q}.
\]
The last identity was derived in Chapter~11. This proves the uniform-error bound.

\emph{Prepare the convergence argument.} Small later errors and possibly large earlier errors need different treatment. Fix a stage $J$ after which errors are small. The finitely many errors before $J$ have fixed values. Their influence will be reduced by more and more powers of $q$ as $n$ grows. The later errors may be numerous, but each is small and their geometric weights have a bounded sum. This suggests splitting the sum at $J$.

The order of choices is important. First a target tolerance is supplied. Next choose $J$ so later disturbances are sufficiently small. Only after $J$ is fixed do we choose how far out to go to suppress its finite prefix. Choosing a new $J$ for each term in the final comparison would not establish one threshold for an entire tail.

Finally suppose $\eps_n\to0$. Let $\delta>0$ be an arbitrary target tolerance. Choose an integer $J\ge1$ such that
\[
 \eps_j<\frac{(1-q)\delta}{2}\qquad(j\ge J).
\]
For $n>J$, split the error sum into an early finite part and a later part:
\begin{align*}
 e_n
 &\le q^{n-J}\left(q^J e_0+
       \sum_{j=0}^{J-1}q^{J-1-j}\eps_j\right)
       +\sum_{j=J}^{n-1}q^{n-1-j}\eps_j\\
 &\le q^{n-J}C_J+\frac{(1-q)\delta}{2}
        \sum_{k=0}^{n-1-J}q^k\\
 &\le q^{n-J}C_J+\frac\delta2,
\end{align*}
where the parentheses define a fixed finite nonnegative number $C_J$.

To check the factoring in the first line, write the early terms before extracting their common power:
\[
 q^ne_0+q^{n-1}\eps_0+q^{n-2}\eps_1+\cdots+q^{n-J}\eps_{J-1}.
\]
The smallest exponent among them is $n-J$. Factoring $q^{n-J}$ leaves
\[
 C_J=q^Je_0+q^{J-1}\eps_0+q^{J-2}\eps_1+\cdots+\eps_{J-1}.
\]
Once the tolerance and $J$ have been chosen, this number does not depend on $n$. It is finite because it adds finitely many finite nonnegative real values.

For the late part, $j=J$ through $n-1$ gives powers $q^{n-1-J}$ down to $q^0$. Their sum is at most $1/(1-q)$. Multiplying that bound by the allowed late-error size $(1-q)\delta/2$ leaves $\delta/2$. The other half of the tolerance is reserved for the early part. This is why that particular tolerance was requested of the error sequence. Since $0<q<1$, choose $N>J$ large enough that $q^{n-J}C_J<\delta/2$ for every $n\ge N$; if $C_J=0$, this term already vanishes. Then $e_n<\delta$ for $n\ge N$. This is convergence of $y_n$ to $x_*$.
\end{proof}

The finite-prefix split is a reusable technique. Early errors can be large because only finitely many occur and their effects are eventually damped. Later errors are controlled uniformly because they have become small. The proof does not require the accumulated error sums $\eps_0+\cdots+\eps_n$ to stay below one finite bound as $n$ grows. It controls the \emph{weighted} influence of errors, not their unweighted total.

\section{Is the constant an artifact of the proof?}
Take $X=\RR$, $T(x)=qx$, $y_0=0$, and add exactly $\eta$ at each update. The recurrence has the explicit solution
\[
 y_n=\eta\sum_{k=0}^{n-1}q^k
     =\eta\frac{1-q^n}{1-q}.
\]
This follows either from induction or by substituting the formula into $y_{n+1}=qy_n+\eta$.

For the induction step, assume $y_n=\eta(1-q^n)/(1-q)$. Then
\begin{align*}
 y_{n+1}
 &=q\,\frac{\eta(1-q^n)}{1-q}+\eta\\
 &=\frac{\eta(q-q^{n+1})+\eta(1-q)}{1-q}\\
 &=\frac{\eta(1-q^{n+1})}{1-q}.
\end{align*}
At $n=0$ the formula gives zero, as required. The exact formula is therefore established for every finite step; the limiting value follows afterward from $q^n\to0$. We do not infer the formula merely from the appearance of the first few terms. Since $x_*=0$ and $y_n\ge0$, the uniform-error bound is an equality. Thus its limiting factor $1/(1-q)$ cannot be decreased for all systems satisfying these assumptions.

A sharp bound need not predict convergence of every allowed approximate sequence. For example, let $T(x)=0$ and choose $y_n=(-1)^n\eta$ with $\eta>0$. This map satisfies the contraction inequality for any chosen $q\in(0,1)$, and
\[
 d(y_{n+1},T(y_n))=\eta.
\]
Yet $y_n$ alternates between two different values and does not converge. Bounded errors give a distance estimate, not convergence by themselves. Vanishing errors provide the extra conclusion in Theorem~\ref{thm:inexact}.

\section{A numerical experiment with exact arithmetic}
Choose $T(x)=(3x+1)/4$, whose fixed point is $1$, and use the biased update
\[
 y_{n+1}=\frac{3y_n+1}{4}+\frac1{100},\qquad y_0=0.
\]
The constants are $q=3/4$ and $\eta=1/100$. The approximate recurrence is affine:
\[
 y_{n+1}=\frac34 y_n+\frac{13}{50}.
\]
Its own fixed point is $(13/50)/(1/4)=26/25$. Subtracting this value from both sides gives
\[
 y_n=\frac{26}{25}\left(1-\left(\frac34\right)^n\right).
\]
The theorem's error certificate is
\[
 \abs{y_n-1}\le
 \left(\frac34\right)^n+
 \frac1{25}\left(1-\left(\frac34\right)^n\right).
\]
Notice that the true error can first decrease toward zero and later increase toward $1/25$. A plot that stops near the smallest error would obscure the bias. The formula and bound tell us what a finite plot alone cannot establish.

The companion computes this experiment using rational numbers, checks the exact closed form at every tested index, and compares the exact error with the exact certificate. Decimal output is only for display. Changing the sign pattern of the errors gives a second experiment: compare constant bias, alternating errors, and a sequence of errors tending to zero. The theorem predicts a bound for all three, but convergence to the original fixed point only under the vanishing-error condition.

\subsection*{What ``sharp'' means in this example}
Suppose someone proposes replacing the eventual error factor $1/(1-q)$ by a smaller constant $K$ that should work for every allowed iteration with this fixed $q$. Choose $\eta>0$ in the constant-bias example. Its actual error tends to $\eta/(1-q)$, which is strictly larger than $K\eta$. Thus sufficiently late errors exceed the proposed smaller bound. One permitted example refutes the universal improvement.

This does not say the original bound is attained in every application. In the biased affine example approaching $26/25$, some early errors cancel, so the certificate is conservative. A bound is sharp over a class when the class contains an example preventing a uniformly stronger bound. For a particular smaller class with additional structure, an improvement may still be possible. That distinction supplies a precise next research question: what extra assumption would rule out the example that obstructs an improved constant?

\section{The finished investigation as a short paper}
A mathematical account of this project has four connected parts. The introduction identifies the mismatch between exact and approximate updates and distinguishes stability from convergence. The theory states the domain, contraction factor, update-error model, unfolded estimate, sharpness example, and vanishing-error result. The results present reproducible experiments as illustrations of these established statements, not as their replacements. The conclusion explains the limitations: approximate iterates must stay in $X$, $q<1$ is essential to these geometric estimates, and persistent bounded errors need not disappear.

The contribution statement for this course project should say that it reconstructs and analyzes a standard perturbation argument. It should not claim a novel extraction of a theorem merely because the learner or assistant did not previously know it. For a research submission, compare the exact statement and assumptions with primary literature, trace citations to their actual results, and revise the novelty claim after that comparison. A search that returns no match is evidence about that search, not a proof of novelty.

\section{Your independent capstone}
Choose one thread and keep the first deliverable narrow enough that every assertion can be checked. An assignments project can add integer capacities to tasks and derive the appropriate shortage condition by replacing a capacity-$c$ task with $c$ distinct slots. An iteration project can analyze a specified decaying-error schedule. A sign-sequence project can give a reproducible classification at one additional finite length and investigate which symmetries preserve its conditions. These are starting points, not claims of new results.

Submit a precise statement, a prerequisite map, a proof or counterexample, executable checks when relevant, an assumption-and-evidence record, and a short explanation of what is genuinely new to the literature versus new only to you. The oral component should ask you to explain one delicate step, change an assumption, and predict what fails. Your ability to answer matters more than the fluency of the saved AI conversation.

An effective final review uses the assistant in three passes with different tasks. First request a faithful reconstruction of the argument. Next request counterexamples to stronger variants and scrutiny of each hypothesis. Finally ask for a correspondence check between the prose, formulas, code, and formal statements. These passes may still share blind spots; they do not constitute independent human peer review. Their value lies in the concrete obligations they expose for you to check.

\begin{lab}{From a plausible assertion to a defensible result}
Start with the false claim ``bounded update errors do not affect the limiting fixed point.'' Produce the constant-bias counterexample before asking for a repair. Derive the one-step inequality yourself, ask the assistant to propose the unfolded form, and verify its indexing at $n=1,2,3$. Then prove the induction step, test the sharp example, and explain why alternating bounded errors can prevent convergence. Finally close the chapter and reconstruct the finite-prefix argument for vanishing errors. This completes the collaboration loop: you can now explain the result rather than merely possess its generated proof.
\end{lab}

\section*{Exercises}
\begin{exercise}\label{ex:18.1}Write out the error bounds at $n=1,2,3$ before using summation notation. Identify the weight multiplying the earliest and the latest update errors.\end{exercise}
\begin{exercise}\label{ex:18.2}For $T(x)=qx$, $y_0=0$, and $y_{n+1}=qy_n+\eta$, prove the explicit formula and explain why the asymptotic factor $1/(1-q)$ is sharp.\end{exercise}
\begin{exercise}\label{ex:18.3}Verify every assumption and the nonconvergence conclusion in the alternating-error example $T(x)=0$, $y_n=(-1)^n\eta$. Explain why a valid error bound does not settle convergence.\end{exercise}
\begin{exercise}\label{ex:18.4}Suppose $\eps_j\le a q^j$ with $a\ge0$ and the same $q\in(0,1)$ as the contraction. Derive $e_n\le q^n e_0+a n q^{n-1}$ and prove that this upper bound tends to zero.\end{exercise}
\begin{exercise}\label{ex:18.5}In a finite assignment instance, each task $r$ has a nonnegative integer capacity $c_r$. A student assigned to $r$ uses one unit of capacity. Prove that all students can be assigned exactly when $|S|\le\sum_{r\in N(S)}c_r$ for every set $S$ of students.\end{exercise}
\begin{exercise}\label{ex:18.6}Prepare a final reading-and-investigation record for the biased iteration in this chapter: exact question, minimal prerequisites, proof skeleton, counterexample to the initial claim, one exact numerical check, and the permitted novelty statement. Then name one further question whose answer is not established by the existing proof.\end{exercise}


\appendix
\renewcommand{\chapterkind}{APPENDIX}


\chapter{Solutions: statements, maps, and structure}\label{app:sola}
These solutions are part of the learning material, not merely an answer key. Compare the point where your argument stopped with the corresponding inference below. After reading a solution, reproduce its mechanism on a changed example before treating the idea as learned.

\section{Week 1}
\sol{1.1}{At $x=4$, $2(x+3)=2(4+3)=14$, whereas $2x+3=2\cdot4+3=11$. In the first expression we add three before doubling, so the added three is also doubled. In the second we double only the original input and then add three. The difference is three, not a contradiction: these are different functions of the same input.}
\sol{1.2}{The conclusion asks for an odd sum. By the definition, we should aim to write that sum as twice an integer plus one.

Name the even input $a$ and the odd input $b$. Their definitions supply integers $k,\ell$ with $a=2k$ and $b=2\ell+1$. The witnesses may be different because the two inputs were not required to be equal. Substitute and regroup:
\begin{align*}
 a+b&=2k+(2\ell+1)\\
 &=2k+2\ell+1\\
 &=2(k+\ell)+1.
\end{align*}
The number $k+\ell$ is an integer, so the final expression meets the definition of odd. It is the requested witness. For example, $a=4,b=-3$ have witnesses $k=2,\ell=-2$; the sum is one, witnessed by $k+\ell=0$. The example illustrates the proof but does not replace its arbitrary-input argument.}
\sol{1.3}{The proposed calculation uses the same odd number twice. Squaring $2k+1$ therefore covers pairs of equal inputs only. It does not cover, for example, the product of three and five. Correct algebra can still have insufficient scope.

For the general claim, name the inputs independently: $a=2k+1$ and $b=2\ell+1$, with integer $k,\ell$. Then distribute one bracket at a time:
\begin{align*}
 ab&=(2k+1)(2\ell+1)\\
 &=2k(2\ell+1)+(2\ell+1)\\
 &=4k\ell+2k+2\ell+1\\
 &=2(2k\ell+k+\ell)+1.
\end{align*}
The expression in parentheses is an integer because products and sums of integers are integers. This gives the defining form of an odd integer and proves the claim for every allowed pair. Setting $k=\ell$ now recovers the square case as one consequence, not as a substitute for the general proof.}
\sol{1.4}{The sum is
\[
 (2a_1+1)+(2a_2+1)+(2a_3+1)
 =(2\cdot2+1)+(2(-1)+1)+(2\cdot4+1)
 =5-1+9=13.
\]
There are three terms because the index runs from one through three, not because the upper limit is a multiplicative factor.}
\sol{1.5}{A zero product requires $x=0$ or $x-2=0$. Thus the only candidates are $0$ and $2$. Substitution verifies both. Dividing by $x$ discards the valid candidate $0$, because division is permitted only when $x\ne0$. A division-based proof is complete only after separately checking the excluded case.}
\sol{1.6}{Use
\[
 f(n)=(n-1)(n-2)\cdots(n-10).
\]
At an integer $n$ from one through ten, one factor is zero. At $n=11$, all factors are positive, giving $10\cdot9\cdots1=3\,628\,800$. The finite observations do not determine the behavior of this formula outside the tested inputs.}

\section{Week 2}
\sol{2.1}{The converse is ``If $x^2=9$, then $x=3$.'' It is false at $x=-3$. The contrapositive is ``If $x^2\ne9$, then $x\ne3$.'' It is true: $x=3$ would force $x^2=9$ and contradict its assumption. The original implication is also true by substitution.}
\sol{2.2}{Work from the outer promise inward. The original sentence says that \emph{every} integer $n$ has at least one satisfactory response $m$. Its failure requires \emph{some} integer $n$ for which no response works. Thus the first two words of the negation are ``There exists,'' and the response clause must say ``for every integer $m$.''

For a proposed response to work, both $m>n$ and evenness are required. It fails when at least one requirement fails: $m\le n$ or $m$ is not even. Equality belongs to the first failure because the original comparison was strict. The full negation is
\[
 \exists n\in\ZZ\ \forall m\in\ZZ:\quad
 (m\le n)\lor\neg(m\text{ is even}).
\]
In words: there is an integer $n$ such that every integer $m$ either is at most $n$ or is not even. The ``or'' is inclusive. An $m$ that fails both conditions is still an unsuccessful response; there is no reason to exclude it. We are describing how the original statement would fail, not asserting that this negation is true.}
\sol{2.3}{The first statement is true: after receiving $x$, choose $y=-x$. Their sum is zero. The second is false. If a single $y$ worked for every $x$, using $x=0$ would give $y=0$, while using $x=1$ would give $y=-1$. One number cannot satisfy both requirements. Equivalently, after any proposed $y$, choose $x=1-y$, whose sum with $y$ is one.}
\sol{2.4}{The value $x=-3$ satisfies $x^2=9>4$ but not $x>2$. The corrected conclusion is $x<-2$ or $x>2$, equivalently $|x|>2$. To check the classification, when $-2\le x\le2$ the magnitude of $x$ is at most two, so $x^2\le4$. Outside that interval the magnitude exceeds two, so squaring gives $x^2>4$.}
\sol{2.5}{Let both students accept only task $1$. Each has an acceptable task, but distinct assignments would require two different tasks drawn from a one-element set. The stronger condition checks groups: the acceptable tasks of every set of students must number at least as many as those students. For this pair it reads $1\ge2$, which fails. Chapter~9 proves this groupwise condition is also sufficient in the finite setting.}
\sol{2.6}{Choose a positive integer $N>1/\eps$. For any integer $n\ge N$, positivity gives $1/n\le1/N<\eps$. The witness $N$ may depend on the tolerance already supplied. It must work for every later $n\ge N$, so choosing a different $N$ after seeing each $n$ would not establish the stated eventual guarantee.}

\section{Week 3}
\sol{3.1}{The union is $\{1,2,3,4\}$, the intersection is $\{2,4\}$, and $A\setminus B=\{1\}$. A union lists each member once even if it belongs to both sets. The difference retains members of the first set that do not belong to the second.}
\sol{3.2}{We must prove surjectivity of $g\circ f$, so start from the \emph{output} end. Let $f:A\to B$ and $g:B\to C$ be surjective, and choose an arbitrary target $c\in C$. The goal is to find $a\in A$ such that $g(f(a))=c$.

Surjectivity of $g$ applies to $c$, because $c$ belongs to its codomain $C$. It supplies some $b\in B$ with $g(b)=c$. Surjectivity of $f$ now applies to this particular $b$, because $b$ belongs to its codomain $B$. It supplies some $a\in A$ with $f(a)=b$. Substitution checks the candidate:
\[
 (g\circ f)(a)=g(f(a))=g(b)=c.
\]
Every target $c$ can be reached in this way, so the composition is surjective. Witnesses are found in reverse order from function execution. We do not assume either witness is unique; surjectivity promises existence only.}
\sol{3.3}{If $2m=2n$ for integers $m,n$, division by the nonzero number two gives $m=n$. Thus the map is injective. It is not surjective onto $\ZZ$, because no integer maps to one. Change its codomain to $2\ZZ=\{2k:k\in\ZZ\}$. Every member of this codomain is the image of its defining integer $k$, so the new map is also surjective and hence bijective. Its formula and domain did not change.}
\sol{3.4}{For every input $x$,
\begin{align*}
 x\in f^{-1}(S\cup T)
 &\Longleftrightarrow f(x)\in S\cup T\\
 &\Longleftrightarrow f(x)\in S\ \text{or}\ f(x)\in T\\
 &\Longleftrightarrow x\in f^{-1}(S)\ \text{or}\ x\in f^{-1}(T)\\
 &\Longleftrightarrow x\in f^{-1}(S)\cup f^{-1}(T).
\end{align*}
The two sets have exactly the same members, which proves equality. No injectivity or surjectivity is used.}
\sol{3.5}{For a left inverse, the given maps are $f:A\to B$ and $h:B\to A$, with $h(f(a))=a$ for every $a\in A$. To prove injectivity, choose $x,y\in A$ and assume $f(x)=f(y)$. Applying the same function $h$ to equal objects preserves equality:
\[
 h(f(x))=h(f(y)).
\]
Use the left-inverse identity once at $x$ and once at $y$. The displayed equality becomes $x=y$, as required. Applying $h$ was legitimate because both outputs of $f$ belong to its domain $B$.

For a right inverse, the identity is instead $f(h(b))=b$ for every $b\in B$. To prove surjectivity, choose any target $b\in B$. The element $h(b)$ belongs to $A$ and is a candidate preimage. Its output is exactly $f(h(b))=b$, so the candidate works. The first identity permits recovering original inputs; the second permits reaching arbitrary targets. They justify different conclusions.}
\sol{3.6}{If $z\in f(S\cap T)$, some $x\in S\cap T$ satisfies $f(x)=z$. That same $x$ belongs to both sets, so $z\in f(S)$ and $z\in f(T)$. Therefore $f(S\cap T)\subseteq f(S)\cap f(T)$. For failure of equality, use $f(x)=x^2$, $S=\{1\}$, and $T=\{-1\}$ in the real domain. The intersection of the input sets is empty, so its image is empty; each separate image is $\{1\}$. The reverse direction would have to combine two possibly different preimages into one common preimage, which is not generally possible.}

\section{Week 4}
\sol{4.1}{Write $a=2k$ and $b=2\ell$ with integer witnesses. Then $ab=4k\ell$. Since $k\ell$ is an integer, the product is divisible by four by definition.}
\sol{4.2}{The original claim has assumption ``$3n+1$ is even'' and conclusion ``$n$ is odd.'' Its contrapositive assumes the conclusion fails and must prove the original assumption fails.

If the integer $n$ is not odd, the even/odd classification gives $n=2k$ for an integer $k$. Substitute this into the quantity in the original assumption:
\[
 3n+1=3(2k)+1=6k+1=2(3k)+1.
\]
Since $3k$ is an integer, this is odd and hence not even. Thus ``$n$ is not odd'' implies ``$3n+1$ is not even,'' proving the contrapositive. By its equivalence to the original implication, the required result follows. The temporary assumption was the negation of the conclusion, not the original hypothesis.}
\sol{4.3}{Suppose for contradiction that $r+u=s$ is rational. Then $u=s-r$. For rational numbers $s=a/b$ and $r=c/d$ with nonzero integer denominators,
\[
 s-r=\frac{ad-bc}{bd}
\]
is rational, since its numerator and denominator are integers and $bd\ne0$. Thus $u$ would be rational, contrary to the assumption. Hence $r+u$ is irrational.}
\sol{4.4}{Existence: $x=2$ satisfies $2=3\cdot2-4$. Uniqueness: any solution satisfies $x=3x-4$, so subtracting $x$ gives $0=2x-4$, then $2x=4$ and $x=2$. Every solution equals the exhibited one. Rearranging to find a candidate is not, by itself, a reason to omit its substitution check.}
\sol{4.5}{The missing possibility is the negative root. Since $\sqrt2$ denotes the nonnegative square root, the equation $x^2=2$ becomes
\[
 (x-\sqrt2)(x+\sqrt2)=0.
\]
The zero-product property gives $x=\sqrt2$ or $x=-\sqrt2$, and both square to two. An extra assumption $x\ge0$ would select the first value.}
\sol{4.6}{Because $2\sqrt2$ is nonnegative and has square $4\cdot2=8$, it equals the nonnegative square root $\sqrt8$. If $\sqrt8$ were rational, division by two would make $\sqrt2$ rational. This contradicts Chapter~4. The reduction is complete because it checks both the identity between the roots and preservation of rationality under division by a nonzero integer.}

\section{Week 5}
\sol{5.1}{For $n=0$, the empty sum is zero and $n(n+1)/2=0$. Suppose the formula holds at $k\ge0$. Then
\begin{align*}
 \sum_{j=1}^{k+1}j
 &=\sum_{j=1}^{k}j+(k+1)\\
 &=\frac{k(k+1)}2+(k+1)\\
 &=\frac{k(k+1)+2(k+1)}2
 =\frac{(k+1)(k+2)}2.
\end{align*}
This is exactly the formula at $k+1$. The base and transition establish it for every natural number.}
\sol{5.2}{The proposed formula is $y_n=3\cdot2^n$. At zero it gives three, as required. If it holds at $n$, then $y_{n+1}=2y_n=2(3\cdot2^n)=3\cdot2^{n+1}$. Induction proves the formula. The exponent counts how many doublings have occurred, not the number of terms in a sum.}
\sol{5.3}{The divisions are
\[
 1071=2\cdot462+147,\qquad
 462=3\cdot147+21,\qquad
 147=7\cdot21+0.
\]
The quotients are $2,3,7$, and the remainders are $147,21,0$. Replacing a pair by its divisor and remainder preserves its common divisors. At the end the common divisors are exactly the divisors of $21$, so the greatest common divisor is $21$. The last nonzero remainder, not the last quotient, gives the answer.}
\sol{5.4}{Encode plus and minus as $+1$ and $-1$. The initial product is $+1$. Flipping two entries multiplies the product by $(-1)^2=1$, so it remains $+1$. A configuration with exactly three minus signs has product $(-1)^3=-1$. Since the invariant differs, that configuration is unreachable by any number of allowed moves.}
\sol{5.5}{Subtract the candidate fixed point one:
\[
 x_{n+1}-1=\frac{x_n+1}{2}-1=\frac{x_n-1}{2}.
\]
Induction gives $x_n-1=2^{-n}(x_0-1)$: it holds at zero, and each update halves the preceding expression. For every finite $n$, $2^{-n}$ is positive and hence nonzero. If $x_0-1\ne0$, their product cannot be zero. Thus approaching one does not mean reaching one after finitely many exact updates.}
\sol{5.6}{Starting from $P(0)$ and moving by two reaches $0,2,4,\ldots$ only. It gives no information about odd indices. Add the starting case $P(1)$; the same transition then gives $1,3,5,\ldots$. Every natural number lies in one of these two chains, so both starting cases together cover the full domain.}

\section{Week 6}
\sol{6.1}{The required integers are $-12,-7,-2,3,8$. Each differs from $-2$ by an integer multiple of five. The next value, $13$, lies outside the requested interval. The representative among $0,1,2,3,4$ is $3$, since $3-(-2)=5$.}
\sol{6.2}{An operation on classes is well-defined only if different names for the same input classes give the same output class. Suppose $a'$ names the same class as $a$, and $b'$ the same class as $b$. By the definition of congruence, write
\[
 a'=a+mk,\qquad b'=b+m\ell
\]
for integers $k,\ell$. The output classes are equal exactly when their representatives differ by a multiple of $m$. Compute that difference:
\begin{align*}
 a'b'-ab
 &=(a+mk)(b+m\ell)-ab\\
 &=ab+am\ell+mkb+m^2k\ell-ab\\
 &=m(a\ell+bk+mk\ell).
\end{align*}
The expression in parentheses is an integer, so the difference is divisible by $m$. Therefore $[a'b']_m=[ab]_m$. We have proved equality of classes, not equality of the underlying products, which can be different integers. Both representatives were allowed to change independently.}
\sol{6.3}{Label positions by $0,1,2,3$ modulo four. Rotation by two sends $j$ to $j+2$ modulo four. Applying it twice sends $j$ to $j+4$, the original position. Hence it is its own inverse. It is not the identity because, for example, it sends position $0$ to position $2$. An operation whose square is the identity need not itself be the identity.}
\sol{6.4}{Reflexivity holds because a finite set has the same cardinality as itself. Symmetry holds because $|A|=|B|$ implies $|B|=|A|$. Transitivity holds because $|A|=|B|$ and $|B|=|C|$ imply $|A|=|C|$. Therefore the relation is an equivalence relation. It identifies sets according to size, not according to their actual members.}
\sol{6.5}{The proposed rule takes an input class modulo $m$ and returns a class modulo $d$. Assume $m,d$ are positive integers. Different representatives of the same input class must produce the same output class.

First suppose $d$ divides $m$. Write $m=dt$ for an integer $t$. If $[a]_m=[b]_m$, then $a-b=mk$ for an integer $k$. Substituting $m=dt$ gives $a-b=d(tk)$, so $d$ divides $a-b$. This is exactly $[a]_d=[b]_d$. Thus any two representatives of one input class give the same output class.

Conversely, suppose the rule is well-defined. The integers zero and $m$ represent the same class modulo $m$, because their difference is a multiple of $m$. Their output classes must therefore agree: $[0]_d=[m]_d$. By the definition of equality of congruence classes, $d$ divides $m$. This proves necessity.

For example, passing from classes modulo six to classes modulo three works. Passing from classes modulo six to classes modulo four does not: zero and six are the same input class but give distinct classes modulo four. The failing example identifies precisely what well-definedness is checking.}
\sol{6.6}{Use the family $\{\{1\},\{2\}\}$ ordered by inclusion. Each member is maximal because neither is properly contained in another member of the family. Neither contains the other, so there is no maximum. Add $\{1,2\}$. It contains both previous members and itself, so it is now the maximum. The ambient universe does not determine which candidates are present in the ordered family.}


\chapter{Solutions: counting, assignment, and approximation}\label{app:solb}
\section{Week 7}
\sol{7.1}{The edges are $ab,bc,ca,cd$. The degrees at $a,b,c,d$ are respectively $2,2,3,1$. Their sum is eight, twice the four edges. Exactly two vertices, $c$ and $d$, have odd degree. The edge $cd$ contributes once to the degree of $c$ and once to that of $d$; this is the same double-counting mechanism as the general theorem.}
\sol{7.2}{There are four possible remainder classes, represented by $0,1,2,3$. Place each of the five integers in its remainder class. If no two shared a class, each class would contain at most one of the five integers, allowing at most four integers altogether. This contradiction proves a shared remainder. The inputs need not be positive or distinct; repeated integers already give a pair in the same class.}
\sol{7.3}{Use the directed cycle $a\to b\to c\to a$. Every vertex has an outgoing edge, so none is a sink. There are no loops or opposite pairs, so this is an oriented graph. One vertex cannot have an outgoing edge in a loopless one-vertex graph. With two vertices, at most one edge is allowed, and its endpoint is a sink; with no edges both are sinks. Thus three is the smallest possible order.}
\sol{7.4}{At $a$, the first neighborhood is $\{b,c\}$ and the second is $\{d\}$. The two routes through $b$ and $c$ reach the same endpoint, counted once. At $b$, the first neighborhood is $\{d\}$ and the second is empty. At $c$ the same holds. At $d$ both are empty. Thus $a,b,c$ fail the inequality $|N_2^+(v)|\ge|N_1^+(v)|$, whereas the sink $d$ satisfies it as $0\ge0$.}
\sol{7.5}{If every degree exceeded $2m/n$, summing over the $n$ vertices would give a degree sum strictly larger than $n(2m/n)=2m$. This contradicts the handshake identity. Hence at least one degree is at most the average. The conclusion is existential. In a star, one central vertex is connected to $n-1$ other vertices, called leaves, with no edges between the leaves. The center has degree $n-1$ and every leaf has degree one; for large $n$, neither a maximum-degree estimate near the average nor equal degrees follows.}
\sol{7.6}{Take the transitive tournament with $a\to b$, $a\to c$, and $b\to c$. Vertex $a$ reaches both other vertices in one step, so it has the required reachability property. But its first neighborhood has size two and its exact second neighborhood is empty: $c$ is excluded because it is already a first neighbor. Thus this reachability property alone does not imply the second-neighborhood inequality for that vertex.}

\section{Week 8}
\sol{8.1}{Let Ada accept tasks $1,2$ and Ben accept only task $1$. The matching consisting of Ada--$1$ is maximal: Ben--$1$ shares the used task and Ada--$2$ shares the used student, so neither can simply be added. It is not maximum because the matching Ada--$2$, Ben--$1$ has size two. Maximal means no edge can be added without revising the current choices; maximum compares against every possible matching.}
\sol{8.2}{The unmatched student is $c$. Use the alternating path
\[
 c\ --\ 2\ --\ a\ --\ 1\ --\ b\ --\ 3.
\]
The edges $c$--$2$, $a$--$1$, and $b$--$3$ are outside the matching; $a$--$2$ and $b$--$1$ are inside it. Task $3$ is unmatched. Flipping the path gives $c$--$2$, $a$--$1$, $b$--$3$, a matching of size three. Each internal vertex exchanges one matched edge for another, while the endpoints become matched.}
\sol{8.3}{Let $S$ be all three students. Their neighbor set is $N(S)=\{1,2\}$, so $|N(S)|=2<3=|S|$. Distinct assignments of all three would require three distinct members of a two-element set, which is impossible. Counting each student's two options separately counts the same tasks repeatedly; feasibility depends on the number of distinct options available to the group.}
\sol{8.4}{The flip proof treats each internal vertex as incident to exactly two path edges, one matched and one unmatched. Repeated visits could invalidate that local count if one flipped an arbitrary walk without first simplifying it. A shortest directed route in the alternating-search graph has no repeated vertex: removing the segment between two visits yields a shorter directed route with the same endpoints. Directions in this search graph encode the allowed matched/unmatched moves, so the shortened route is still alternating. Equivalently, the breadth-first predecessor tree supplies a simple path directly.}
\sol{8.5}{A matching uses each left vertex at most once, so its size is at most $|L|$. Starting with size $k$, and increasing it by one at each augmentation, allows at most $|L|-k$ augmentations. From the empty matching the bound is $|L|$. Without a finite upper bound, a quantity increasing by one can increase indefinitely; strict progress alone is not a termination proof.}
\sol{8.6}{The search starts at one unmatched left vertex $u$; other matched and unmatched vertices are explored by the rules in Chapter~8. Every reached right vertex was reached from a reached left vertex along an edge, so $B\subseteq N(A)$. For the converse, let $a\in A$ and let $b$ be any neighbor of $a$. If $ab$ is unmatched, the search explores that left-to-right edge and reaches $b$. If $ab$ is matched, then $a\ne u$. The only way the search reaches this other left vertex is through its unique matched right partner, which must be $b$. Thus $b$ was already reached. In either case $b\in B$, proving $N(A)\subseteq B$. Exhaustion of the search is needed to say all eligible unmatched edges have been explored.}

\section{Week 9}
\sol{9.1}{Let $N(a)=N(b)=\{1\}$ and $N(c)=\{2,3\}$. Every student has a neighbor, and the full left set has all three tasks as neighbors. Nevertheless $S=\{a,b\}$ has only one neighbor. No assignment can give both $a$ and $b$ distinct acceptable tasks. Checking only individual students and the whole set misses this intermediate obstruction.}
\sol{9.2}{The complete check is
\begin{center}
\begin{tabular}{@{}lll@{}}\toprule
$S$ & $N(S)$ & $|N(S)|\ge|S|$\\\midrule
$\varnothing$ & $\varnothing$ & $0\ge0$\\
$\{a\}$ & $\{1\}$ & $1\ge1$\\
$\{b\}$ & $\{1,2\}$ & $2\ge1$\\
$\{c\}$ & $\{2,3\}$ & $2\ge1$\\
$\{a,b\}$ & $\{1,2\}$ & $2\ge2$\\
$\{a,c\}$ & $\{1,2,3\}$ & $3\ge2$\\
$\{b,c\}$ & $\{1,2,3\}$ & $3\ge2$\\
$\{a,b,c\}$ & $\{1,2,3\}$ & $3\ge3$\\\bottomrule
\end{tabular}
\end{center}
Both $\{a\}$ and $\{a,b\}$ are nonempty proper tight subsets. A matching is $a$--$1$, $b$--$2$, $c$--$3$. The table verifies the hypothesis; the matching itself is a shorter certificate of feasibility for this particular instance.}
\sol{9.3}{Every member of $N(S)\cup N(T)$ either belongs to $N(S)$ or belongs to $N(T)$ without belonging to $N(S)$. Thus
\[
 N(S)\cup N(T)=N(S)\ \mathbin{\dot\cup}\ (N(T)\setminus N(S)),
\]
where the dotted union indicates disjointness. Counting a disjoint union adds its sizes. We did not assume that $N(S)$ and $N(T)$ themselves are disjoint; subtracting the overlap is what creates a disjoint decomposition.}
\sol{9.4}{The difference $|N(S)|-|S|$ is an integer. If positive, it is at least one, giving $|N(S)|\ge|S|+1$. For arbitrary reals, positivity allows values strictly between zero and one. For example, $1.2>1$ does not imply $1.2\ge2$. The one-unit slack comes from integer counting.}
\sol{9.5}{We will apply Hall's theorem, so we must establish its condition for an arbitrary set $S$ of students. If $S$ is empty, the condition reads $0\ge0$. The same counting argument below includes that case.

Count the edges with left endpoint in $S$. Each student in $S$ accepts exactly $d$ tasks and therefore contributes exactly $d$ such edges. No edge has two left endpoints, so this counts each edge once. The total is $d|S|$.

Every right endpoint of a counted edge lies in $N(S)$. Each task in that set is acceptable to at most $d$ students in the whole graph, and hence to at most $d$ members of $S$. Counting the same edges by their right endpoints therefore gives at most $d|N(S)|$. Combining the two counts,
\[
 d|S|\le d|N(S)|.
\]
Divide by the positive number $d$ to obtain $|S|\le|N(S)|$. This holds for every subset of the finite left set. The graph is bipartite by the assignment model, so all the hypotheses of the finite Hall theorem are met. It supplies a matching saturating the students. If $d$ were zero, division and the conclusion would both require reconsideration.}
\sol{9.6}{A model portfolio can use the feasible graph of Exercise~9.2 with assignment $a$--$1$, $b$--$2$, $c$--$3$. Its comparison graph is the one in Exercise~9.1, with shortage certificate $S=\{a,b\}$ and $N(S)=\{1\}$.

Here is a complete constructive sufficiency proof in the general finite setting. Assume $|N(S)|\ge|S|$ for every $S\subseteq L$ and start from the empty matching. If the current matching leaves a student $u$ unmatched, run the alternating search from $u$. If it reaches an unmatched task, take a simple alternating path to it and flip the path. This increases the matching size by one. Otherwise let $A,B$ be the reached left and right sets. The search rules give $N(A)=B$, as proved in Exercise~8.6. Every right vertex in $B$ is matched, since otherwise the search would have found an augmenting endpoint. Its matched left partner is reached. Conversely, every reached left vertex other than $u$ is reached through its matched partner. These pairings form a bijection between $B$ and $A\setminus\{u\}$, so $|N(A)|=|B|=|A|-1$, contradicting Hall's condition. Therefore a nonsaturating matching always admits an augmentation. Size can increase at most $|L|$ times, so the process ends with all students matched.

Finally, maximality forbids adding one edge while leaving all present choices intact. Maximum size allows earlier choices to be revised. Ada--$1$ in Exercise~8.1 demonstrates the distinction. The portfolio should state which definition is used whenever it calls a matching ``optimal.''}

\section{Week 10}
\sol{10.1}{Matrix multiplication gives
\[
 \begin{pmatrix}2&1\\0&3\end{pmatrix}
 \begin{pmatrix}1\\-2\end{pmatrix}
 =\begin{pmatrix}2\cdot1+1(-2)\\0\cdot1+3(-2)\end{pmatrix}
 =\begin{pmatrix}0\\-6\end{pmatrix}.
\]
Equivalently, take one copy of the first column $(2,0)$ and minus two copies of the second column $(1,3)$. Their sum is $(0,-6)$. This displays how columns encode the action on basis vectors.}
\sol{10.2}{The kernel consists of pairs satisfying $x+y=0$. Solving for $y$ gives $(x,y)=(t,-t)$ for arbitrary real $t$. For output five, solve $x+y=5$ to obtain $(t,5-t)$. One solution is $(0,5)$, and every other solution differs from it by $(t,-t)$ in the kernel. The map loses one degree of information: it remembers the sum but not the split between coordinates.}
\sol{10.3}{For $T(x,y)=(x+y,x-y)$, define $U(a,b)=((a+b)/2,(a-b)/2)$. Then
\[
 U(T(x,y))=\left(\frac{x+y+x-y}{2},\frac{x+y-x+y}{2}\right)=(x,y).
\]
Conversely,
\[
 T(U(a,b))=\left(\frac{a+b+a-b}{2},\frac{a+b-a+b}{2}\right)=(a,b).
\]
Both compositions are identities, so $U$ is an inverse, not merely a map that works for one selected input.}
\sol{10.4}{For $u=(1,2)$ and $v=(2,-1)$, the dot product is $1\cdot2+2(-1)=0$. Both squared lengths are $1+4=5$, so both lengths are $\sqrt5$. Cauchy--Schwarz reads $0\le\sqrt5\sqrt5=5$. Orthogonality does not mean either vector is zero.}
\sol{10.5}{Start from the squared length because its dot product expands using addition:
\begin{align*}
 \norm{u+v}_2^2
 &=(u+v)\cdot(u+v)\\
 &=u\cdot u+u\cdot v+v\cdot u+v\cdot v\\
 &=\norm{u}_2^2+2(u\cdot v)+\norm{v}_2^2\\
 &\le\norm{u}_2^2+2|u\cdot v|+\norm{v}_2^2\\
 &\le\norm{u}_2^2+2\norm{u}_2\norm{v}_2+\norm{v}_2^2\\
 &=(\norm{u}_2+\norm{v}_2)^2.
\end{align*}
The first inequality uses that a real number is at most its absolute value. The next applies Cauchy--Schwarz. The final equality recognizes the expansion of a square.

Both $\norm{u+v}_2$ and $\norm{u}_2+\norm{v}_2$ are nonnegative. Taking their nonnegative square roots therefore preserves the comparison and yields $\norm{u+v}_2\le\norm{u}_2+\norm{v}_2$. Without nonnegativity, a comparison of squares would not by itself compare the unsquared quantities in the same order.}
\sol{10.6}{Assume $a_1u_1+\cdots+a_ku_k=0$. To recover one coefficient, choose its index $j$ and take the dot product of both sides with $u_j$. Linearity gives
\[
 a_1(u_1\cdot u_j)+\cdots+a_k(u_k\cdot u_j)=0\cdot u_j=0.
\]
For every $i\ne j$, orthogonality means $u_i\cdot u_j=0$. All of those terms disappear. The only remaining equation is
\[
 a_j(u_j\cdot u_j)=a_j\norm{u_j}_2^2=0.
\]
The vector $u_j$ is nonzero, so its squared length is positive. Division by that squared length gives $a_j=0$. Because $j$ was an arbitrary index, the argument recovers every coefficient. This is linear independence, not just a statement about one particular coefficient choice.

Nonzeroness matters. If one of the vectors were zero, a nonzero coefficient multiplying it could appear in a zero combination. The proof would then fail precisely at the division step.}

\section{Week 11}
\sol{11.1}{Name the input errors $\Delta x=x-a$ and $\Delta y=y-b$, where the symbol $\Delta$ is simply a label for a difference. Subtraction and regrouping give
\[
 (2x-y)-(2a-b)=2(x-a)-(y-b)=2\Delta x-\Delta y.
\]
Now use the triangle inequality and the absolute-value scaling rule:
\begin{align*}
 |2\Delta x-\Delta y|
 &\le|2\Delta x|+|{-\Delta y}|\\
 &=2|\Delta x|+|\Delta y|\\
 &\le2(0.02)+0.03=0.07.
\end{align*}
The signs of the input errors were not specified, so a guaranteed estimate must allow the worst permitted combination. For $\Delta x=0.02$ and $\Delta y=-0.03$, the output error is $0.04+0.03=0.07$, attaining the bound. For $\Delta x=0.02$ and $\Delta y=0.03$, it is $0.04-0.03=0.01$. In this expression, same-sign errors can partially cancel, while opposite-sign errors reinforce because the second error is subtracted.}
\sol{11.2}{For $a_n=3/(n+1)$ with $n\ge0$, the error from zero is $3/(n+1)$. Given $\eps>0$, choose a positive integer $N>3/\eps$. Then $n\ge N$ implies $n+1>N$ and hence $|a_n-0|=3/(n+1)<3/N<\eps$. The same $N$ works for every later index, which is required for convergence. This choice is sufficient, even though it need not be the smallest possible threshold.}
\sol{11.3}{For positive coordinate dimension $d$, let $d_\infty(x,y)=\max_{1\le i\le d}|x_i-y_i|$. It is nonnegative. It equals zero exactly when every coordinate difference equals zero, hence exactly when $x=y$. Symmetry follows from $|x_i-y_i|=|y_i-x_i|$. For the triangle inequality, each coordinate satisfies
\[
 |x_i-z_i|\le|x_i-y_i|+|y_i-z_i|
 \le d_\infty(x,y)+d_\infty(y,z).
\]
Taking the maximum over $i$ preserves this common bound. Thus it is a metric.}
\sol{11.4}{For any $m>n$, insert the consecutive iterates between the endpoints. Repeated triangle inequalities give
\[
 d(x_m,x_n)\le d(x_m,x_{m-1})+\cdots+d(x_{n+1},x_n)
 \le\sum_{k=n}^{m-1}2(1/3)^k.
\]
The geometric-sum bound from the chapter yields
\[
 \sum_{k=n}^{m-1}2(1/3)^k
 \le\frac{2(1/3)^n}{1-1/3}
 =3(1/3)^n=3^{1-n}.
\]
This bound does not depend on $m$, so it controls all later endpoints together. For any positive tolerance, choose $N$ so $3^{1-N}$ is smaller than that tolerance. When both indices are at least $N$, put the smaller one in the role of $n$; if they are equal their distance is zero. This proves the Cauchy condition. Completeness would be an additional requirement to conclude convergence inside the chosen space.}
\sol{11.5}{The sequence $x_n=2+2^{-(n+1)}$ lies in $(2,3)$ for all $n\ge0$. Its real limit is two because its distance from two is $2^{-(n+1)}\to0$. A convergent real sequence is Cauchy: two sufficiently late terms are each within $\eps/2$ of the limit and hence within $\eps$ of each other. If this sequence converged to a point of $(2,3)$, uniqueness of real limits would force that point to be two, which is outside the interval. Thus this is a Cauchy sequence without a limit in the space.}
\sol{11.6}{One is an upper bound for $(0,1)$. No $b<1$ is an upper bound: if $b<0$, the point $1/2$ exceeds it; if $0\le b<1$, the point $(b+1)/2$ lies in $(0,1)$ and exceeds $b$. Therefore one is the least upper bound. It is not a maximum because it is not a member. Indeed every $x\in(0,1)$ is exceeded by $(x+1)/2$, another member of the interval.}

\section{Week 12}
\sol{12.1}{For $T(x)=(3x+1)/4$, subtraction gives $|T(x)-T(y)|=(3/4)|x-y|$. Thus $q=3/4$ is a contraction factor on the complete space $\RR$. The fixed-point equation gives $4x=3x+1$, so the fixed point is one. Starting at zero, subtract one at each step to get $x_{n+1}-1=(3/4)(x_n-1)$, hence $x_n=1-(3/4)^n$. The formula also verifies the starting value and recurrence directly.}
\sol{12.2}{The residual certificate is
\[
 d(z,x_*)\le\frac{d(z,Tz)}{1-q}
 \le\frac{0.002}{1-0.75}=0.008.
\]
The residual is not itself a bound on the true error unless the factor $1/(1-q)$ is included or another argument justifies a stronger estimate.}
\sol{12.3}{For $T(x)=(x^2+1)/3$ on $[1/3,2/3]$, the squared input lies in $[1/9,4/9]$, so the image lies in $[10/27,13/27]$. Both endpoints lie between $1/3=9/27$ and $2/3=18/27$, proving the self-map condition. For inputs in this smaller interval,
\[
 |T(x)-T(y)|=\frac{|x-y||x+y|}{3}
 \le\frac49|x-y|.
\]
Thus a contraction factor is $4/9$, smaller than $2/3$ on $[0,1]$. Bounds using the smaller factor apply after the approximate point and the required iterations are known to lie in this smaller domain. Merely writing a smaller interval does not prove that membership.}
\sol{12.4}{For $T(x)=x/2$ on $(0,1)$, the map is a self-map with factor $1/2$, but the space is incomplete and its real fixed point zero lies outside it. For the identity on $\RR$, the factor one is not strictly below one and every point is fixed, so uniqueness fails. For $T(x)=x+1$ on $\RR$, the same factor-one boundary admits no fixed point. For $T(x)=x/2+1$ restricted to $[0,1]$, the contraction estimate holds but the map is not a self-map: $T(1)=3/2$. These examples violate different hypotheses and should not be described as counterexamples to the theorem itself.}
\sol{12.5}{The fixed-point equation $x=ax+b$ gives $x_*=b/(1-a)$; the denominator is nonzero because $|a|<1$. The map has contraction factor $|a|$ on $\RR$. Subtracting the fixed-point equation from the recurrence yields $x_{n+1}-x_*=a(x_n-x_*)$, hence $x_n-x_*=a^n(x_0-x_*)$. If $a<0$ and the starting error is nonzero, its sign alternates at each step while its magnitude decreases by the factor $|a|$. If the starting error is zero, all iterates are already fixed.}
\sol{12.6}{The contraction theorem supplies both fixed points, because the two maps are self-maps of the same nonempty complete space and have contraction factors at most $q<1$. Their identities are $T(x_T)=x_T$ and $S(x_S)=x_S$.

We need to compare points obtained from different maps, while the contraction bound compares outputs of the same map. Insert the intermediate point $T(x_S)$ to separate those two issues:
\begin{align*}
 d(x_T,x_S)
 &=d(T(x_T),S(x_S))\\
 &\le d(T(x_T),T(x_S))+d(T(x_S),S(x_S))\\
 &\le q\,d(x_T,x_S)+\eta.
\end{align*}
The first term uses the contraction property of $T$. The second uses the assumed discrepancy bound at the allowed input $x_S$. Writing $D=d(x_T,x_S)$, the result is $D\le qD+\eta$. Subtract $qD$ to get $(1-q)D\le\eta$, then divide by $1-q>0$. Thus $d(x_T,x_S)\le\eta/(1-q)$.

The proof explains why a uniform comparison of maps can control their fixed points: one inserted point isolates the change of input from the change of map.}


\chapter{Solutions: probability, formalization, and investigation}\label{app:solc}
\section{Week 13}
\sol{13.1}{The expectation is $0(1/2)+2(1/3)+5(1/6)=2/3+5/6=3/2$. The possible values two and five are each at least $3/2$. The expectation is a weighted average, not necessarily a possible observed value.}
\sol{13.2}{Two labelings place all three vertices on the same side, giving cut size zero. The other six labelings split one vertex from the other two, giving cut size two. The average over all eight equally likely labelings is $(2\cdot0+6\cdot2)/8=3/2$, half of the three edges. Cut size one or three never occurs for a triangle.}
\sol{13.3}{Place vertex $1$ on the left by the tie rule. Vertex $2$ has one assigned neighbor on the left, so place it on the right. Vertex $3$ has its assigned neighbor $2$ on the right, so place it on the left. Vertex $4$ has both assigned neighbors, $1$ and $3$, on the left, so place it on the right. Every cycle edge crosses, giving a cut of size four. The general guarantee is only at least two, but the rule can do better on a particular instance.}
\sol{13.4}{Let $A_1,\ldots,A_4$ be the bad events. Their union is the event that at least one bad event occurs. The union bound gives
\[
 \Prob(A_1\cup A_2\cup A_3\cup A_4)
 \le\sum_{j=1}^{4}\Prob(A_j)\le4(0.2)=0.8.
\]
The complementary event, in which all four are avoided, therefore has probability at least $1-0.8=0.2$. Its probability is positive, so it contains at least one outcome. This is an existence proof, not a claim that every outcome avoids the bad events.

With the bound $0.25$ instead, the same calculation only gives avoidance probability at least $1-4(0.25)=0$. Every event already has nonnegative probability, so this supplies no existence information. To show that the conclusion can genuinely fail, use a uniform four-outcome space and let each bad event contain one different outcome. Each has probability $1/4$, but their union is the whole space. No outcome avoids them all.}
\sol{13.5}{Let $X=0$ with probability $1/2$ and $X=10$ with probability $1/2$. Its expectation is five, but one possible value is zero. For a finite probability space, at least one outcome of positive probability has value at least the expectation. Otherwise every positive-weight term would lie strictly below that average, and their weighted average would also lie strictly below it, a contradiction. The argument ensures existence, not uniform performance.}
\sol{13.6}{Markov's inequality gives $\Prob(X\ge12)\le\E X/12=1/4$. Equality is attained by $X=12$ with probability $1/4$ and $X=0$ with probability $3/4$. The expectation is three and the event $X\ge12$ has probability exactly $1/4$. Nonnegativity is part of the inequality's hypothesis.}

\section{Week 14}
\sol{14.1}{For $a=(1,1,-1)$, the nonzero aperiodic shifts give
\[
 A_1=1\cdot1+1(-1)=0,\qquad A_2=1(-1)=-1.
\]
The nonzero periodic correlations are
\[
 C_1=1\cdot1+1(-1)+(-1)1=-1,
 \qquad C_2=1(-1)+1\cdot1+(-1)1=-1.
\]
Both also follow from $C_1=A_1+A_2$ and $C_2=A_2+A_1$. The zero periodic correlation is three, although it was not among the requested nonzero shifts.}
\sol{14.2}{Every correlation summand is a product of two entries. Negating all entries changes $a_j a_k$ to $(-a_j)(-a_k)=a_j a_k$. Each summand is unchanged, with or without wrapped indices, so the entire aperiodic and periodic correlations are unchanged. This symmetry preserves the conditions and can justify identifying paired examples, provided counts are then explicitly reported modulo that symmetry.}
\sol{14.3}{With the convention $H_{ij}=a_{j-i\bmod4}$, the first two rows are $(1,1,1,-1)$ and $(-1,1,1,1)$. Their dot product is $-1+1+1-1=0$, an off-diagonal entry of $HH^{\mathsf T}$. The dot product of the first row with itself is $1+1+1+1=4$, a diagonal entry. These two checks illustrate the construction; the periodic-correlation argument in the chapter establishes all entries together.}
\sol{14.4}{Let $1\le s<n$. Periodic indexing reduces indices modulo $n$, so a shift of $n-s$ is the same as a shift of $-s$. Consequently,
\[
 C_{n-s}=\sum_{j=0}^{n-1}a_j a_{(j-s)\bmod n}.
\]
Set $k=(j-s)\bmod n$. As $j$ runs through every residue, so does $k$, and conversely $j=(k+s)\bmod n$. Substitute both occurrences of the old index:
\begin{align*}
 C_{n-s}
 &=\sum_{k=0}^{n-1}a_{(k+s)\bmod n}a_k\\
 &=\sum_{k=0}^{n-1}a_k a_{(k+s)\bmod n}\\
 &=C_s.
\end{align*}
The first equality reindexes all terms without losing any; the second reverses the real factors in each product; the last recognizes the definition. The index range is empty when $n=1$, so there is then no nonzero-shift assertion to check. Real multiplication is relevant: a correlation defined using complex conjugation would require checking its own convention rather than copying the final equality blindly.}
\sol{14.5}{The order must be a square, and if it exceeds two it must be divisible by four. Eight is not a square. Nine and twenty-five are squares but not divisible by four. Sixteen and thirty-six satisfy both necessary conditions and remain candidates. Neither condition constructs a matrix at these orders. In particular, surviving a necessary-condition test is not evidence of existence on its own.}
\sol{14.6}{One sequence with all required correlations verified is a certificate of existence at its length. A correctly implemented exhaustive search of all $2^n$ sign sequences, with no qualifying result, establishes nonexistence at that fixed length, subject to the correctness of the search and arithmetic. It says nothing by itself about an untested length. A conclusion for all lengths requires a proof applying to arbitrary $n$, or a proved reduction showing that every possible example would occur among the finitely many cases searched.}

\section{Week 15}
The following scripts use the definitions in the companion namespace. They are intended to be read alongside their ordinary mathematical arguments. The reproducibility appendix distinguishes the reviewed source from an actually executed Lean check.
\sol{15.1}{The expressions $P,Q$ are propositions. The expression $P\to Q$ is the proposition that a proof of $P$ can be transformed into a proof of $Q$. The name $h$ denotes such a transformation. When \code{hp : P}, the application \code{h hp} has type \code{Q}; it supplies a proof of the conclusion from the given proof of the premise.}
\sol{15.2}{A conjunction already contains a proof of each component:
\par\noindent\begin{minipage}{\linewidth}
\begin{lstlisting}
theorem and_left (P Q : Prop) : And P Q -> P := by
  intro h
  exact h.left
\end{lstlisting}
\end{minipage}\par

The only inference is to extract the first component from the assumed pair of proofs. No claim about the truth of $P$ without that hypothesis is made.}
\sol{15.3}{Supply five and check its property:
\par\noindent\begin{minipage}{\linewidth}
\begin{lstlisting}
theorem witness_five : Exists (fun n : Nat => n + 2 = 7) := by
  exact Exists.intro 5 rfl
\end{lstlisting}
\end{minipage}\par

The equality reduces to the numerical identity $5+2=7$. The witness's type is natural numbers, not an unspecified real or integer domain.}
\sol{15.4}{Applying a function to equal objects preserves equality, but equality of outputs need not imply equality of inputs. From \code{h : g x = g y}, applying \code{congrArg g h} would, when the types allow it, give equality of another pair of outputs after applying $g$ again. It would not cancel $g$. The missing property is injectivity of $g$, which explicitly guarantees that equal outputs have equal inputs. A constant map supplies a counterexample to unrestricted cancellation.}
\sol{15.5}{The mathematical proof applies $g$ to an equality after $f$, then invokes injectivity of the composition:
\par\noindent\begin{minipage}{\linewidth}
\begin{lstlisting}
theorem inj_of_comp {A B C : Type}
    (f : A -> B) (g : B -> C)
    (hcomp : Inj (fun x => g (f x))) : Inj f := by
  intro x y hxy
  exact hcomp x y (congrArg g hxy)
\end{lstlisting}
\end{minipage}\par

The hypothesis concerns only the composition. There is no need to assume $g$ is injective outside the image of $f$.}
\sol{15.6}{Natural-number subtraction is truncated at zero, so \code{(7 - 10 : Nat)} computes to zero. Integer subtraction allows negative results, so \code{(7 - 10 : Int)} computes to minus three. A change of domain can change the operations themselves, not just the collection of candidate values. Translating an integer formula into naturals without checking this distinction can change the theorem even when the displayed symbols look similar.}

\section{Week 16}
\sol{16.1}{After unpacking \code{hg c}, the witness is \code{b : B} and \code{hb} proves \code{g b = c}. After unpacking \code{hf b}, the witness is \code{a : A} and \code{ha} proves \code{f a = b}. The final witness must belong to $A$ because the composition starts there. Rewriting the final expression $g(f(a))$ first by \code{ha} and then by \code{hb} gives $c$.}
\sol{16.2}{Represent union by disjunction of membership predicates:
\par\noindent\begin{minipage}{\linewidth}
\begin{lstlisting}
def join {A : Type} (s t : A -> Prop) : A -> Prop :=
  fun x => Or (s x) (t x)

theorem preimage_join {A B : Type} (f : A -> B)
    (s t : B -> Prop) (x : A) :
    Iff (preimage f (join s t) x)
      (join (preimage f s) (preimage f t) x) := by
  rfl
\end{lstlisting}
\end{minipage}\par

Both sides unfold to the proposition that $s(f(x))$ or $t(f(x))$ holds. This is the membership proof from Exercise~3.4 expressed through definitions.}
\sol{16.3}{The first and last components of the nested conjunction supply the two claims:
\par\noindent\begin{minipage}{\linewidth}
\begin{lstlisting}
theorem second_not_self {V : Type}
    (r : V -> V -> Prop) (v w : V)
    (h : Second r v w) : Not (w = v) := by
  exact h.left

theorem second_route {V : Type}
    (r : V -> V -> Prop) (v w : V)
    (h : Second r v w) :
    Exists (fun u : V => And (r v u) (r u w)) := by
  exact h.right.right
\end{lstlisting}
\end{minipage}\par

The projections reflect the grouping in the definition. Changing that grouping requires changing the projections, though not the underlying mathematical meaning.}
\sol{16.4}{The new theorem assumes an impossible proposition as an additional input. It proves only that the conclusion follows \emph{if} that input were available, using the logical rule that falsehood implies every proposition. It does not provide a proof of the original theorem, whose hypothesis list did not include \code{False}. A statement diff exposes this change even if the proof contains no unfinished placeholder.}
\sol{16.5}{First compare theorem signatures with the intended mathematical statement, including the types of every variable; this detects a changed domain. Next unfold the definitions used in the conclusion and compare their membership conditions; this detects a missing exclusion such as the first-neighbor condition. Finally run the checker, inspect completion warnings and proof dependencies, and investigate any placeholder axiom; this detects an unfinished proof. All three checks are needed because a correct proof of the wrong definition can have no placeholder.}
\sol{16.6}{A Python enumeration can check all sign sequences at length four using its particular code and arithmetic. A Lean proof of a finite proposition might establish that a specified four-entry sequence satisfies an encoded correlation condition. A universal Lean theorem might quantify over every natural length and characterize all lengths admitting such a sequence. These are different scopes even when all programs finish without an error. For any of them, the encoded conditions must still match the intended correlation convention. A finite proposition in Lean does not become universal merely because it was formally proved.}

\section{Week 17}
\sol{17.1}{The constant sequence $x_n=c$ has infinitely many indices inside any positive-radius ball centered at $c$, but only one distinct value. The proof needs infinitely many indices so that it can choose increasing indices. It does not need distinct points; a repeated point is entirely compatible with a convergent subsequence and may make convergence simpler.}
\sol{17.2}{An arbitrary choice could repeat an earlier index or place the selected indices out of order, so the resulting list might not be a subsequence. After choosing $n_{k-1}$, choose $n_k$ from $S_k\setminus\{0,1,\ldots,n_{k-1}\}$. The excluded set is finite while $S_k$ is infinite, so a permissible index remains. This guarantees $n_k>n_{k-1}$ without losing membership in $S_k$.}
\sol{17.3}{For a Cauchy sequence in the discrete metric, apply the Cauchy condition with tolerance $1/2$. Beyond some index, every pair of terms must have distance less than $1/2$, hence distance zero. The sequence is therefore eventually constant and converges in $\NN$, proving completeness. Balls of radius $1/2$ are singletons, so no finite family covers $\NN$; total boundedness fails. The sequence $x_n=n$ has no convergent subsequence because every pair of distinct selected terms remains distance one apart. Completeness controls Cauchy sequences; it does not say every sequence has a Cauchy subsequence.}
\sol{17.4}{Total boundedness asks for a finite cover by balls at \emph{each requested radius}. Fix $r>0$. Since $b-a>0$, choose a positive integer $N$ satisfying $N>(b-a)/r$. This implies the spacing $h=(b-a)/N$ satisfies $0<h<r$.

Use the $N+1$ centers
\[
 c_k=a+kh,\qquad k=0,1,\ldots,N.
\]
The first center is $a$, the last is $b$, and all the others lie between them. If $x=b$, it is already a center. Otherwise, $x\in[a,b)$ lies in one of the finitely many consecutive intervals $[c_k,c_{k+1}]$. Its distance to $c_k$ is at most $h<r$, so it belongs to the radius-$r$ ball about that center. Thus every point is covered.

The list is finite and all its centers belong to the space, as the definition requires. The integer $N$, and therefore the number and locations of centers, may depend on $a,b,r$. Once those are fixed, the same list covers every $x\in[a,b]$; we did not choose a new list after seeing each point.}
\sol{17.5}{The route $v\to u\to w$ makes $w$ two-edge reachable from $v$. However, the additional edge $v\to w$ makes its shortest directed distance one, so $w$ is excluded from the exact second neighborhood. In this graph $N_1^+(v)=\{u,w\}$ and $N_2^+(v)=\varnothing$. Replacing exact second-neighbor membership by mere existence of a two-edge route changes the counted set and can change the truth of an inequality involving its size. It is a semantic change, not a harmless syntax simplification.}
\sol{17.6}{For contraction, a suitable request is: ``Using my knowledge of triangle inequalities and induction, teach the Cauchy condition, completeness, and the geometric-tail estimate needed to prove that repeated applications of this specified contraction reach its unique fixed point. After each concept, show its use in that proof.'' For the subsequence theorem, ask: ``Using the metrics and completeness I already know, explain finite small-ball covers, infinite pigeonhole selection, and increasing subsequence indices, then build the nested-selection argument.'' The first needs a contraction factor and iterative geometric error propagation; the second needs total boundedness and finite-cover extraction. Neither requires the entire collection of concepts in the other.

Without a completed repository build, the reading record may report the statements and definitions read, the source identifiers inspected, the dependency roles identified, and any proofs reconstructed independently. It may not report that the repository was compiled, that its validation procedure passed, or that all imported proofs were verified. Those are separate empirical and mathematical claims.}

\section{Week 18}
\sol{18.1}{The first three bounds are
\begin{align*}
 e_1&\le qe_0+\eps_0,\\
 e_2&\le q^2e_0+q\eps_0+\eps_1,\\
 e_3&\le q^3e_0+q^2\eps_0+q\eps_1+\eps_2.
\end{align*}
At step $n$, the earliest error $\eps_0$ has weight $q^{n-1}$ and the most recent error $\eps_{n-1}$ has weight one. Earlier errors have passed through more contraction steps. Writing out these cases is a practical check on the exponent $n-1-j$ in the general formula.}
\sol{18.2}{At $n=0$, the formula $y_n=\eta(1-q^n)/(1-q)$ gives zero. If it holds at $n$, then
\[
 y_{n+1}=q\eta\frac{1-q^n}{1-q}+\eta
 =\eta\frac{q-q^{n+1}+1-q}{1-q}
 =\eta\frac{1-q^{n+1}}{1-q}.
\]
Induction proves it. Since $q^n\to0$, the distance from the true fixed point zero tends to $\eta/(1-q)$. For $\eta>0$, any universally smaller multiplicative factor would fail on this example for sufficiently large $n$. Thus the asymptotic factor is sharp under these assumptions.}
\sol{18.3}{The space $\RR$ is nonempty and complete, and $T(x)=0$ is a self-map. It satisfies $|T(x)-T(y)|=0\le q|x-y|$ for any chosen $q\in(0,1)$. The fixed point is zero. For $y_n=(-1)^n\eta$ with $\eta>0$, the update discrepancy is $|y_{n+1}-T(y_n)|=\eta$, so the uniform-error hypothesis holds. Nevertheless even-index terms equal $\eta$ and odd-index terms equal $-\eta$. Their distance is $2\eta$, so the sequence is not Cauchy and cannot converge. A bound on all distances to the fixed point does not force those distances, or the sequence itself, to converge.}
\sol{18.4}{Every weight in the error sum is positive. Substitute the upper bound $\eps_j\le aq^j$ term by term:
\begin{align*}
 \sum_{j=0}^{n-1}q^{n-1-j}\eps_j
 &\le\sum_{j=0}^{n-1}a q^{n-1-j}q^j\\
 &=\sum_{j=0}^{n-1}a q^{n-1}\\
 &=a n q^{n-1}.
\end{align*}
The exponent becomes $n-1-j+j=n-1$ in every term, and there are $n$ terms. The theorem therefore gives $e_n\le q^n e_0+a n q^{n-1}$.

We still must show that the factor $n$ does not prevent convergence to zero. It is not enough to say that $q^{n-1}$ tends to zero, because it is being multiplied by a growing quantity. Choose $r=(q+1)/2$, so $q<r<1$, and set $t=q/r$, so $0<t<1$. For each positive integer $n$, every one of $1,t,\ldots,t^{n-1}$ is at least $t^{n-1}$. Hence
\[
 n t^{n-1}\le1+t+\cdots+t^{n-1}\le\frac1{1-t}.
\]
Multiply by $r^{n-1}$ and use $q=rt$:
\[
 0\le n q^{n-1}=r^{n-1}n t^{n-1}
 \le\frac{r^{n-1}}{1-t}.
\]
The denominator is a fixed positive number and $r^{n-1}\to0$, so the upper bound tends to zero by the geometric-limit argument. The initial term $q^n e_0$ also tends to zero. For any requested tolerance, make each of these two terms smaller than half the tolerance; one threshold beyond both suffices. Thus the full error bound tends to zero.}
\sol{18.5}{A task of capacity two can serve two students, while a task of capacity zero serves none. The desired assignment is no longer an injective function into task names; several students may use the same task within its capacity. We translate the problem into an ordinary matching by distinguishing its individual slots.

\emph{Necessity.} Suppose a valid assignment exists, and choose any student set $S$. Every student in $S$ uses one unit of a task in $N(S)$. Task $r$ can contribute at most $c_r$ units, so the total available to the group is at most $\sum_{r\in N(S)}c_r$. Serving $|S|$ students therefore requires
\[
 |S|\le\sum_{r\in N(S)}c_r.
\]
This is a statement about every group, not just the entire class.

\emph{Sufficiency.} Replace task $r$ by distinct vertices $(r,1),\ldots,(r,c_r)$. For example, a capacity-two task $p$ and capacity-one task $q$ become three slots $(p,1),(p,2),(q,1)$. A student accepting $p$ is connected to both $p$ slots. No slot is created for a zero-capacity task.

For general capacities, the resulting bipartite graph is finite. A student set $S$ has as neighboring slots precisely all slots belonging to tasks in $N(S)$. Slots of different tasks have different first coordinates, so their sets are disjoint. The number of these neighbors is therefore exactly $\sum_{r\in N(S)}c_r$. The assumed inequality is Hall's condition for this expanded graph.

Hall's theorem gives each student a distinct acceptable slot. Discard the slot's second coordinate and assign the student to its task. At most $c_r$ students can be assigned to $r$, because that task has only $c_r$ distinct slots. Every student is served acceptably, proving sufficiency. The translation preserves both acceptability and capacity, which is why solving the expanded problem solves the original one.}
\sol{18.6}{The exact question is whether bounded additive update errors preserve convergence to the fixed point of $T(x)=(3x+1)/4$, and what bound replaces that false claim. The prerequisites are real arithmetic, metric distance, the contraction theorem, induction, and finite geometric sums. The proof skeleton is: insert the exact update; obtain $e_{n+1}\le qe_n+\eps_n$; unfold by induction; sum the nonnegative geometric weights; then test a constant bias.

With bias $1/100$, the recurrence has limit $26/25$, not the original fixed point one, refuting the initial claim. At $n=1$, the exact iterate is $13/50$, the true error is $37/50$, and the general certificate is $3/4+1/100=19/25$. Thus $37/50\le38/50$, an exact check. The full formula $y_n=(26/25)(1-(3/4)^n)$ is proved in the chapter, not inferred from that one check.

The permitted novelty statement is that the project reconstructs and studies a standard contraction perturbation bound with reproducible examples; it does not establish a new theorem relative to the literature. One further question is whether a specified noncontractive map admits a comparable estimate under a different stability hypothesis. The present proof does not answer it because its tail bounds require a factor strictly below one.}


\chapter{Running the laboratories and checking your work}\label{app:practical}
\section{The smallest amount of Python needed}
The core proofs can be read without programming. For the computational laboratories, the companion uses Python 3.10 or later and only its standard library. The following vocabulary is sufficient to begin; the official tutorial explains lists and sets in more detail~\cite{pystructures}.

A \emph{list} is an ordered collection written in square brackets. In \code{a = [1, 1, 1, -1]}, the first entry is \code{a[0]}, not \code{a[1]}. Python starts list indices at zero. A \emph{set} records membership without repeated elements; \code{\{0, 1\}} is the set containing zero and one. The empty set is written \code{set()}, because empty braces have a different meaning in Python. The expression \code{len(a)} is the number of entries.

A loop repeats a block for successive values. Indentation identifies the repeated block. For example:
\par\noindent\begin{minipage}{\linewidth}
\begin{lstlisting}
a = [1, 1, 1, -1]
value = 0
for j in range(3):
    value = value + a[j] * a[j + 1]
print(value)
\end{lstlisting}
\end{minipage}\par

The values of \code{j} are $0,1,2$; the upper endpoint in \code{range(3)} is excluded. The successive additions are $1,1,-1$, so the printed result is one, the aperiodic correlation $A_1$. The expression \code{value = value + ...} updates a stored variable; it is not a mathematical equality asserted for a fixed value on both sides. Python uses \code{==} to test equality.

A function packages a calculation. The following educational example assumes valid indices:
\par\noindent\begin{minipage}{\linewidth}
\begin{lstlisting}
def correlation(a, shift):
    value = 0
    for j in range(len(a) - shift):
        value = value + a[j] * a[j + shift]
    return value
\end{lstlisting}
\end{minipage}\par

The \code{return} line gives the output. Unlike the teaching fragment, the companion validates its inputs and rejects empty sign sequences or invalid aperiodic shifts. It keeps periodic and aperiodic calculations as different named functions so that their conventions cannot be confused silently.

\section{An exact fraction is not a rounded decimal}
Use \code{Fraction} when the intended inputs are rational. It stores a numerator and denominator and performs rational arithmetic rather than immediately converting to a floating-point approximation~\cite{pyfractions}.
\par\noindent\begin{minipage}{\linewidth}
\begin{lstlisting}
from fractions import Fraction
x = Fraction(1, 3)
y = (x * x + 1) / 3
print(y)
\end{lstlisting}
\end{minipage}\par

This prints \code{10/27}. Use \code{Fraction(1, 10)} for the exact rational one tenth. Passing a floating-point value such as \code{0.1} first represents that floating-point input, not an intention to supply the exact decimal fraction.

The numerical tables in the sign laboratory use integer arithmetic. The iteration checks use fractions. This removes one source of roundoff uncertainty, but does not prove that a loop enumerates the intended objects or that the mathematical assumptions hold. Exact arithmetic and correct modeling are separate requirements.

\subsection*{Running a command is different from entering a calculation}
The archive contains source files. Extract it into one folder before running them. A \emph{working directory} is the folder in which a terminal command currently looks for files. Open a terminal in the extracted folder, or change to it with a command such as \code{cd path/to/extracted/folder}, replacing that example path by the actual location.

Check the installed Python interpreter with \code{python --version} or \code{python3 --version}. Use whichever command names the compatible interpreter on your system. A command such as \code{python labs.py} runs the entire saved file. By contrast, the examples beginning \code{from labs import ...} are Python statements: put them in a small \code{.py} file in the same folder and run that file, or enter them at Python's interactive prompt after starting the interpreter.

The interactive prompt is usually displayed as \code{>>>}. It indicates where Python expects a statement; do not copy the prompt characters into a source file. A terminal's shell prompt expects commands for programs, whereas Python's prompt expects Python expressions and statements. An error caused by pasting into the wrong prompt says nothing about the mathematical claim being tested.

The saved \code{labs.py} and \code{verify.py} files contain the full examples and checks. A short teaching fragment illustrates a language feature; it should not be mistaken for a second complete replacement of those programs.

\section{Three reproducible starting experiments}
Extract the source package and open a terminal in its directory. The command \code{python labs.py} runs the demonstrations; on systems where the interpreter is named \code{python3}, use that name instead. The files require no package installation beyond a compatible Python interpreter.

For assignments, execute:
\par\noindent\begin{minipage}{\linewidth}
\begin{lstlisting}
from labs import assign
answer = assign([{0, 1}, {0}, {1, 2}], 3)
print(answer)
\end{lstlisting}
\end{minipage}\par

Student indices $0,1,2$ correspond to $a,b,c$ in Chapter~8. Task indices $0,1,2$ correspond to the book's tasks $1,2,3$. The output pairs use these zero-based names. A true \code{saturates\_left} flag comes with a full assignment. A false flag comes with \code{shortage\_left} and \code{shortage\_neighbors}; check the claimed set of neighbors and its size yourself. The function does not claim to return a maximum-size partial matching when it stops with an impossibility certificate.

For signs, execute:
\par\noindent\begin{minipage}{\linewidth}
\begin{lstlisting}
from labs import aperiodic, periodic, sign_counts
print(aperiodic([1, 1, 1, -1], 1))
print(periodic([1, 1, 1, -1], 1))
print(sign_counts(4))
\end{lstlisting}
\end{minipage}\par

The first two outputs are one and zero. The last output reports counts at each positive length up to four. Counts refer to all labeled sign sequences, with no division by negation, reversal, or rotation symmetries. Increasing the maximum length doubles the number of additional candidates at each next length; choose a finite range deliberately.

For the final investigation, execute:
\par\noindent\begin{minipage}{\linewidth}
\begin{lstlisting}
from labs import biased_iteration
for row in biased_iteration(8):
    print(row)
\end{lstlisting}
\end{minipage}\par

Each row contains the index, exact iterate, exact error from one, and the proven upper bound. The strings such as \code{13/50} preserve the exact fractions. Convert to decimals only when preparing a display and retain the exact values in the evidence record.

\section{What the verification program checks}
The command \code{python verify.py} writes \code{verification.json}. It compares the matching algorithm with both Hall's subset condition and an independent exhaustive assignment search for every bipartite graph whose left and right sizes each range from zero through four. It verifies the explicit certificates, not just whether two Boolean answers agree.

It also checks exact second-neighborhood membership in small oriented graphs; the deterministic cut guarantee on small undirected graphs; Cauchy--Schwarz on a finite collection of integer vectors; elementary finite-map statements; the sign-correlation identities and count table; and the rational iteration formulas and bounds. The checked ranges are recorded explicitly in the JSON file. These are finite regression tests of the implementations and examples, not replacements for the general proofs in the text. The additional command \code{python verify\_revision.py} checks worked examples expanded in this reader edition, including the walk deletion, projection coefficient, and signed-error calculation. The combined build runs both test files.

Run the file normally, not with Python's optimization option that disables assertions. If a test fails, do not edit the expected answer merely to make the output pass. Identify whether the error lies in the definition, algorithm, test oracle, or textbook statement, and reconstruct the relevant mathematical obligation.

\section{The Lean companion and its limits}
The file \code{proofs.lean} contains the complete small examples and exercise solutions from Chapters~15--16, within the namespace \code{MathCourse}. The file \code{lean-toolchain} pins Lean 4.19.0. It does not pin the toolchain of the external research repository; that project must use its own configuration and dependencies.

After installing the official Lean toolchain support, run
\par\noindent\begin{minipage}{\linewidth}
\begin{lstlisting}
lean proofs.lean
python verify.py --lean
\end{lstlisting}
\end{minipage}\par

The second command requires Lean execution in addition to the numerical tests. It fails rather than silently reporting a pass when Lean is unavailable. Inspect the produced \code{lean\_check.log} and the theorem dependency output. The intentionally impossible-context theorem remains a conditional statement, not an unconditional contradiction, even if it uses no extra axiom.

\textbf{Execution status of this edition.} The Lean examples have been reviewed against their mathematical arguments and the documented language constructs, but no Lean executable was available in the authoring environment, and a toolchain could not be obtained there. They were therefore \emph{not compiler-checked in this edition}. The companion is supplied for execution with the pinned toolchain; its acceptance is not represented as an already completed test. The full OpenAI research repository was also not built or independently verified for this book.

For additional practice after these lessons, \emph{Mathematics in Lean} provides an interactive textbook and exercises~\cite{mil}. Use it as a continuation once the statement and proof-state ideas in this course are familiar, not as an unexplained prerequisite imposed before the first lesson.

\section{Rebuilding the book}
The consolidated source is \code{math\_ai.tex}; it contains the chapters, exercises, solutions, bibliography, and TikZ drawings without external image files. A TeX distribution must provide the packages declared in its preamble, including AMS mathematics, New TX fonts, TikZ, \code{tcolorbox}, and \code{listings}. Font binaries are not part of the source package; use those provided by the TeX distribution.

Run \code{python build.py} to compile the document and run the Python verification suite. This script requires \code{pdflatex}. It compiles multiple times to resolve the table of contents and cross-references, rejects missing references and overfull layout boxes, and stops on a failed test. Use \code{python build.py --lean} to require the additional Lean check. The PDF can also be built manually with repeated \code{pdflatex -halt-on-error math\_ai.tex} commands until references stabilize.

The source has no network dependency during compilation. Public research references are hyperlinks for reading, not files fetched during the build. A source archive is useful only if its extracted copy rebuilds; that clean-extraction check is part of the release procedure for this edition.


\chapter{A compact guide to notation and mathematical judgment}\label{app:guide}
\section{A notation map with places to return}
Use this map to return to an explanation, not as a vocabulary test. Similar-looking symbols can have different roles; in particular, function preimages need no inverse function, and a set difference is not numerical subtraction.
\begin{longtable}{@{}p{1.25in}p{3.95in}@{}}
\toprule Symbol & Meaning and first detailed use\\\midrule\endhead
$\NN,\ZZ,\QQ,\RR$ & Natural numbers including zero, integers, rational numbers, real numbers (Chapter~1).\\
$\sum_{j=a}^{b}x_j$ & Add one term for each integer index from $a$ through $b$ (Chapters~1, 5). An empty sum is zero.\\
$\forall,\exists$ & For every; there exists. Order controls permitted dependence (Chapter~2).\\
$\Rightarrow,\Leftrightarrow$ & Implication; both implications. Neither means that the statement has already been proved (Chapter~2).\\
$\neg,\land,\lor$ & Not; and; inclusive or (Chapter~2).\\
$\in,\subseteq$ & Membership; inclusion of sets (Chapter~3).\\
$A\cup B,A\cap B$ & Union and intersection (Chapter~3).\\
$A\setminus B$ & Members of $A$ that are not in $B$ (Chapter~3).\\
$f:A\to B$ & A function with specified domain and codomain (Chapter~3).\\
$f(S), f^{-1}(T)$ & Image of an input set; preimage of an output set (Chapter~3).\\
$g\circ f$ & Apply $f$ first, then $g$ (Chapter~3).\\
$\id_A$ & Identity function on $A$ (Chapter~3).\\
$a\equiv b\pmod m$ & $m$ divides $a-b$ (Chapter~6).\\
$[a]_m$ & The whole equivalence class of $a$, not merely one representative (Chapter~6).\\
$|S|$ & Number of members of a finite set (Chapter~3). Context distinguishes this from absolute value, introduced in Chapter~2.\\
$N(S)$ & Union of acceptable-task neighbors of a student set (Chapter~8).\\
$N_1^+(v),N_2^+(v)$ & First out-neighbors and exact second out-neighbors (Chapter~7).\\
$\ker T$ & Inputs sent to the zero vector by a linear map (Chapter~10).\\
$u\cdot v$ & Dot product: multiply corresponding entries and add (Chapter~10).\\
$A^{\mathsf T}$ & Transpose, exchanging rows and columns (Chapter~10).\\
$\norm x_2$ & Euclidean length (Chapter~10).\\
$d(x,y)$ & Distance in the specified metric; the choice of metric is part of the problem (Chapter~11).\\
$x_n\to x_*$ & Eventual closeness at every positive tolerance (Chapter~11).\\
$\eps,\delta$ & Names frequently used for positive tolerances, not fixed universal constants (Chapters~2, 11).\\
$\E X,\Prob(E)$ & Expectation and probability in the finite model (Chapter~13).\\
$A_s,C_s$ & Aperiodic and periodic correlations; check the index ranges (Chapter~14).\\
$B(c,r)$ & Points whose distance from $c$ is strictly less than $r$ (Chapter~17).\\
\bottomrule
\end{longtable}

\section{The underlying structures are a connected map}
A set tells us which objects are allowed. A relation specifies how pairs are connected or compared. An operation combines objects. A map translates one description into another. Additional structure determines which questions make sense: addition permits linear combinations, order permits extremal choices, distance permits approximation, and probability permits weighted averaging.

Ask what information a proof actually uses. Hall's theorem uses a finite adjacency relation and cardinalities, not coordinates of the vertices. The contraction theorem uses distances and completeness, not vector addition. Cauchy--Schwarz uses the algebra and positivity of a dot product. The sign-sequence argument translates pair products into row orthogonality and then into a counting restriction. Seeing what is \emph{not} used helps identify possible generalizations without guessing them indiscriminately.

A representation change is useful only when its correspondence is controlled. An inverse converts equal outputs to equal inputs. A quotient forgets distinctions shown to be irrelevant. A matrix represents a linear action in chosen coordinates. An indicator turns a counting question into a sum of zero-one values. These are different constructions, but each creates a language in which a desired step becomes easier to express and justify.

\section{Choosing a next move instead of listing tricks}
When a universal claim is doubtful, look for a small or boundary counterexample before trying a long proof. The false function-composition converse failed on one unused intermediate input. The false convergence claim failed with a constant bias. Small examples were selected to expose a missing mechanism, not to accumulate reassuring cases.

When existence is difficult, ask what a failure would force. A failed assignment search exposes a shortage. A minimum counterexample can impose extra structure. An average can certify that at least one acceptable object exists. These approaches prove existence in different ways; choose according to whether the objects admit a repair process, an extremal reduction, or a useful average.

When a formula is complicated, ask what quantity the argument actually needs to control. An explicit fixed point may be unnecessary if a residual bound gives the required accuracy. A full classification may be unnecessary if a necessary condition already rules out the application. Conversely, do not let an easier surrogate replace the original target without recording that change.

When a bound appears, examine equality and stability. The Cauchy--Schwarz equality case identifies aligned vectors. The inexact-iteration example shows that $1/(1-q)$ is not merely an algebraic inconvenience. A mathematically informed question asks both ``Can this constant be improved?'' and ``Which example would prevent an improvement?''

\section{A reading difficulty should become a precise question}
Suppose you understand the theorem statement but not a line of its proof. Do not begin by asking the assistant to rewrite the whole chapter. Copy the exact line, the assumptions currently available, and the last inference you understand. Then ask which definition or earlier fact connects them. An effective answer may be a single expanded calculation or a small example, rather than a new proof technique.

For example: ``I know the two neighborhoods can overlap. Why can their union be counted as $|N(S)|+|N(T)\setminus N(S)|$?'' The needed explanation is a disjoint decomposition, not a general course on graph theory. Or: ``I know each update error gets multiplied by $q$ later. Why is its weight $q^{n-1-j}$?'' The needed explanation is to count the subsequent updates, as in Chapter~18.

After the answer, return to the original proof and restate the repaired inference yourself. Next alter a small feature: allow the sets to be equal, change the first update index, or remove the positivity of a divisor. Explain what changes. This is a test of whether the explanation supports reasoning beyond the exact wording supplied by the assistant.

The same method applies to a definition. Ask first what kinds of objects it concerns, then what requirements an example must meet, then for an example and a near miss. Finally identify whether the proof is unpacking an assumed instance of the definition or trying to establish that a new object satisfies it. Those two directions were already present in the first parity proof.

\section{A reading record that remains useful after the conversation}
A finished reading record contains the exact claim and source version; a plain-language interpretation; the objects and hypotheses; a diagram or paragraph describing the proof's dependencies; and one expanded delicate step. Beside each borrowed result, state its hypotheses and why they hold here. Beside each computation, state the finite range and arithmetic model. Beside each formal result, state which proposition was checked and which assumptions it uses.

The final page should contain your own explanation and one justified variation of the question. A saved chat log is not an adequate substitute: it mixes guesses, corrections, and results without necessarily identifying which statements survived scrutiny. The record should let another reader reconstruct the final argument without reproducing the entire conversation.

\section{A suggested capstone rubric}
\begin{center}
\begin{tabularx}{\textwidth}{@{}p{1.15in}p{0.45in}X@{}}\toprule
Dimension & Weight & Evidence\\\midrule
Claim and scope & 20\% & Exact domains, hypotheses, conclusion, and distinction between original and revised questions.\\
Argument & 35\% & Complete proof or counterexample, with delicate inferences and dependencies explained.\\
Understanding & 20\% & Unaided explanation and a correct response when an assumption or example is changed.\\
Reproducibility & 15\% & Executable checks where appropriate, recorded versions and ranges, or a justified noncomputational verification route.\\
Context & 10\% & Accurate attribution, restrained novelty statement, and transparent account of AI assistance and remaining limits.\\\bottomrule
\end{tabularx}
\end{center}
A changed theorem, a missing decisive inference, or an unsupported claim of novelty must be repaired; presentation cannot compensate for those defects. This rubric measures a research apprenticeship, not whether a famous problem was solved.

\section{What completion should feel like}
You do not need to remember every displayed formula. You should be able to recognize what a statement requires, locate the obstruction to your understanding, and ask for a concept or inference precise enough to remove it. You should be able to check the resulting explanation on an example, reconstruct its role in the whole proof, and state what has not yet been established.

The next unfamiliar theorem will still contain unfamiliar mathematics. The difference is that you now have a route into it: identify the structures, expose the dependency chain, learn the missing pieces, and return to the argument. Collaboration becomes useful when the assistant's output creates understanding and checkable evidence rather than merely replacing an empty page with mathematical-looking text.


\backmatter
\begin{thebibliography}{99}
\raggedright
\addcontentsline{toc}{chapter}{References and further reading}
\bibitem{mcs} E. Lehman, F. T. Leighton, and A. R. Meyer, \emph{Mathematics for Computer Science}. MIT OpenCourseWare, course reading materials. A further source for logic, induction, graphs, and probability. \url{https://ocw.mit.edu/courses/6-1200j-mathematics-for-computer-science-spring-2024/pages/readings/}.
\bibitem{oairepo} OpenAI, \emph{Math repository}, catalogue and verification descriptions. \url{https://github.com/openai/math}.
\bibitem{oai173} OpenAI, ``Seymour second-neighborhood'' formalization scope, Family 173. \url{https://raw.githubusercontent.com/openai/math/main/lean/docs/173.md}.
\bibitem{oai179} OpenAI, ``Circulant Hadamard'' formalization scope, Family 179. \url{https://raw.githubusercontent.com/openai/math/main/lean/docs/179.md}.
\bibitem{oaimain} OpenAI, second-neighborhood theorem entry point. \url{https://raw.githubusercontent.com/openai/math/main/lean/OAI/Combinatorics/SecondNeighborhood/Main.lean}.
\bibitem{oaireduction} OpenAI, second-neighborhood counterexample reduction. \url{https://raw.githubusercontent.com/openai/math/main/lean/OAI/Combinatorics/SecondNeighborhood/Reduction.lean}.
\bibitem{oaiconfig} OpenAI, comparison configuration for the second-neighborhood theorem. \url{https://raw.githubusercontent.com/openai/math/main/lean/ComparatorChallenges/SeymourSecondNeighborhood.json}.
\bibitem{leaninstall} Lean, official installation instructions. \url{https://lean-lang.org/install/}.
\bibitem{leanprop} Lean, \emph{Theorem Proving in Lean 4}, ``Propositions and Proofs.'' \url{https://lean-lang.org/theorem_proving_in_lean4/Propositions-and-Proofs/}.
\bibitem{leanquant} Lean, \emph{Theorem Proving in Lean 4}, ``Quantifiers and Equality.'' \url{https://lean-lang.org/theorem_proving_in_lean4/Quantifiers-and-Equality/}.
\bibitem{leantactics} Lean, \emph{Theorem Proving in Lean 4}, ``Tactics.'' \url{https://lean-lang.org/theorem_proving_in_lean4/Tactics/}.
\bibitem{leanaxioms} Lean, \emph{The Lean Language Reference}, ``Axioms.'' \url{https://lean-lang.org/doc/reference/latest/Axioms/}.
\bibitem{leanvalidate} Lean, \emph{The Lean Language Reference}, ``Validating a Lean Proof.'' \url{https://lean-lang.org/doc/reference/latest/ValidatingProofs/}.
\bibitem{mil} J. Avigad and P. Massot, \emph{Mathematics in Lean}, with contributors. Interactive textbook and exercises. \url{https://leanprover-community.github.io/mathematics_in_lean/}.
\bibitem{pystructures} Python Software Foundation, \emph{The Python Tutorial}, ``Data Structures.'' \url{https://docs.python.org/3/tutorial/datastructures.html}.
\bibitem{pyfractions} Python Software Foundation, \emph{The Python Standard Library}, ``fractions---Rational numbers.'' \url{https://docs.python.org/3/library/fractions.html}.
\end{thebibliography}
\noindent\small Public documentation and repository pages were consulted on 7 October 2026. These links are reading references, not a frozen copy of the external repository or evidence of a completed repository build. The executable classroom companion is version-pinned separately. The mathematical proofs written in the course are pedagogical derivations of standard results unless explicitly described otherwise.
\end{document}
